% Leçon 24 — Grammaire : 越来越； ; 一方面……,另一方面…… — Chinois 1ère A — Cours signé OnBuch+
% Programme officiel MINESEC. Compiler avec : tectonic ZHA24-grammaire.tex
\newcommand{\DOCMATIERE}{Chinois}
\newcommand{\DOCNIVEAU}{1ère A}
\newcommand{\DOCMODULE}{Module 3 : Citoyenneté et ouverture au monde}
\newcommand{\DOCLECON}{ZHA24}
\newcommand{\DOCTITRE}{Leçon 24 — Grammaire : 越来越； ; 一方面……,另一方面……}
% OnBuch+ — préambule « cours signés » ENRICHI (compilé avec tectonic / XeLaTeX)
% Système visuel « Pop » d'OnBuch+ : fond crème (jamais blanc), bordures encre
% épaisses, ombres « dures » décalées, titres Archivo Black en capitales.
% Version MATHÉMATIQUES : graphes (pgfplots/tikz), boîtes pédagogiques
% (définition, propriété, méthode, attention, le savais-tu, exercice résolu,
% exercice, TP, bilan), page de garde et signature OnBuch+ ancrée en bas.
%
% Macros attendues dans main.tex AVANT \input{preamble} :
%   \DOCMATIERE  \DOCNIVEAU  \DOCMODULE  \DOCLECON (numéro)  \DOCTITRE
\documentclass[11pt]{article}

\usepackage{fontspec}
\usepackage[french]{babel} % français = langue principale ; le chinois via \zh{..} / textebox (xeCJK)
\usepackage{xeCJK}
\usepackage{pifont} % \ding{51} \ding{55} utilisés par les modèles
\usepackage[a4paper,margin=2cm,top=3.4cm,bottom=2.8cm,headheight=1.4cm,footskip=1.3cm]{geometry}
\usepackage{amsmath,amssymb,amsfonts}
\usepackage{mathtools} % \xmapsto, \coloneqq... souvent utilisés par les modèles
\usepackage{cancel} % simplifications barrées (bilans de forces)
% \abs{z} (module, valeur absolue) et \norm{u} (norme), utilisés par les modèles
\providecommand{\abs}{}\let\abs\relax\DeclarePairedDelimiter{\abs}{\lvert}{\rvert}
\providecommand{\norm}{}\let\norm\relax\DeclarePairedDelimiter{\norm}{\lVert}{\rVert}
\usepackage[table]{xcolor} % [table] : fournit aussi \rowcolors (alternance de lignes)
\usepackage{pagecolor}
\usepackage{tikz}
\usetikzlibrary{arrows.meta,positioning,calc,shapes.geometric,shapes.misc,shapes.arrows,decorations.pathmorphing,decorations.pathreplacing,decorations.markings,patterns,shadows,mindmap,trees,fit,backgrounds,matrix,intersections,angles,quotes}
\usepackage{pgfplots}
\usepackage[version=4]{mhchem} % équations nucléaires \ce{^{235}_{92}U}
\usepackage[european,siunitx]{circuitikz} % schémas électriques
\pgfplotsset{compat=1.18}
\usepgfplotslibrary{fillbetween,groupplots} % aires entre courbes (calcul intégral)
\usepackage{enumitem}
\usepackage{fancyhdr}
% [most] charge toutes les bibliothèques tcolorbox (listings, minted, raster,
% documentation...) et gonfle inutilement la mémoire TeX sur un document long
% (~300 boîtes) : on ne charge que ce qui est réellement utilisé.
\usepackage[skins,breakable]{tcolorbox}
\usepackage{graphicx}
\usepackage{colortbl}
\usepackage{booktabs}
\usepackage{tabularx}
\usepackage{diagbox} % case d'en-tête diagonale des tableaux à double entrée
% Colonnes extensibles centrée / gauche / droite (C, L, R), souvent utilisées par les modèles
\newcolumntype{C}{>{\centering\arraybackslash}X}
\newcolumntype{L}{>{\raggedright\arraybackslash}X}
\newcolumntype{R}{>{\raggedleft\arraybackslash}X}
\usepackage{array}
\usepackage{multicol}
\usepackage{multirow}
\usepackage{siunitx}
% parse-numbers=false : accepte aussi \SI{2,5 \times 10^{-3}}{...}
\sisetup{locale=FR,output-decimal-marker={,},group-minimum-digits=4,parse-numbers=false}
\DeclareSIUnit\ohm{\ensuremath{\symup{\Omega}}} % U+2126 absent de la police
\DeclareSIUnit\division{div} % réglages d'oscilloscope (ms/div, V/div)
\DeclareSIUnit\tr{tr} % tours (vitesse de rotation, ex. \SI{1500}{\tr\per\minute})
\DeclareSIUnit\molar{M} % concentration molaire (ex. \SI{3}{\molar}, \SI{1}{\micro\molar})
\DeclareSIUnit\an{an} % année (le modèle confond parfois avec \year, un compteur TeX) — ex. \SI{5}{\kilo\meter\per\an}
\DeclareSIUnit\FCFA{FCFA} % franc CFA (ex. \SI{500}{\FCFA\per\kilo\gram})
\DeclareSIUnit\degreeBrix{\textdegree Brix} % teneur en sucre d'un sirop/jus (réfractométrie)
\DeclareSIUnit\month{mois} % le modèle utilise parfois \month, un compteur TeX, comme unité de durée
\usepackage{qrcode}
% \degree : commande fréquemment utilisée par les modèles pour un angle ou une
% température en dehors de \SI{}{} (non fournie par les paquets chargés).
\providecommand{\degree}{^{\circ}}
% \qed (fin de démonstration), utilisé par les modèles sans amsthm
\providecommand{\qed}{\hfill$\square$}
% \barre : double barre des valeurs interdites dans un tableau de variations (tkz-tab)
\providecommand{\barre}{\ensuremath{\|}}
% environnement proof (amsthm non chargé) utilisé par les modèles
\ifdefined\proof\else\newenvironment{proof}[1][Démonstration]{\par\noindent\textit{#1.}\ }{\qed\par}\fi
\usepackage{lastpage}
\usepackage{hyperref}
\hypersetup{hidelinks}

% ============================================================
% Palette « Pop » OnBuch+ — mêmes hex que la classe P (plus_ui.dart)
% ============================================================
\definecolor{poporange}{HTML}{F59321}
\definecolor{popdark}{HTML}{E07A0C}
\definecolor{popcream}{HTML}{FFF6E8}
\definecolor{popink}{HTML}{1C1714}
\definecolor{poppurple}{HTML}{7A5AE0}
\definecolor{popgreen}{HTML}{2E9E6A}
\definecolor{poppink}{HTML}{E0457B}
\definecolor{popblue}{HTML}{2F7DE1}
\definecolor{popgold}{HTML}{C98A12}
\definecolor{popmuted}{HTML}{8A7F76}
\definecolor{popline}{HTML}{EADFD2}
\definecolor{popgray}{HTML}{8A7F76}
\colorlet{popgrey}{popgray}
% Alias tolérants (le modèle nomme parfois les couleurs différemment)
\colorlet{popwhite}{white}
\colorlet{popblack}{popink}
\colorlet{popred}{poppink}
\colorlet{popyellow}{popgold}
% Teintes claires (fonds de tableaux/encadrés) — le modèle nomme parfois
% les couleurs avec un suffixe L (ex. popblueL) sans les avoir définies.
\colorlet{poporangeL}{poporange!20}
\colorlet{popdarkL}{popdark!20}
\colorlet{poppurpleL}{poppurple!20}
\colorlet{popgreenL}{popgreen!20}
\colorlet{poppinkL}{poppink!20}
\colorlet{popblueL}{popblue!20}
\colorlet{popgoldL}{popgold!20}
\colorlet{popredL}{popred!20}
\colorlet{popgrayL}{popgray!20}
% Teintes claires pour fonds de tableaux / zones
\colorlet{popblueL}{popblue!10!white}
\colorlet{poporangeL}{poporange!14!white}
\colorlet{popgreenL}{popgreen!12!white}
\colorlet{poppurpleL}{poppurple!12!white}
\colorlet{poppinkL}{poppink!12!white}
\colorlet{popgoldL}{popgold!14!white}

\pagecolor{popcream}

% ============================================================
% Typographie
% ============================================================
\defaultfontfeatures{Ligatures=TeX}
\setmainfont{PlusJakartaSans-Medium.ttf}[Path=../../../fonts/,BoldFont=PlusJakartaSans-Bold.ttf]
% Symboles, API et emojis absents de la police principale : repli sur DejaVu Sans (caractères actifs)
\newfontface\symfont{DejaVuSans.ttf}[Path=../../../fonts/]
\makeatletter
\newcommand\symactive[1]{\catcode#1=13 \begingroup\lccode`\~=#1 \lowercase{\endgroup\def~}{{\symfont\char#1}}}
\newcommand\symblank[1]{\catcode#1=13 \begingroup\lccode`\~=#1 \lowercase{\endgroup\def~}{}}
\newcommand\symhyph[1]{\catcode#1=13 \begingroup\lccode`\~=#1 \lowercase{\endgroup\def~}{\mbox{-}}}
\newcount\symc
\newcommand\symrange[2]{\symc=#1 \@whilenum\symc<#2 \do{\expandafter\symactive\expandafter{\the\symc}\advance\symc by 1 }}
\newcommand\symblankrange[2]{\symc=#1 \@whilenum\symc<#2 \do{\expandafter\symblank\expandafter{\the\symc}\advance\symc by 1 }}
\makeatother
\symrange{"0250}{"02FF}\symactive{"2713}\symactive{"2714}\symactive{"2717}\symactive{"2718}\symactive{"2610}\symactive{"2611}\symactive{"2612}
\symrange{"2600}{"27BF}
\catcode"2705=13 \begingroup\lccode`\~="2705 \lowercase{\endgroup\def~}{{\symfont\char"2713}}\symblank{"26BD}\symblank{"26A1}
\symrange{"2460}{"257F}\symactive{"223C}\symactive{"2248}\symactive{"2260}
\symactive{"01F8}\symactive{"01F9}\symactive{"1E3E}\symactive{"1E3F}
\symhyph{"2011}\symhyph{"2E1A}\symblankrange{"1F300}{"1FAFF}\symblank{"FE0F}
% Caractères chinois : WenQuanYi Zen Hei (sous-ensemble GB2312 + pinyin), gras/italique simulés
\setCJKmainfont{WenQuanYiZenHei-subset.ttf}[Path=../../../fonts/,AutoFakeBold=3,AutoFakeSlant=0.2]
\xeCJKsetup{CJKspace=false}
% Pinyin : voyelles à caron (ǎ ǐ ǒ ǔ ǖ ǘ ǚ ǜ) absentes de la police principale -> caractères actifs
\newfontface\pyfont{WenQuanYiZenHei-subset.ttf}[Path=../../../fonts/]
\newcommand\pyactive[1]{\catcode#1=13 \begingroup\lccode`\~=#1 \lowercase{\endgroup\def~}{{\pyfont\char#1}}}
\pyactive{"01CD}\pyactive{"01CE}\pyactive{"01CF}\pyactive{"01D0}\pyactive{"01D1}\pyactive{"01D2}\pyactive{"01D3}\pyactive{"01D4}\pyactive{"01D5}\pyactive{"01D6}\pyactive{"01D7}\pyactive{"01D8}\pyactive{"01D9}\pyactive{"01DA}\pyactive{"01DB}\pyactive{"01DC}
% Maths en sans-serif (Fira Math) avec chiffres et lettres droites de Plus
% Jakarta, pour que les formules se fondent dans le texte.
\usepackage[math-style=ISO,bold-style=ISO]{unicode-math}
\setmathfont{FiraMath-Regular.otf}[Path=../../../fonts/]
\setmathfont{PlusJakartaSans-Medium.ttf}[Path=../../../fonts/,range={up/num,up/{Latin,latin},\mathup}]
\usepackage{icomma} % virgule décimale sans espace en mode math
% Caractères mathématiques Unicode absents des polices de texte (Plus Jakarta,
% Archivo Black) : rendus par leur équivalent mathématique, en texte comme en
% formule — sinon ils s'affichent en carré vide (ex. « Arithmétique dans ℤ »).
\usepackage{newunicodechar}
\newunicodechar{ℝ}{\ensuremath{\mathbb{R}}}
\newunicodechar{ℤ}{\ensuremath{\mathbb{Z}}}
\newunicodechar{ℂ}{\ensuremath{\mathbb{C}}}
\newunicodechar{ℕ}{\ensuremath{\mathbb{N}}}
\newunicodechar{ℚ}{\ensuremath{\mathbb{Q}}}
\newunicodechar{𝔻}{\ensuremath{\mathbb{D}}}
\newunicodechar{α}{\ensuremath{\alpha}}
\newunicodechar{β}{\ensuremath{\beta}}
\newunicodechar{γ}{\ensuremath{\gamma}}
\newunicodechar{δ}{\ensuremath{\delta}}
\newunicodechar{ε}{\ensuremath{\varepsilon}}
\newunicodechar{θ}{\ensuremath{\theta}}
\newunicodechar{λ}{\ensuremath{\lambda}}
\newunicodechar{μ}{\ensuremath{\mu}}
\newunicodechar{σ}{\ensuremath{\sigma}}
\newunicodechar{τ}{\ensuremath{\tau}}
\newunicodechar{φ}{\ensuremath{\varphi}}
\newunicodechar{ψ}{\ensuremath{\psi}}
\newunicodechar{ω}{\ensuremath{\omega}}
\newunicodechar{Σ}{\ensuremath{\Sigma}}
% majuscules produites par \MakeUppercase dans les titres (α-aminés -> Α)
\newunicodechar{Α}{\ensuremath{\alpha}}
\newunicodechar{Β}{\ensuremath{\beta}}
\newunicodechar{Γ}{\ensuremath{\Gamma}}
\newunicodechar{Δ}{\ensuremath{\Delta}}
\newunicodechar{Ε}{\ensuremath{\varepsilon}}
\newunicodechar{Θ}{\ensuremath{\Theta}}
\newunicodechar{Λ}{\ensuremath{\Lambda}}
\newunicodechar{Μ}{\ensuremath{\mu}}
\newunicodechar{Τ}{\ensuremath{\tau}}
\newunicodechar{Φ}{\ensuremath{\Phi}}
\newunicodechar{Ψ}{\ensuremath{\Psi}}
\newunicodechar{Ω}{\ensuremath{\Omega}}
\newunicodechar{↦}{\ensuremath{\mapsto}}
\newunicodechar{∈}{\ensuremath{\in}}
\newunicodechar{∉}{\ensuremath{\notin}}
\newunicodechar{∧}{\ensuremath{\wedge}}
\newunicodechar{⇔}{\ensuremath{\Leftrightarrow}}
\newunicodechar{⟺}{\ensuremath{\Longleftrightarrow}}
\newunicodechar{⇒}{\ensuremath{\Rightarrow}}
\newunicodechar{⟹}{\ensuremath{\Longrightarrow}}
\newunicodechar{∘}{\ensuremath{\circ}}
\newunicodechar{∪}{\ensuremath{\cup}}
\newunicodechar{∩}{\ensuremath{\cap}}
\newunicodechar{≡}{\ensuremath{\equiv}}
\newunicodechar{⊂}{\ensuremath{\subset}}
\newunicodechar{⊥}{\ensuremath{\perp}}
\newunicodechar{∃}{\ensuremath{\exists}}
\newunicodechar{∀}{\ensuremath{\forall}}
\newunicodechar{∛}{\ensuremath{\sqrt[3]{\,}}}
\newunicodechar{∜}{\ensuremath{\sqrt[4]{\,}}}
\newunicodechar{⃗}{\ensuremath{{}^{\to}}}
\newunicodechar{ⁿ}{\ifmmode^{n}\else\textsuperscript{\ensuremath{n}}\fi}
\newunicodechar{ˣ}{\ifmmode^{x}\else\textsuperscript{\ensuremath{x}}\fi}
\newunicodechar{ᵇ}{\ifmmode^{b}\else\textsuperscript{\ensuremath{b}}\fi}
\newunicodechar{ʳ}{\ifmmode^{r}\else\textsuperscript{\ensuremath{r}}\fi}
\newunicodechar{ᵘ}{\ifmmode^{u}\else\textsuperscript{\ensuremath{u}}\fi}
\newunicodechar{ʸ}{\ifmmode^{y}\else\textsuperscript{\ensuremath{y}}\fi}
\newunicodechar{ᵃ}{\ifmmode^{a}\else\textsuperscript{\ensuremath{a}}\fi}
\newunicodechar{ᵉ}{\ifmmode^{e}\else\textsuperscript{\ensuremath{e}}\fi}
\newunicodechar{ᵏ}{\ifmmode^{k}\else\textsuperscript{\ensuremath{k}}\fi}
\newunicodechar{ᵖ}{\ifmmode^{p}\else\textsuperscript{\ensuremath{p}}\fi}
\newunicodechar{ᴺ}{\ifmmode^{N}\else\textsuperscript{\ensuremath{N}}\fi}
\newunicodechar{ᶜ}{\ifmmode^{c}\else\textsuperscript{\ensuremath{c}}\fi}
\newunicodechar{ʰ}{\ifmmode^{h}\else\textsuperscript{\ensuremath{h}}\fi}
\newunicodechar{ᵠ}{\ifmmode^{q}\else\textsuperscript{\ensuremath{q}}\fi}
\newunicodechar{ₙ}{\ifmmode_{n}\else\textsubscript{\ensuremath{n}}\fi}
\newunicodechar{ₐ}{\ifmmode_{a}\else\textsubscript{\ensuremath{a}}\fi}
\newunicodechar{ᵢ}{\ifmmode_{i}\else\textsubscript{\ensuremath{i}}\fi}
\newunicodechar{ᵣ}{\ifmmode_{r}\else\textsubscript{\ensuremath{r}}\fi}
\newunicodechar{ₖ}{\ifmmode_{k}\else\textsubscript{\ensuremath{k}}\fi}
\newunicodechar{ᵦ}{\ifmmode_{\beta}\else\textsubscript{\ensuremath{\beta}}\fi}
\newunicodechar{ₓ}{\ifmmode_{x}\else\textsubscript{\ensuremath{x}}\fi}
\newfontfamily\popdisplay{ArchivoBlack-Regular.ttf}[Path=../../../fonts/]
\newfontfamily\poplogo{SpaceGrotesk-Bold.ttf}[Path=../../../fonts/]
\newfontfamily\popsemi{PlusJakartaSans-SemiBold.ttf}[Path=../../../fonts/]
\newfontfamily\popxbold{PlusJakartaSans-ExtraBold.ttf}[Path=../../../fonts/]
\newfontfamily\popxboldls{PlusJakartaSans-ExtraBold.ttf}[Path=../../../fonts/,LetterSpace=3.0]

\setlist[enumerate]{leftmargin=1.7em,itemsep=5pt,topsep=4pt}
\setlist[itemize]{leftmargin=1.7em,itemsep=4pt,topsep=4pt,label={\color{poporange}\textbullet}}
\setlength{\parindent}{0pt}
\setlength{\parskip}{5pt}
\renewcommand{\arraystretch}{1.25}

% pgfplots : style Pop par défaut
\pgfplotsset{
  every axis/.append style={axis line style={popink,line width=1pt},tick style={popink},
    grid style={popline,line width=0.5pt},label style={font=\small},tick label style={font=\footnotesize}},
  popcurve/.style={poporange,line width=1.8pt,no markers,smooth},
}
% Même style disponible dans un \draw TikZ simple (hors pgfplots)
\tikzset{popcurve/.style={poporange,line width=1.8pt}}
\tikzset{popgrid/.style={popline,very thin}}
\colorlet{popcurve}{poporange} % le nom du style est parfois utilisé comme couleur

% ============================================================
% Pill / squiggle
% ============================================================
\newcommand{\poppill}[3]{%
  \tikz[baseline=-0.55ex]\node[draw=popink,line width=1.2pt,rounded corners=2.6pt,%
    fill=#2,inner xsep=6.5pt,inner ysep=3pt]{%
    \popxboldls\fontsize{7}{7}\selectfont\color{#3}\MakeUppercase{#1}};%
}
\newcommand{\popsquiggle}[1]{%
  \tikz[baseline=-0.4ex]{
    \draw[#1,line width=2.6pt,line cap=round]
      plot [smooth] coordinates {(0,0.09) (0.34,0.01) (0.68,0.21) (1.02,0.01) (1.36,0.10)};
  }%
}
% Mot-clé surligné (marqueur orange sous le texte)
\newcommand{\cle}[1]{{\popxbold\color{popdark}#1}}

% ============================================================
% En-tête / pied
% ============================================================
\pagestyle{fancy}
\fancyhf{}
\renewcommand{\headrulewidth}{0pt}
\renewcommand{\footrulewidth}{0pt}
\lhead{\raisebox{-0.15ex}{\poplogo\fontsize{13}{13}\selectfont\color{popink}OnBuch\color{poporange}+}}
\rhead{\poppill{\DOCMATIERE\ \textbullet\ \DOCNIVEAU\ \textbullet\ Leçon \DOCLECON}{white}{popink}}
\lfoot{\popxbold\fontsize{7.6}{7.6}\selectfont\color{popmuted}\MakeUppercase{Cours signé OnBuch+}\ \textbullet\ {\popsemi\DOCTITRE}}
\rfoot{\tikz[baseline=-0.6ex]\node[rounded corners=8.5pt,fill=popink,minimum height=17pt,minimum width=17pt,inner xsep=5pt,inner sep=0]{\popxbold\fontsize{8}{8}\selectfont\color{popcream}\thepage\,/\,\pageref*{LastPage}};}

% ============================================================
% Titres
% ============================================================
\newcommand{\chaptitre}[2]{%
  \begin{tcolorbox}[enhanced,colback=poporange,colframe=popink,boxrule=2.6pt,arc=11pt,
      drop shadow=popink,left=15pt,right=15pt,top=12pt,bottom=11pt,before skip=3pt,after skip=18pt]
    \poppill{#2}{popink}{popcream}\par\vspace{9pt}
    {\raggedright\hyphenpenalty=10000\exhyphenpenalty=10000\popdisplay\fontsize{22}{24}\selectfont\color{popcream}\MakeUppercase{#1}\par}
  \end{tcolorbox}
}
% Section numérotée : pastille orange avec le numéro + titre Archivo
\newcounter{popsec}
\newcommand{\coursec}[1]{%
  \stepcounter{popsec}\setcounter{popsub}{0}%
  \par\vspace{16pt}\noindent
  \begin{minipage}{\linewidth}
  \tikz[baseline=-0.7ex]\node[draw=popink,line width=1.6pt,fill=poporange,rounded corners=4pt,minimum size=22pt,inner sep=0]{\popdisplay\fontsize{12}{12}\selectfont\color{popcream}\thepopsec};%
  \hspace{8pt}{\popdisplay\fontsize{15}{17}\selectfont\color{popink}\MakeUppercase{#1}}\par
  \vspace{4pt}\noindent{\color{poporange}\rule{36pt}{3.6pt}}
  \end{minipage}\par\nopagebreak\vspace{8pt}
}
\newcounter{popsub}
\newcommand{\courssub}[1]{%
  \stepcounter{popsub}%
  \par\vspace{9pt}\noindent{\popxbold\fontsize{11.5}{13}\selectfont\color{popink}{\color{poporange}\thepopsec.\thepopsub}\hspace{6pt}#1}\par\nopagebreak\vspace{4pt}
}

% ============================================================
% Boîtes pédagogiques — même recette : fond blanc, bordure épaisse colorée,
% ombre dure encre, badge plein. Titre optionnel [#1].
% ============================================================
\tcbset{popbase/.style={breakable,enhanced,colback=white,boxrule=2.3pt,arc=9pt,
  shadow={3pt}{-3pt}{0pt}{popink},left=14pt,right=14pt,top=11pt,bottom=11pt,before skip=11pt,after skip=15pt,
  fontupper=\popsemi\fontsize{10.6}{14.5}\selectfont\color{popink},
  fontlower=\popsemi\fontsize{10.6}{14.5}\selectfont\color{popink}}}
% Coche dessinée (le glyphe ✓ n'existe pas dans les polices du document)
\newcommand{\popcheck}[1]{\tikz[baseline=-0.5ex]\draw[#1,line width=1.6pt,line cap=round,line join=round](-0.13,0.0)--(-0.04,-0.09)--(0.13,0.1);}
\providecommand{\coche}{\popcheck{popgreen}}
\providecommand{\resultat}[1]{\par\smallskip\noindent\fcolorbox{popgreen}{popgreenL}{\parbox{\dimexpr\linewidth-2\fboxsep-2\fboxrule\relax}{\textbf{Résultat.} #1}}\par\smallskip}
\newcommand{\popboxhead}[3]{\poppill{#1}{#2}{white}\ifx\relax#3\relax\else\hspace{6pt}{\popxbold\fontsize{10.6}{12}\selectfont\color{popink}#3}\fi\par\vspace{7pt}}

\newtcolorbox{aretenir}[1][]{popbase,colframe=popgreen,before upper={\popboxhead{À retenir}{popgreen}{#1}}}
% Texte en chinois (document de lecture, dialogue, extrait) : cadre
% d'étude. Un titre optionnel (source, auteur) s'affiche dans l'en-tête.
\newtcolorbox{textebox}[1][]{popbase,colframe=popblue,before upper={\popboxhead{文本 · Texte}{popblue}{#1}},after upper={}}
% Marques ✓ / ✗ (forme correcte / fautive) souvent utilisées par les modèles
\providecommand{\cross}{\textcolor{poppink}{\ensuremath{\times}}}
\providecommand{\xmark}{\cross}\providecommand{\crossmark}{\cross}\providecommand{\cmark}{\ensuremath{\checkmark}}
% Ligne de réponse / justification à compléter (inventée par les modèles)
\providecommand{\justification}[1]{\par\noindent\textit{Justification~:} \dotfill\par}
% \textipa (package tipa non chargé) : la police ne couvre pas l'API -> simple passage
\providecommand{\textipa}[1]{#1}
% Soulignement (ulem non chargé) : \uline, \ul
\providecommand{\uline}[1]{\underline{#1}}\providecommand{\ul}[1]{\underline{#1}}
% \glossaire{\item ...} inventée par les modèles : liste de glossaire
\providecommand{\glossaire}[1]{\begin{itemize}#1\end{itemize}}
% \alert (beamer) inventée par les modèles : texte mis en évidence
\providecommand{\alert}[1]{\textbf{\textcolor{poppink}{#1}}}
% variantes ASCII d'ordinaux écrites en macro
\providecommand{\iere}{\textsuperscript{re}}\providecommand{\ieme}{\textsuperscript{e}}\providecommand{\ere}{\textsuperscript{re}}\providecommand{\eme}{\textsuperscript{e}}\providecommand{\er}{\textsuperscript{er}}\providecommand{\ie}{\textsuperscript{e}}
% Ordinaux français écrits en macro par les modèles : 1\ière, 3\ième
\newcommand{\ière}{\textsuperscript{re}}\newcommand{\ième}{\textsuperscript{e}}
% \exemple (inventée par les modèles) : étiquette « Exemple : »
\providecommand{\exemple}{\textbf{Exemple~:}\ }
% \square utilisable aussi hors mode math (cases à cocher)
\AtBeginDocument{\let\squareMath\square\renewcommand{\square}{\ensuremath{\squareMath}}}
% Mot ou phrase en chinois à l'intérieur d'un paragraphe français
\newcommand{\zh}[1]{\ifmmode\text{#1}\else{#1}\fi}
\let\eng\zh \let\esp\zh \let\all\zh % alias tolérés
% flèches d'intonation inventées par les modèles
\providecommand{\textsearrow}{}\renewcommand{\textsearrow}{\ensuremath{\searrow}}\providecommand{\textnearrow}{}\renewcommand{\textnearrow}{\ensuremath{\nearrow}}
% \fr{..} : mot français (inventé par les modèles)
\providecommand{\fr}[1]{\textit{#1}}
% zhbox : encadré de texte chinois inventé par les modèles
\newenvironment{zhbox}{\begin{quote}}{\end{quote}}
% \courssubsub : titre de niveau 3 inventé par les modèles
\providecommand{\courssubsub}[1]{\par\medskip\noindent\textbf{#1}\par\smallskip}
% \vs : « versus » inventé par les modèles
\providecommand{\vs}{\textit{vs}\ }
% \blank : case à compléter inventée par les modèles
\providecommand{\blank}[1][]{\underline{\hspace{1.6em}}}
\newtcolorbox{exemplebox}[1][]{popbase,colframe=poporange,before upper={\popboxhead{Exemple}{poporange}{#1}}}
\newtcolorbox{definition}[1][]{popbase,colframe=popblue,before upper={\popboxhead{Définition}{popblue}{#1}}}
\newtcolorbox{propriete}[1][]{popbase,colframe=poppurple,before upper={\popboxhead{Propriété}{poppurple}{#1}}}
\newtcolorbox{methode}[1][]{popbase,colframe=popgold,before upper={\popboxhead{Méthode}{popgold}{#1}}}
\newtcolorbox{attention}[1][]{popbase,colframe=poppink,before upper={\popboxhead{Attention}{poppink}{#1}}}
\newtcolorbox{experience}[1][]{popbase,colframe=popblue,colback=popblueL,before upper={\popboxhead{Expérience}{popblue}{#1}}}
% « Le savais-tu ? » : carte violette pleine (contexte, culture, Cameroun)
\newtcolorbox{savaistu}[1][]{popbase,colback=poppurple,colframe=popink,
  fontupper=\popsemi\fontsize{10.4}{14.2}\selectfont\color{white},
  before upper={\poppill{Le savais-tu\,?}{white}{poppurple}\ifx\relax#1\relax\else\hspace{6pt}{\popxbold\color{white}#1}\fi\par\vspace{7pt}}}
% Exercice résolu : énoncé en haut, \tcblower puis la solution sur fond crème
\newcounter{popexo}\newcounter{popexr}
\newtcolorbox{exoresolu}[1][]{popbase,colframe=poporange,colbacklower=popcream,
  segmentation style={popink,line width=1.2pt,dashed},
  before upper={\stepcounter{popexr}\popboxhead{Exercice résolu \thepopexr}{poporange}{#1}},
  before lower={\poppill{Solution}{popink}{popcream}\par\vspace{6pt}}}
% Exercice (non résolu en place) — la correction est dans la partie Corrigés
\newtcolorbox{exercice}[1][]{popbase,colframe=popink,
  before upper={\stepcounter{popexo}\popboxhead{Exercice \thepopexo}{popink}{#1}}}
% Corrigé
\newtcolorbox{corrige}[1][]{popbase,colframe=popgreen,colback=popgreenL,
  before upper={\popboxhead{Corrigé}{popgreen}{#1}}}
% Objectifs / prérequis / bilan
\newtcolorbox{objectifs}{popbase,colframe=popink,colback=poporangeL,
  before upper={\popboxhead{Objectifs de la leçon}{popink}{}}}
\newtcolorbox{prerequis}{popbase,colframe=popink,colback=popblueL,
  before upper={\popboxhead{Prérequis}{popblue}{}}}
\newtcolorbox{bilan}{popbase,colframe=popink,colback=white,boxrule=2.8pt,
  before upper={\popboxhead{Fiche bilan}{poporange}{L'essentiel à connaître pour l'examen}}}
% Figure encadrée avec légende
\newtcolorbox{popfigure}[1][]{enhanced,breakable=false,colback=white,colframe=popline,boxrule=1.4pt,arc=8pt,
  left=10pt,right=10pt,top=10pt,bottom=8pt,before skip=10pt,after skip=12pt,halign=center,
  lower separated=false,halign lower=center,
  fontlower=\popsemi\fontsize{9}{11.5}\selectfont\color{popmuted},
  #1}
\newcommand{\legende}[1]{\par\vspace{6pt}{\popsemi\fontsize{9}{11.5}\selectfont\color{popmuted}#1}}

% ============================================================
% Signature OnBuch+
% ============================================================
\newcommand{\poplogoinline}[1]{{\poplogo\fontsize{#1}{#1}\selectfont\color{popink}OnBuch\color{poporange}+}}
\newcommand{\signatureonbuch}{%
  \begin{tcolorbox}[enhanced,colback=popink,colframe=popink,boxrule=2.2pt,arc=10pt,
      drop shadow=poporange,left=14pt,right=14pt,top=10pt,bottom=10pt,before skip=0pt,after skip=0pt]
    \tikz[baseline=-0.7ex]\node[circle,fill=popgreen,draw=popcream,line width=1.3pt,minimum size=22pt,inner sep=0]{\popcheck{white}};%
    \hspace{9pt}\begin{minipage}[c]{0.62\linewidth}
      {\popxbold\fontsize{12}{13}\selectfont\color{popcream}Signé\ \textbullet\ {\poplogo OnBuch\color{poporange}+}}\par\vspace{2pt}
      {\raggedright\popsemi\fontsize{8}{10.5}\selectfont\color{popline}\MakeUppercase{Équipe pédagogique OnBuch+}\\\MakeUppercase{Conforme au programme officiel MINESEC}\par}
    \end{minipage}\hfill
    {\poplogo\fontsize{22}{22}\selectfont\color{popcream}OnBuch\color{poporange}+}
  \end{tcolorbox}%
}

% ============================================================
% Page de garde
% #1 titre · #2 matière · #3 niveau · #4 module · #5 n° leçon · #6 durée ·
% #7 description / objectifs (texte ou itemize)
% ============================================================
\newcommand{\pagegarde}[7]{%
  \thispagestyle{empty}
  \newgeometry{a4paper,margin=1.7cm,top=1.6cm,bottom=1.4cm}%
  \begin{tikzpicture}[remember picture,overlay]
    \fill[poporange!18] ([xshift=-3.2cm,yshift=-3.6cm]current page.north east) circle (4.6cm);
    \fill[poppurple!14] ([xshift=2.4cm,yshift=7.6cm]current page.south west) circle (3.8cm);
  \end{tikzpicture}%
  {\poplogo\fontsize{30}{30}\selectfont\color{popink}OnBuch\color{poporange}+}\hfill
  \raisebox{0.6ex}{\poppill{Cours signé \textbullet\ Premium}{popink}{popcream}}\par
  \vspace{1pt}\popsquiggle{poppurple}\par
  \vspace{8pt}
  \begin{tcolorbox}[enhanced,colback=poporange,colframe=popink,boxrule=2.8pt,arc=13pt,
      drop shadow=popink,left=18pt,right=18pt,top=12pt,bottom=13pt,after skip=10pt]
    \poppill{#2 \textbullet\ #3}{popink}{popcream}\hspace{5pt}\poppill{#4}{white}{popink}\par\vspace{8pt}
    {\popdisplay\fontsize{14}{14}\selectfont\color{popink}LEÇON #5}\par\vspace{4pt}
    {\raggedright\hyphenpenalty=10000\exhyphenpenalty=10000\popdisplay\fontsize{25}{27}\selectfont\color{popcream}\MakeUppercase{#1}\par}
    \vspace{8pt}
    {\popxbold\fontsize{9.5}{9.5}\selectfont\color{popink}\MakeUppercase{Durée conseillée : #6}}
  \end{tcolorbox}
  \begin{tcolorbox}[enhanced,colback=white,colframe=popink,boxrule=2.2pt,arc=10pt,
      drop shadow=popink,left=16pt,right=16pt,top=10pt,bottom=10pt,after skip=8pt]
    \poppill{Dans cette leçon}{poppurple}{white}\par\vspace{5pt}
    \popsemi\fontsize{9.3}{12.6}\selectfont\color{popink}#7
  \end{tcolorbox}
  \noindent\begin{minipage}[t]{0.487\linewidth}
    \begin{tcolorbox}[enhanced,colback=popgreen,colframe=popink,boxrule=2.2pt,arc=10pt,
        drop shadow=popink,left=11pt,right=11pt,top=10pt,bottom=10pt,height=3.1cm,valign=center]
      \tcbox[colback=white,colframe=popink,boxrule=1.3pt,arc=5pt,boxsep=0pt,nobeforeafter,
        left=3pt,right=3pt,top=3pt,bottom=3pt,valign=center]{\qrcode[height=1.9cm]{https://play.google.com/store/apps/details?id=com.luvvix.onbuch&hl=fr}}%
      \hspace{8pt}\begin{minipage}[c]{0.5\linewidth}
        {\popdisplay\fontsize{11}{12.5}\selectfont\color{white}\MakeUppercase{Télécharge OnBuch}}\par\vspace{4pt}
        {\raggedright\popsemi\fontsize{8.4}{11}\selectfont\color{white}Gratuit \textbullet\ Google Play\\Tuteur IA, annales, concours, résultats\par}
      \end{minipage}
    \end{tcolorbox}
  \end{minipage}\hfill
  \begin{minipage}[t]{0.487\linewidth}
    \begin{tcolorbox}[enhanced,colback=poppurple,colframe=popink,boxrule=2.2pt,arc=10pt,
        drop shadow=popink,left=11pt,right=11pt,top=10pt,bottom=10pt,height=3.1cm,valign=center]
      \tikz[baseline=-0.7ex]\node[circle,fill=white,draw=popink,line width=1.3pt,minimum size=1.6cm,inner sep=0]{\popdisplay\fontsize{24}{24}\selectfont\color{poppurple}+};%
      \hspace{8pt}\begin{minipage}[c]{0.6\linewidth}
        {\popdisplay\fontsize{11}{12.5}\selectfont\color{white}\MakeUppercase{Passe à OnBuch+}}\par\vspace{4pt}
        {\raggedright\popsemi\fontsize{8.4}{11}\selectfont\color{white}Tous les cours signés, Tuteur IA illimité, fiches PDF — onglet OnBuch+ de l'app\par}
      \end{minipage}
    \end{tcolorbox}
  \end{minipage}
  \vspace{8pt}
  \popcontenu
  \vfill
  \signatureonbuch
  \restoregeometry
  \newpage
}

% Rangée « Ce cours contient » (page de garde)
\newcommand{\poptile}[3]{% #1 couleur #2 gros texte #3 légende
  \begin{tcolorbox}[enhanced,colback=#1,colframe=popink,boxrule=2pt,arc=9pt,shadow={2.5pt}{-2.5pt}{0pt}{popink},
      left=6pt,right=6pt,top=9pt,bottom=9pt,halign=center,valign=center,height=1.75cm,width=\linewidth,nobeforeafter]
    {\popdisplay\fontsize{12.5}{13}\selectfont\color{white}\MakeUppercase{#2}}\par\vspace{3pt}
    {\centering\popsemi\fontsize{7.8}{9.5}\selectfont\color{white}#3\par}
  \end{tcolorbox}}
\newcommand{\popcontenu}{%
  \par\noindent{\popxboldls\fontsize{7.5}{7.5}\selectfont\color{popmuted}CE COURS SIGNÉ CONTIENT}\par\vspace{7pt}
  \noindent\begin{minipage}[t]{0.235\linewidth}\poptile{popblue}{Cours}{complet, illustré, pas à pas}\end{minipage}\hfill
  \begin{minipage}[t]{0.235\linewidth}\poptile{poporange}{Méthodes}{et exercices résolus}\end{minipage}\hfill
  \begin{minipage}[t]{0.235\linewidth}\poptile{poppink}{Exercices}{type examen + corrigés}\end{minipage}\hfill
  \begin{minipage}[t]{0.235\linewidth}\poptile{popgreen}{Fiche bilan}{l'essentiel à retenir}\end{minipage}\par}

% Sommaire
\newtcolorbox{sommaire}{enhanced,colback=white,colframe=popink,boxrule=2.2pt,arc=10pt,
  drop shadow=popink,left=18pt,right=18pt,top=13pt,bottom=13pt,before skip=6pt,after skip=20pt,
  before upper={\poppill{Sommaire}{poppurple}{white}\par\vspace{9pt}\popsemi\fontsize{10.6}{14.5}\selectfont\color{popink}}}

% Signature de fin de cours — poussée tout en bas de la dernière page
\newcommand{\findecours}{%
  \par\vfill\nopagebreak
  \begin{center}\popsquiggle{poporange}\end{center}\vspace{4pt}
  \signatureonbuch
}
\AtBeginDocument{\def\llbracket{\mathopen{[\mkern-3.5mu[}}\def\rrbracket{\mathclose{]\mkern-3.5mu]}}} % intervalles d'entiers (glyphe ⟦ absent de la police)
% Glyphes absents de Fira Math : redéfinis à partir de glyphes disponibles
\AtBeginDocument{%
  \def\setminus{\mathbin{\backslash}}\let\smallsetminus\setminus
  \def\vdots{\mathord{\vbox{\baselineskip3.2pt\lineskiplimit0pt\kern4pt\hbox{.}\hbox{.}\hbox{.}}}}%
  \def\ddots{\mathinner{\mkern1mu\raise6.5pt\hbox{.}\mkern2mu\raise3.6pt\hbox{.}\mkern2mu\raise0.7pt\hbox{.}\mkern1mu}}%
  \def\triangle{\mathord{\scalebox{0.9}{$\symup{\Delta}$}}}%
  \def\nabla{\mathord{\raisebox{\depth}{\scalebox{1}[-1]{$\symup{\Delta}$}}}}%
  \def\checkmark{\popcheck{popdark}}%
  \def\star{\mathbin{\ast}}\def\bigstar{\mathord{\ast}}%
  \def\diamond{\mathbin{\scalebox{0.6}{$\Diamond$}}}%
  \def\bigcup{\mathop{\mathchoice{\raisebox{-0.3em}{\scalebox{1.7}{$\cup$}}}{\scalebox{1.25}{$\cup$}}{\cup}{\cup}}\displaylimits}%
  \def\bigcap{\mathop{\mathchoice{\raisebox{-0.3em}{\scalebox{1.7}{$\cap$}}}{\scalebox{1.25}{$\cap$}}{\cap}{\cap}}\displaylimits}%
  \def\lesseqqgtr{\mathrel{\lessgtr}}\def\bigoplus{\mathop{\oplus}\displaylimits}%
}
\providecommand{\ssi}{si et seulement si} % abréviation fréquente
\AtBeginDocument{\providecommand{\wideparen}{\overparen}} % arc de cercle

%% FIGFIX-BEGIN v5 — correctifs globaux des figures (docs/onbuch-generation/tools/figfix)
\usepackage{adjustbox}\usepackage{etoolbox}\tcbuselibrary{raster}
% toute figure TikZ/pgfplots plus large que la ligne est réduite à la largeur disponible
\BeforeBeginEnvironment{tikzpicture}{\begin{adjustbox}{max width=\linewidth}}
\AfterEndEnvironment{tikzpicture}{\end{adjustbox}}
% légendes pgfplots lisibles par-dessus les courbes
\pgfplotsset{every axis/.append style={legend style={fill=white,fill opacity=0.92,text opacity=1,draw=black!15,inner sep=3pt}}}
% étiquettes d'arêtes (syntaxe edge["..."]) avec fond blanc
\tikzset{every edge quotes/.append style={fill=white,inner sep=1.2pt}}
%% FIGFIX-END
\begin{document}

\pagegarde{Leçon 24 — Grammaire : 越来越； ; 一方面……,另一方面……}{Chinois}{1ère A}{Module 3 : Citoyenneté et ouverture au monde}{ZHA24}{2 h}{Cette leçon de grammaire présente deux structures fondamentales du chinois niveau 1ère : 越来越 pour exprimer un changement progressif (« de plus en plus ») et 一方面……另一方面…… pour articuler deux aspects d'une même situation (« d'une part… d'autre part »). Schémas, emplois, comparaison avec le français et entraînement systématique préparent les élèves à les employer sans faute à l'écrit comme à l'oral.
\vspace{4pt}
\begin{multicols}{2}\small
\begin{itemize}
\item À la fin de cette leçon, je sais construire la structure 越来越 + adjectif/verbe psychologique pour exprimer un changement progressif
\item À la fin de cette leçon, je sais que 越来越 ne s'emploie jamais avec 很 ni avec un autre degré devant l'adjectif
\item À la fin de cette leçon, je sais utiliser 越来越 avec des verbes comme 喜欢, 想, 怕 (verbes d'état psychologique)
\item À la fin de cette leçon, je sais construire 一方面……另一方面…… pour présenter deux aspects parallèles d'une situation
\item À la fin de cette leçon, je sais placer correctement 一方面 et 另一方面 par rapport au sujet et au verbe
\item À la fin de cette leçon, je sais comparer ces structures avec « de plus en plus » et « d'une part… d'autre part » en français
\end{itemize}\end{multicols}}

\begin{sommaire}
\begin{multicols}{2}
\begin{enumerate}
\item Observation guidée des exemples
\item Structure 1 : 越来越 — schéma et emplois
\item Structure 2 : 一方面……另一方面……
\item Comparaison avec le français et pièges des francophones
\item Exercices d'application et de transformation
\item Activité d'intégration
\item Méthodes et exercices résolus
\item Exercices
\item Corrigés des exercices
\item Fiche bilan
\end{enumerate}
\end{multicols}
\end{sommaire}

\begin{objectifs}
\begin{itemize}
\item À la fin de cette leçon, je sais construire la structure 越来越 + adjectif/verbe psychologique pour exprimer un changement progressif
\item À la fin de cette leçon, je sais que 越来越 ne s'emploie jamais avec 很 ni avec un autre degré devant l'adjectif
\item À la fin de cette leçon, je sais utiliser 越来越 avec des verbes comme 喜欢, 想, 怕 (verbes d'état psychologique)
\item À la fin de cette leçon, je sais construire 一方面……另一方面…… pour présenter deux aspects parallèles d'une situation
\item À la fin de cette leçon, je sais placer correctement 一方面 et 另一方面 par rapport au sujet et au verbe
\item À la fin de cette leçon, je sais comparer ces structures avec « de plus en plus » et « d'une part… d'autre part » en français
\item À la fin de cette leçon, je sais repérer et corriger les erreurs fréquentes des francophones (很 + 越来越, ordre des mots, oubli de la seconde partie)
\item À la fin de cette leçon, je sais employer les deux structures dans des phrases et un court texte argumentatif
\end{itemize}
\end{objectifs}

\begin{prerequis}
\begin{itemize}
\item Les adjectifs qualificatifs et la négation 不 (Seconde)
\item L'adverbe de degré 很 et sa fonction de liaison sujet-adjectif (Seconde)
\item La structure sujet + verbe + complément et l'ordre de base des mots en chinois
\item Les verbes d'état psychologique : 喜欢, 想, 爱, 怕 (début d'année)
\item Les connecteurs simples 因为……所以…… et 但是 (début d'année)
\end{itemize}
\end{prerequis}

\newpage

\coursec{Observation guidée des exemples}

\courssub{Situation d'accroche : de Yaoundé au marché Mokolo}

À Yaoundé, pendant la saison des pluies, les élèves du lycée remarquent que les embouteillages deviennent \emph{de plus en plus} longs et que le prix des tomates au marché Mokolo monte \emph{de plus en plus} vite. Le matin, le taxi met chaque semaine un peu plus de temps pour arriver à l'école. Comment dire tout cela en chinois ? Comment exprimer ce changement qui avance petit à petit, jour après jour ?

Et quand un élève de Première hésite entre poursuivre ses études à Douala ou partir au Nigeria, il pèse le pour et le contre : \emph{d'une part} Douala est proche de sa famille, \emph{d'autre part} le Nigeria offre d'autres opportunités. Comment construire cette argumentation équilibrée en chinois ?

Ces deux outils, \zh{越来越} (yuè lái yuè, « de plus en plus ») et \zh{一方面……另一方面……} (yì fāngmiàn… lìng yì fāngmiàn…, « d'une part… d'autre part… »), permettent d'exprimer le \cle{changement progressif} et l'\cle{argumentation équilibrée}, deux compétences clés de l'examen. Dans cette première étape, nous n'allons pas encore apprendre la règle : nous allons \textbf{observer} des phrases authentiques, comme des détectives, pour découvrir nous-mêmes le fonctionnement de ces deux structures.

\textbf{Problématique :} comment le chinois exprime-t-il la progression (« de plus en plus ») et l'équilibre entre deux arguments (« d'une part… d'autre part… ») ?

\courssub{Premier corpus : observer 越来越}

Lisons d'abord quatre phrases de la vie quotidienne. Lis-les à voix haute, en faisant attention aux tons, puis observe ce qui se répète.

\begin{textebox}[Phrase 1 — La météo]
天气越来越热了。\\
Tiānqì yuè lái yuè rè le.\\
« Il fait de plus en plus chaud. »
\end{textebox}

\begin{textebox}[Phrase 2 — Les prix au marché]
西红柿越来越贵了。\\
Xīhóngshì yuè lái yuè guì le.\\
« Les tomates sont de plus en plus chères. »
\end{textebox}

\begin{textebox}[Phrase 3 — Les études]
他的中文越来越好了。\\
Tā de Zhōngwén yuè lái yuè hǎo le.\\
« Son chinois devient de meilleur en meilleur. »
\end{textebox}

\begin{exemplebox}[Phrase 4 — Les sentiments]
我越来越喜欢中文了。\\
Wǒ yuè lái yuè xǐhuan Zhōngwén le.\\
« J'aime de plus en plus le chinois. »
\end{exemplebox}

\begin{center}
\begin{tabularx}{\linewidth}{|X|X|}
\hline
\rowcolor{popblueL} Phrase chinoise (caractères + pinyin) & Traduction française \\
\hline
天气\textbf{越来越}热了。 Tiānqì \textbf{yuè lái yuè} rè le. & Il fait \emph{de plus en plus} chaud. \\
\hline
西红柿\textbf{越来越}贵了。 Xīhóngshì \textbf{yuè lái yuè} guì le. & Les tomates sont \emph{de plus en plus} chères. \\
\hline
他的中文\textbf{越来越}好了。 Tā de Zhōngwén \textbf{yuè lái yuè} hǎo le. & Son chinois devient \emph{de meilleur en meilleur}. \\
\hline
我\textbf{越来越}喜欢中文了。 Wǒ \textbf{yuè lái yuè} xǐhuan Zhōngwén le. & J'aime \emph{de plus en plus} le chinois. \\
\hline
\end{tabularx}
\end{center}

\courssub{Questions d'observation sur 越来越}

Avant toute règle, réponds à ces questions en regardant le tableau ci-dessus :

\begin{enumerate}
\item Quel groupe de mots se répète dans les quatre phrases ? Souligne-le.
\item Où se place-t-il dans la phrase : avant ou après le sujet ? avant ou après l'adjectif ?
\item Quel mot se trouve \textbf{juste après} \zh{越来越} dans chaque phrase ? Quelle est sa nature : un nom, un adjectif, un verbe ?
\item Quel petit mot se trouve souvent à la fin de la phrase ?
\item Comment le français traduit-il cette structure ? Trouve au moins deux formulations possibles.
\end{enumerate}

\begin{methode}[Méthode d'observation d'une structure grammaticale]
Pour observer une structure nouvelle, procède toujours en trois étapes :
\begin{enumerate}
\item \textbf{Souligner} la structure qui se répète dans tous les exemples (ici \zh{越来越}).
\item \textbf{Encadrer} le mot qui la suit immédiatement et identifier sa nature grammaticale (ici : 热 rè « chaud », 贵 guì « cher », 好 hǎo « bon », 喜欢 xǐhuan « aimer » — des adjectifs ou un verbe d'état).
\item \textbf{Comparer} avec la traduction française : quelle expression française correspond à la structure chinoise ? Y a-t-il un mot chinois qui n'a pas d'équivalent en français, ou l'inverse ?
\end{enumerate}
Cette méthode « souligner — encadrer — comparer » te servira pour toutes les leçons de grammaire de l'année.
\end{methode}

\courssub{Hypothèses des élèves sur 越来越}

Après l'observation, les élèves formulent généralement les hypothèses suivantes. Vérifie si les tiennes sont proches :

\begin{itemize}
\item \textbf{Hypothèse 1 :} \zh{越来越} signifie « de plus en plus » et exprime un changement qui progresse avec le temps.
\item \textbf{Hypothèse 2 :} la structure se place \textbf{avant} l'adjectif ou le verbe, jamais après.
\item \textbf{Hypothèse 3 :} le mot qui suit est toujours un \cle{adjectif} (热, 贵, 好) ou un \cle{verbe d'état / verbe psychologique} (喜欢 « aimer », 想 « avoir envie de / manquer à », 怕 « craindre »), jamais un nom.
\item \textbf{Hypothèse 4 :} le \zh{了} final semble marquer que le changement est déjà en cours, qu'on le constate maintenant.
\end{itemize}

\begin{attention}[Piège repéré dès l'observation : l'interdiction de 很]
En français, on peut dire « de plus en plus \emph{très} cher ». En chinois, c'est impossible : on ne met \textbf{jamais} \zh{很} (hěn, « très ») entre \zh{越来越} et l'adjectif. On dit \zh{越来越贵} (de plus en plus cher), mais jamais \zh{越来越很贵}. Pourquoi ? Parce que \zh{越来越} contient déjà l'idée d'intensité qui grandit ; ajouter \zh{很} ferait double emploi. Retiens aussi que \zh{很} devant un adjectif isolé est souvent juste un « lien » neutre, pas un vrai « très ».
\end{attention}

\courssub{Tableau de vocabulaire : lexique du corpus 1}

\begin{center}
\begin{tabularx}{\linewidth}{|X|X|X|X|}
\hline
\rowcolor{popblueL} Caractères & Pinyin & Nature & Traduction \\
\hline
天气 & tiānqì & nom & temps, météo \\
\hline
越来越 & yuè lái yuè & structure & de plus en plus \\
\hline
热 & rè & adj. & chaud \\
\hline
西红柿 & xīhóngshì & nom & tomate \\
\hline
贵 & guì & adj. & cher \\
\hline
中文 & Zhōngwén & nom & chinois (langue) \\
\hline
好 & hǎo & adj. & bon, bien \\
\hline
喜欢 & xǐhuan & verbe & aimer, apprécier \\
\hline
了 & le & part. & marqueur d'aspect (changement, réalisation) \\
\hline
\end{tabularx}
\end{center}

\courssub{Second corpus : observer 一方面……另一方面……}

Passons maintenant aux phrases d'argumentation. Observe comment le locuteur présente deux aspects d'une même situation.

\begin{textebox}[Phrase 5 — L'apprentissage du chinois]
学中文一方面有意思，另一方面有用。\\
Xué Zhōngwén yì fāngmiàn yǒu yìsi, lìng yì fāngmiàn yǒuyòng.\\
« Apprendre le chinois, d'une part c'est intéressant, d'autre part c'est utile. »
\end{textebox}

\begin{textebox}[Phrase 6 — Le choix des études]
一方面我想在杜阿拉学习，另一方面我想去尼日利亚。\\
Yì fāngmiàn wǒ xiǎng zài Dù'ālā xuéxí, lìng yì fāngmiàn wǒ xiǎng qù Nírìlìyà.\\
« D'une part je voudrais étudier à Douala, d'autre part je voudrais aller au Nigeria. »
\end{textebox}

\begin{textebox}[Phrase 7 — La vie en ville]
城市生活一方面很方便，另一方面太贵了。\\
Chéngshì shēnghuó yì fāngmiàn hěn fāngbiàn, lìng yì fāngmiàn tài guì le.\\
« La vie en ville, d'une part elle est très pratique, d'autre part elle est trop chère. »
\end{textebox}

\begin{center}
\begin{tabularx}{\linewidth}{|X|X|}
\hline
\rowcolor{popgreenL} Phrase chinoise (caractères + pinyin) & Traduction française \\
\hline
学中文\textbf{一方面}有意思，\textbf{另一方面}有用。 Xué Zhōngwén \textbf{yì fāngmiàn} yǒu yìsi, \textbf{lìng yì fāngmiàn} yǒuyòng. & Apprendre le chinois, \emph{d'une part} c'est intéressant, \emph{d'autre part} c'est utile. \\
\hline
\textbf{一方面}我想在杜阿拉学习，\textbf{另一方面}我想去尼日利亚。 \textbf{Yì fāngmiàn} wǒ xiǎng zài Dù'ālā xuéxí, \textbf{lìng yì fāngmiàn} wǒ xiǎng qù Nírìlìyà. & \emph{D'une part} je voudrais étudier à Douala, \emph{d'autre part} je voudrais aller au Nigeria. \\
\hline
城市生活\textbf{一方面}很方便，\textbf{另一方面}太贵了。 Chéngshì shēnghuó \textbf{yì fāngmiàn} hěn fāngbiàn, \textbf{lìng yì fāngmiàn} tài guì le. & La vie en ville, \emph{d'une part} elle est pratique, \emph{d'autre part} elle est trop chère. \\
\hline
\end{tabularx}
\end{center}

\courssub{Questions d'observation sur 一方面……另一方面……}

\begin{enumerate}
\item Quels mots se répètent dans les trois phrases ? Combien de « morceaux » la structure comporte-t-elle ?
\item Où se placent ces morceaux : au début de chaque partie de phrase, ou à la fin ?
\item Que remarques-tu sur le premier caractère du deuxième morceau : est-ce le même que dans le premier ? (\zh{一} yì « un » contre \zh{另} lìng « autre ».)
\item Les deux parties de la phrase parlent-elles du même sujet ou de sujets différents ?
\item Quelle expression française correspond à cette structure ?
\end{enumerate}

\textbf{Hypothèses des élèves :}

\begin{itemize}
\item \textbf{Hypothèse 1 :} la structure a deux parties obligatoires : \zh{一方面} au début du premier argument, \zh{另一方面} au début du second.
\item \textbf{Hypothèse 2 :} elle sert à présenter deux aspects, deux avantages ou deux possibilités d'une \textbf{même} situation, comme une balance à deux plateaux.
\item \textbf{Hypothèse 3 :} les deux arguments peuvent être de même valeur (deux avantages : intéressant + utile) ou opposés (pratique + cher), comme « pour » et « contre ».
\item \textbf{Hypothèse 4 :} une virgule chinoise \zh{，} sépare les deux parties.
\end{itemize}

\courssub{Tableau de vocabulaire : lexique du corpus 2}

\begin{center}
\begin{tabularx}{\linewidth}{|X|X|X|X|}
\hline
\rowcolor{popgreenL} Caractères & Pinyin & Nature & Traduction \\
\hline
学 & xué & verbe & apprendre, étudier \\
\hline
一方面 & yì fāngmiàn & locution & d'une part \\
\hline
另 & lìng & adj. & autre \\
\hline
另一方面 & lìng yì fāngmiàn & locution & d'autre part \\
\hline
有意思 & yǒu yìsi & adj. & intéressant \\
\hline
有用 & yǒuyòng & adj. & utile \\
\hline
杜阿拉 & Dù'ālā & n.p. & Douala \\
\hline
尼日利亚 & Nírìlìyà & n.p. & Nigeria \\
\hline
城市 & chéngshì & nom & ville \\
\hline
生活 & shēnghuó & nom/verbe & vie, vivre \\
\hline
方便 & fāngbiàn & adj. & pratique, commode \\
\hline
太 & tài & adv. & trop \\
\hline
\end{tabularx}
\end{center}

\begin{aretenir}[Ce que l'observation nous a déjà appris]
\begin{itemize}
\item \zh{越来越} + adjectif ou verbe d'état (+ \zh{了}) = « de plus en plus… ». \textbf{Jamais} de \zh{很} à l'intérieur.
\item \zh{一方面}…，\zh{另一方面}… = « d'une part…, d'autre part… ». Deux blocs, un même thème, une virgule entre les deux.
\item Le français possède des équivalents presque mot à mot : la difficulté n'est donc pas le sens, mais la \textbf{place exacte} des mots dans la phrase chinoise.
\end{itemize}
\end{aretenir}

\courssub{Schéma visuel des deux structures observées}

\begin{popfigure}
\begin{tikzpicture}[node distance=0.6cm, every node/.style={font=\small, align=center}]
\node[draw=popblue, fill=popblueL, rounded corners, minimum width=2.0cm, minimum height=0.8cm, align=center] (s1){Sujet\\我 / 天气};
\node[draw=poporange, fill=poporangeL, rounded corners, minimum width=2.4cm, minimum height=0.8cm, right=of s1, align=center] (yl){越来越\\yuè lái yuè};
\node[draw=popgreen, fill=popgreenL, rounded corners, minimum width=2.6cm, minimum height=0.8cm, right=of yl, align=center] (adj){Adjectif / Verbe d'état\\热 / 喜欢};
\node[draw=poppurple, fill=poppurpleL, rounded corners, minimum width=1.2cm, minimum height=0.8cm, right=of adj] (le) {了};
\draw[-{Stealth}, thick, popdark] (s1) -- (yl);
\draw[-{Stealth}, thick, popdark] (yl) -- (adj);
\draw[-{Stealth}, thick, popdark] (adj) -- (le);
\node[below=0.3cm of yl, text=popmuted, font=\footnotesize] {« de plus en plus »};
\node[below=0.3cm of le, text=popmuted, font=\footnotesize] {changement constaté};
\end{tikzpicture}
\legende{Schéma 1 : la chaîne 越来越 — Sujet + 越来越 + Adjectif/Verbe d'état + 了.}
\end{popfigure}

\begin{popfigure}
\begin{tikzpicture}[node distance=0.8cm, every node/.style={font=\small, align=center}]
\node[draw=popdark, fill=popcream, rounded corners, minimum width=5cm, minimum height=0.8cm, align=center] (theme){Thème commun\\学中文 / 城市生活};
\node[draw=popblue, fill=popblueL, rounded corners, minimum width=4.5cm, minimum height=1.0cm, below left=1.0cm and 0.5cm of theme, align=center] (a){一方面 + Argument 1\\有意思 / 很方便};
\node[draw=poporange, fill=poporangeL, rounded corners, minimum width=4.5cm, minimum height=1.0cm, below right=1.0cm and 0.5cm of theme, align=center] (b){另一方面 + Argument 2\\有用 / 太贵了};
\draw[-{Stealth}, thick, popdark] (theme.south west) -- (a.north);
\draw[-{Stealth}, thick, popdark] (theme.south east) -- (b.north);
\draw[{Stealth}-{Stealth}, thick, poppurple, dashed] (a.east) -- (b.west) node[midway, below, font=\footnotesize, text=poppurple] {balance / virgule};
\end{tikzpicture}
\legende{Schéma 2 : la balance 一方面……另一方面…… — deux arguments autour d'un même thème.}
\end{popfigure}

\courssub{Exercice résolu d'observation}

\begin{exoresolu}[Repérage des structures]
Énoncé — Dans les phrases suivantes, souligne la structure étudiée, encadre le mot qui la suit, puis traduis en français : (1) 雅温得的雨越来越大了。 (2) 踢足球一方面很好玩，另一方面很累。
\tcblower
Solution détaillée — (1) La structure est \zh{越来越} (Yǎwēndé de yǔ \textbf{yuè lái yuè} dà le). Le mot encadré qui suit est \zh{大} (dà, « grand / fort »), un adjectif ; le \zh{了} final marque le changement constaté. Traduction : « La pluie à Yaoundé devient de plus en plus forte. » (2) La structure est \zh{一方面……另一方面……} (Tī zúqiú \textbf{yì fāngmiàn} hěn hǎowán, \textbf{lìng yì fāngmiàn} hěn lèi). Les deux arguments encadrés sont \zh{很好玩} (« très amusant ») et \zh{很累} (« très fatigant ») ; le thème commun est « jouer au football ». Traduction : « Jouer au football, d'une part c'est très amusant, d'autre part c'est très fatigant. » Remarque : ici \zh{很} est possible car il accompagne les adjectifs dans une phrase descriptive ordinaire ; c'est seulement avec \zh{越来越} qu'il est interdit.
\end{exoresolu}

\begin{exercice}[À toi d'observer]
Observe les trois phrases suivantes et réponds aux questions : (1) 喀麦隆的足球越来越有名了。 (2) 学汉语一方面难，另一方面有意思。 (3) 我越来越想念中国了。
\begin{enumerate}
\item Souligne les structures et encadre le mot qui suit \zh{越来越}.
\item Quelle est la nature des mots encadrés ?
\item Traduis les trois phrases en français.
\item Dans la phrase (2), les deux arguments sont-ils du même côté (pour/pour) ou opposés (pour/contre) ?
\end{enumerate}
\end{exercice}

\begin{corrige}[Exercice « À toi d'observer »]
(1) Structures : \zh{越来越} dans les phrases 1 et 3 ; \zh{一方面……另一方面……} dans la phrase 2. Mots encadrés : \zh{有名} (yǒumíng, « célèbre ») et \zh{想念} (xiǎngniàn, « manquer à, regretter »). (2) \zh{有名} est un adjectif, \zh{想念} est un verbe d'état (verbe psychologique) : les deux sont acceptés après \zh{越来越}. (3) Traductions : « Le football camerounais devient de plus en plus célèbre. » / « Apprendre le chinois, d'une part c'est difficile, d'autre part c'est intéressant. » / « La Chine me manque de plus en plus. » (4) Les arguments sont opposés : un « contre » (难, difficile) et un « pour » (有意思, intéressant) — c'est le cas typique de la balance argumentaire.
\end{corrige}

\coursec{Structure 1 : 越来越 — Règle formelle, emplois et exceptions}

\begin{propriete}[Structure 越来越 : schéma et emplois]
\textbf{Schéma de base :}
\begin{center}
\textbf{Sujet + 越来越 + Adjectif / Verbe d'état + (了)}
\end{center}

\textbf{Emplois principaux :}
\begin{enumerate}
\item \textbf{Changement progressif d'un état ou d'une qualité} (adjectif) : \zh{天气越来越冷了} (Il fait de plus en plus froid).
\item \textbf{Évolution d'un sentiment ou d'une disposition mentale} (verbe psychologique) : \zh{我越来越喜欢这儿了} (J'aime de plus en plus cet endroit), \zh{他越来越想家了} (Il a de plus en plus le mal du pays).
\item \textbf{Constatation d'un fait accompli en cours} : la particule \zh{了} en fin de phrase est \textbf{quasi systématique} pour marquer que le changement est réel, observable maintenant (nouvel état).
\end{enumerate}

\textbf{Ordre des mots strict :} \zh{越来越} se place \textbf{immédiatement avant} l'adjectif ou le verbe. Rien ne s'intercale (pas de \zh{很}, pas de \zh{真}, pas de complément).
\end{propriete}

\begin{exemplebox}[Exemples variés avec 越来越]
\begin{tabularx}{\linewidth}{|X|X|X|}
\hline
\rowcolor{popblueL} Chinois (caractères + pinyin) & Structure & Français \\
\hline
物价越来越高了。 Wùjià yuè lái yuè gāo le. & Suj + 越来越 + Adj + 了 & Les prix sont de plus en plus élevés. \\
\hline
妹妹越来越漂亮了。 Mèimei yuè lái yuè piàoliang le. & Suj + 越来越 + Adj + 了 & Ma petite sœur est de plus en plus jolie. \\
\hline
我们越来越想去北京了。 Wǒmen yuè lái yuè xiǎng qù Běijīng le. & Suj + 越来越 + Verbe (+ compl.) + 了 & Nous avons de plus en plus envie d'aller à Pékin. \\
\hline
这个问题越来越难了。 Zhège wèntí yuè lái yuè nán le. & Suj + 越来越 + Adj + 了 & Ce problème devient de plus en plus difficile. \\
\hline
\end{tabularx}
\end{exemplebox}

\begin{attention}[Erreurs fréquentes des francophones avec 越来越]
\begin{enumerate}
\item \textbf{Insérer 很 / 真 / 非常 :} \zh{*越来越很贵} (Faux). \rightarrow \textbf{Correct :} \zh{越来越贵}.
\item \textbf{Oublier le 了 final :} \zh{*天气越来越热} (Phrase incomplète, sonne « suspendu »). \rightarrow \textbf{Correct :} \zh{天气越来越热了}.
\item \textbf{Mettre 越来越 après l'adjectif :} \zh{*热越来越了} (Faux, ordre inverse). \rightarrow \textbf{Correct :} \zh{越来越热了}.
\item \textbf{L'utiliser avec un nom :} \zh{*越来越钱} (Faux). \rightarrow \textbf{Correct :} \zh{钱越来越多了} (L'argent devient de plus en plus abondant).
\item \textbf{Confondre 越来越 et 越……越…… (structure corrélative, niveau HSK 4) :} \zh{越来越} est autonome ; \zh{越……越……} lie deux verbes/adjectifs (ex: \zh{越吃越胖}). Ne les mélangez pas.
\end{enumerate}
\end{attention}

\coursec{Structure 2 : 一方面……另一方面…… — Règle formelle et argumentation}

\begin{propriete}[Structure 一方面……另一方面…… : schéma et emplois]
\textbf{Schéma de base :}
\begin{center}
\textbf{(Sujet / Thème +) 一方面 + Proposition 1 ， 另一方面 + Proposition 2}
\end{center}

\textbf{Points clés :}
\begin{itemize}
\item \zh{一} (yì) et \zh{另} (lìng) : \zh{一} = « un » (premier aspect) ; \zh{另} = « autre » (deuxième aspect, distinct).
\item \zh{方面} (fāngmiàn) = « aspect, facette, côté ».
\item La virgule \zh{，} est \textbf{obligatoire} pour séparer les deux propositions.
\item Le sujet/thème commun peut être énoncé une seule fois au début (Phrase 5, 7) ou répété dans chaque proposition (Phrase 6).
\end{itemize}

\textbf{Emplois :}
\begin{enumerate}
\item \textbf{Énumération d'avantages / aspects positifs} (pour/pour) : \zh{这个手机一方面便宜，另一方面好用}.
\item \textbf{Énumération d'inconvénients / aspects négatifs} (contre/contre) : \zh{这份工作一方面累，另一方面工资低}.
\item \textbf{Argumentation équilibrée « pour / contre »} (le plus fréquent à l'examen) : \zh{出国留学一方面开阔眼界，另一方面花费大}.
\end{enumerate}
\end{propriete}

\begin{exemplebox}[Variantes de position du sujet]
\begin{tabularx}{\linewidth}{|X|X|}
\hline
\rowcolor{popgreenL} Type & Exemple (caractères + pinyin + français) \\
\hline
Sujet unique au début & 这件衣服\textbf{一方面}颜色好看，\textbf{另一方面}价格合理。\\
& Zhè jiàn yīfu \textbf{yì fāngmiàn} yánsè hǎokàn, \textbf{lìng yì fāngmiàn} jiàgé hélǐ.\\
& « Ce vêtement, d'une part la couleur est belle, d'autre part le prix est raisonnable. » \\
\hline
Sujet répété (insistance) & \textbf{一方面}我喜欢喀麦隆，\textbf{另一方面}我想去中国。\\
& \textbf{Yì fāngmiàn} wǒ xǐhuan Kāmàilóng, \textbf{lìng yì fāngmiàn} wǒ xiǎng qù Zhōngguó.\\
& « D'une part j'aime le Cameroun, d'autre part je veux aller en Chine. » \\
\hline
\end{tabularx}
\end{exemplebox}

\begin{attention}[Pièges et erreurs fréquentes avec 一方面……另一方面……]
\begin{enumerate}
\item \textbf{Oublier 另 (lìng) :} écrire \zh{一方面……一方面……} (Faux). Le second marqueur est \textbf{另}方面.
\item \textbf{Oublier la virgule :} les deux blocs doivent être séparés par \zh{，}.
\item \textbf{Mélanger les connecteurs :} ne pas écrire \zh{一方面……但是……} ou \zh{虽然……另一方面……}. La structure est figée : \zh{一方面……另一方面……}.
\item \textbf{Deux sujets sans lien :} les deux arguments doivent porter sur \textbf{le même thème}. \zh{*一方面我喜欢吃饭，另一方面他去学校} (Incohérent).
\item \textbf{Traduction littérale « un côté… l'autre côté » :} en français, on dit « d'une part… d'autre part » ou « à la fois… et… » (si même valence), pas « un aspect… l'autre aspect » (trop lourd).
\end{enumerate}
\end{attention}

\coursec{Comparaison systématique français / chinois et pièges croisés}

\begin{propriete}[Tableau comparatif : De plus en plus / D'une part… d'autre part]
\begin{center}
\begin{tabularx}{\linewidth}{|X|X|X|}
\hline
\rowcolor{poppurpleL} Notion & Français & Chinois (Structure + Exemple) \\
\hline
Progression & de plus en plus + adj/verbe & \zh{越来越} + Adj/V + \zh{了} \\
& \emph{Il fait de plus en plus chaud.} & \zh{天气越来越热了。} \\
\hline
Intensité forte (statique) & très + adj & \zh{很} / \zh{真} / \zh{非常} + Adj (\textbf{sans} 了 obligatoire) \\
& \emph{Il fait très chaud.} & \zh{天气很热。} \\
\hline
Argumentation équilibrée & d'une part…, d'autre part… & \zh{一方面}…, \zh{另一方面}… \\
& \emph{D'une part c'est cher, d'autre part c'est beau.} & \zh{一方面很贵，另一方面很漂亮。} \\
\hline
Concession / Opposition & bien que… / mais… & \zh{虽然}…, \zh{但是}… (Structure différente, voir leçon ultérieure) \\
\hline
\end{tabularx}
\end{center}
\end{propriete}

\begin{methode}[Traduire « de plus en plus » et « d'une part… d'autre part » sans erreur]
\textbf{Pour « de plus en plus » (越来越) :}
\begin{enumerate}
\item Identifier le sujet.
\item Identifier l'adjectif ou le verbe d'état (aimer, vouloir, craindre, manquer).
\item Construire : Sujet + \zh{越来越} + Adj/V + \zh{了}.
\item \textbf{Vérifier :} pas de \zh{很} / \zh{真} entre \zh{越来越} et l'adjectif.
\end{enumerate}

\textbf{Pour « d'une part… d'autre part » (一方面……另一方面……) :}
\begin{enumerate}
\item Identifier le thème commun.
\item Rédiger l'argument 1 (phrase complète).
\item Rédiger l'argument 2 (phrase complète).
\item Assembler : [Thème +] \zh{一方面} + Arg1 + \zh{，} + \zh{另一方面} + Arg2.
\item \textbf{Vérifier :} présence de \zh{另} (pas \zh{一} deux fois), présence de la virgule \zh{，}.
\end{enumerate}
\end{methode}

\courssub{Exercices d'application et de transformation}

\begin{exercice}[Transformation : Français $\rightarrow$ Chinois]
Traduis en chinois (caractères + pinyin) en utilisant \zh{越来越} ou \zh{一方面……另一方面……}.
\begin{enumerate}
\item La vie à Yaoundé devient de plus en plus chère.
\item Étudier le chinois, d'une part c'est difficile, d'autre part c'est très utile.
\item Ma petite sœur aime de plus en plus le football.
\item Ce portable, d'une part il est léger, d'autre part il n'est pas cher.
\item Le temps (météo) devient de plus en plus froid.
\item Aller à l'université à Douala : d'une part c'est près de la maison, d'autre part les frais sont élevés.
\end{enumerate}
\end{exercice}

\begin{exercice}[Transformation : Restructuration de phrases]
Réécris les phrases suivantes en insérant la structure demandée. Attention à l'ordre des mots et au \zh{了}.
\begin{enumerate}
\item 这本书有意思。 (avec \zh{越来越} + \zh{了})
\item 学外语有用，学外语开阔眼界。 (avec \zh{一方面……另一方面……}, sujet unique)
\item 现在的学生忙。 (avec \zh{越来越} + \zh{了})
\item 住在农村安静，住在农村不方便。 (avec \zh{一方面……另一方面……}, sujet répété)
\end{enumerate}
\end{exercice}

\begin{exercice}[Production guidée : Petit paragraphe argumenté]
Rédige un court paragraphe (4-5 phrases) sur le thème : \emph{« Étudier à l'étranger ou rester au Cameroun ? »}.
\begin{itemize}
\item Utilise \textbf{au moins une fois} \zh{越来越} (ex: \zh{越来越多的学生想出国}).
\item Utilise \textbf{au moins une fois} \zh{一方面……另一方面……} pour peser le pour et le contre.
\item Vocabulaire suggéré : \zh{出国} (chūguó), \zh{留学} (liúxué), \zh{机会} (jīhuì), \zh{家人} (jiārén), \zh{费用} (fèiyòng), \zh{决定} (juédìng).
\end{itemize}
\end{exercice}

\begin{corrige}[Exercices d'application et de transformation]

\textbf{Exercice 1 (Traduction) :}
\begin{enumerate}
\item 雅温得的生活越来越贵了。 (Yǎwēndé de shēnghuó yuè lái yuè guì le.)
\item 学中文一方面难，另一方面很有用。 (Xué Zhōngwén yì fāngmiàn nán, lìng yì fāngmiàn hěn yǒuyòng.)
\item 妹妹越来越喜欢足球了。 (Mèimei yuè lái yuè xǐhuan zúqiú le.)
\item 这个手机一方面很轻，另一方面不贵。 (Zhège shǒujī yì fāngmiàn hěn qīng, lìng yì fāngmiàn bù guì.)
\item 天气越来越冷了。 (Tiānqì yuè lái yuè lěng le.)
\item 在杜阿拉上大学一方面离家近，另一方面费用高。 (Zài Dù'ālā shàng dàxué yì fāngmiàn lí jiā jìn, lìng yì fāngmiàn fèiyòng gāo.)
\end{enumerate}

\textbf{Exercice 2 (Restructuration) :}
\begin{enumerate}
\item 这本书\textbf{越来越有意思了}。 (Zhè běn shū \textbf{yuè lái yuè yǒu yìsi le}.)
\item 学外语\textbf{一方面}有用，\textbf{另一方面}开阔眼界。 (Xué wàiyǔ \textbf{yì fāngmiàn} yǒuyòng, \textbf{lìng yì fāngmiàn} kāikuò yǎnjiè.)
\item 现在的学生\textbf{越来越忙了}。 (Xiànzài de xuésheng \textbf{yuè lái yuè máng le}.)
\item \textbf{一方面}住在农村安静，\textbf{另一方面}住在农村不方便。 (\textbf{Yì fāngmiàn} zhù zài nóngcūn ānjìng, \textbf{lìng yì fāngmiàn} zhù zài nóngcūn bù fāngbiàn.)
\end{enumerate}

\textbf{Exercice 3 (Production guidée — proposition de corrigé) :}
\zh{越来越多的喀麦隆学生想出国留学。一方面出国有更多机会，另一方面费用太高了。一方面离开家人不容易，另一方面在国内也有好大学。我要仔细决定。}
(Pinyin : Yuè lái yuè duō de Kāmàilóng xuésheng xiǎng chūguó liúxué. Yì fāngmiàn chūguó yǒu gèng duō jīhuì, lìng yì fāngmiàn fèiyòng tài gāo le. Yì fāngmiàn líkāi jiārén bù róngyì, lìng yì fāngmiàn zài guónèi yě yǒu hǎo dàxué. Wǒ yào zǐxì juédìng.)
Traduction : « De plus en plus d'étudiants camerounais veulent partir étudier à l'étranger. D'une part, partir offre plus d'opportunités, d'autre part les frais sont trop élevés. D'une part, quitter sa famille n'est pas facile, d'autre part il y a aussi de bonnes universités dans le pays. Je dois bien réfléchir pour décider. »
\end{corrige}

\begin{savaistu}[Le saviez-vous ? — Étymologie et culture]
\textbf{越 (yuè) :} À l'origine, ce caractère représente une hache (戉) avec un son. Il signifie « dépasser, franchir, aller au-delà ». Dans \zh{越来越}, l'idée est « dépasser (le stade actuel), venir (vers l'avenir), dépasser (encore) » : une progression sans fin.

\textbf{另 (lìng) :} Composé de \zh{立} (debout) et \zh{口} (bouche/ordre) + \zh{力} (force) en bas (forme ancienne). Cela évoque « un autre ordre, une autre force » $\rightarrow$ « autre, différent ».

\textbf{Contexte camerounais / chinois :} En Chine, pour exprimer une hésitation polie, on utilise souvent \zh{一方面……另一方面……} pour refuser une invitation sans froisser : \zh{一方面我很想去，另一方面最近太忙了} (« D'une part j'aimerais bien y aller, d'autre part je suis trop occupé ces temps-ci »). C'est l'équivalent culturel du « Je veux bien, mais… » camerounais.
\end{savaistu}

\begin{experience}[Activité orale : « Le grand débat de la Première »]
\textbf{Consigne :} En binôme, tirez un sujet au sort. L'élève A défend le « pour » avec \zh{一方面}, l'élève B le « contre » avec \zh{另一方面}. Puis synthétisez à deux voix avec la structure complète.
\textbf{Sujets :}
\begin{enumerate}
\item Avoir un smartphone au lycée.
\item Manger au restaurant vs manger à la cantine.
\item Apprendre l'anglais vs apprendre le chinois.
\item Habiter en ville (Yaoundé/Douala) vs habiter au village.
\end{enumerate}
\textbf{Exemple de production attendue :}
A : \zh{一方面手机方便查资料。}
B : \zh{另一方面上课容易分心。}
Ensemble : \zh{带手机来学校一方面方便查资料，另一方面容易分心。}
\end{experience}

\coursec{Leçon 24 — Grammaire : 越来越 ; 一方面……另一方面……}

\begin{textebox}[Situation de communication : Au lycée de Yaoundé]
\zh{—— 同学们，最近雅温得的天气怎么样？}\\
\emph{Tóngxuémen, zuìjìn Yǎwēndé de tiānqì zěnme yàng?}\\
« — Camarades, comment est le temps à Yaoundé ces temps-ci ? »\\

\zh{—— 老师，越来越热了！一方面大家想游泳，另一方面要复习考试，真难办。}\\
\emph{Lǎoshī, yuè lái yuè rè le! Yī fāngmiàn dàjiā xiǎng yóuyǒng, lìng yī fāngmiàn yào fùxí kǎoshì, zhēn nán bàn.}\\
« — Professeur, il fait de plus en plus chaud ! D'un côté, tout le monde veut nager, de l'autre, il faut réviser les examens, c'est vraiment embêtant. »\\

\zh{—— 是啊，中文越学越有意思，但是作业也越来越多。一方面想放松，另一方面不敢懈怠。}\\
\emph{Shì a, Zhōngwén yuè xué yuè yǒu yìsi, dànshì zuòyè yě yuè lái yuè duō. Yī fāngmiàn xiǎng fàngsōng, lìng yī fāngmiàn bù gǎn xièdài.}\\
« — Oui, plus on étudie le chinois, plus c'est intéressant, mais les devoirs deviennent de plus en plus nombreux. D'un côté, on veut se détendre, de l'autre, on n'ose pas relâcher ses efforts. »
\end{textebox}

\coursec{Observation guidée des exemples}

Lisez attentivement le dialogue ci-dessus. Repérez deux structures grammaticales récurrentes :
\begin{enumerate}
\item \zh{越来越热}、\zh{越来越多}、\zh{越学越有意思} : elles expriment une \textbf{évolution}, un changement de degré dans le temps ou en fonction d'une action.
\item \zh{一方面……另一方面……}、\zh{一方面……另一方面……} : elles servent à \textbf{présenter deux facettes} d'une même situation, souvent contrastées ou complémentaires.
\end{enumerate}

Ces deux structures sont indispensables pour exprimer une opinion nuancée, décrire une tendance sociale ou argumenter dans une rédaction (niveau HSK 3-4 / Baccalauréat). Nous les étudierons l'une après l'autre.

\coursec{Structure 1 : 越来越 — Schéma et emplois}

\courssub{Le sens général : « de plus en plus »}

\begin{definition}[越来越 yuè lái yuè]
\zh{越来越} est un adverbe de degré qui exprime une \cle{évolution progressive dans le temps} : le degré d'une qualité, d'un état ou d'un sentiment \textbf{augmente (ou diminue) peu à peu}. Il correspond au français « de plus en plus » ; sa forme négative \zh{越来越不} correspond à « de moins en moins ».
\end{definition}

L'idée essentielle est celle d'un \cle{changement} : la situation d'aujourd'hui n'est plus celle d'hier, et celle de demain sera encore différente. \zh{越来越} ne décrit jamais un état fixe, mais une dynamique.

\begin{exemplebox}[Observer avant de comprendre]
\zh{雅温得的天气越来越热了。}\\
\emph{Yǎwēndé de tiānqì yuè lái yuè rè le.}\\
« Le temps à Yaoundé devient de plus en plus chaud. »

\zh{中文越来越有意思。}\\
\emph{Zhōngwén yuè lái yuè yǒu yìsi.}\\
« Le chinois devient de plus en plus intéressant. »

\zh{我越来越想家。}\\
\emph{Wǒ yuè lái yuè xiǎng jiā.}\\
« Ma famille me manque de plus en plus. »
\end{exemplebox}

\courssub{Le schéma de la structure}

\begin{propriete}[Schéma général de 越来越]
\textbf{Sujet + 越来越 + Adjectif / Verbe psychologique (+ 了)}

\begin{itemize}
\item \zh{越来越} se place \textbf{devant} l'adjectif ou le verbe, comme un adverbe ordinaire.
\item L'adjectif / verbe qui suit est \textbf{nu} : jamais de \zh{很}, \zh{非常}, \zh{太} devant lui.
\item La particule \zh{了} finale est \textbf{facultative} mais très fréquente : elle souligne que le changement est déjà perceptible (constat).
\end{itemize}
\end{propriete}

\begin{popfigure}
\begin{tikzpicture}[
  bloc/.style={draw=popblue, fill=popblueL, rounded corners=3pt, minimum height=0.9cm, align=center, font=\small},
  opt/.style={draw=poporange, fill=poporangeL, rounded corners=3pt, minimum height=0.9cm, align=center, font=\small},
  fleche/.style={-{Stealth[length=3mm]}, thick, popdark}
]
\node[bloc, align=center] (s) at (0,0){Sujet\\ \zh{天气}};
\node[bloc, fill=poppurpleL, draw=poppurple] (y) at (3.2,0) {\zh{越来越}};
\node[bloc, fill=popgreenL, draw=popgreen, align=center] (a) at (6.4,0){Adjectif / V. psych.\\ \zh{热} / \zh{喜欢}};
\node[opt] (l) at (9.8,0) {(\zh{了})};
\draw[fleche] (s) -- (y);
\draw[fleche] (y) -- (a);
\draw[fleche, dashed] (a) -- (l);
\draw[fleche, poporange, very thick, -{Stealth[length=4mm]}] (1.2,-1.4) -- (8.8,-0.6)
  node[midway, below, font=\small\itshape, poporange] {le degré monte avec le temps};
\end{tikzpicture}
\legende{Schéma de la structure 越来越 : la flèche ascendante symbolise la progression du degré.}
\end{popfigure}

\begin{popfigure}
\begin{tikzpicture}[
  pt/.style={circle, draw=popblue, fill=popblueL, inner sep=2pt, minimum size=0.55cm, align=center, font=\small},
  fleche/.style={-{Stealth[length=3mm]}, thick, popdark}
]
\draw[fleche] (-0.5,0) -- (10.5,0);
\node[pt] (p1) at (1.5,0) {};
\node[pt] (p2) at (5,0) {};
\node[pt] (p3) at (8.5,0) {};
\node[below=2pt of p1, font=\small] {hier (passé)};
\node[below=2pt of p2, font=\small] {aujourd'hui (présent)};
\node[below=2pt of p3, font=\small] {demain (futur)};
\node[above=2pt of p1, font=\small] {un peu chaud \zh{热}};
\node[above=2pt of p2, font=\small] {plus chaud \zh{更热}};
\node[above=2pt of p3, font=\small] {encore plus chaud \zh{更热}};
\draw[-{Stealth[length=3mm]}, very thick, poporange] (1.5,1.1) to[bend left=15] (8.5,1.9);
\node[font=\small\itshape, poporange] at (5,2.3) {intensité croissante : \zh{越来越热}};
\end{tikzpicture}
\legende{Frise des temps : 越来越 décrit une intensité qui grandit du passé vers le futur.}
\end{popfigure}

\courssub{Emploi 1 : 越来越 + adjectif}

C'est l'emploi le plus fréquent. L'adjectif exprime une qualité dont le degré évolue : \zh{热} (chaud), \zh{好} (bon), \zh{贵} (cher), \zh{多} (nombreux), \zh{漂亮} (beau), \zh{难} (difficile), \zh{方便} (pratique), \zh{堵车} (embouteillé — verbe d'état).

\begin{exemplebox}[越来越 + adjectif]
\zh{东西越来越贵了。}\\
\emph{Dōngxi yuè lái yuè guì le.}\\
« Les choses deviennent de plus en plus chères. »

\zh{杜阿拉的车越来越多了。}\\
\emph{Dùālā de chē yuè lái yuè duō le.}\\
« Il y a de plus en plus de voitures à Douala. »

\zh{他的汉语越来越好了。}\\
\emph{Tā de Hànyǔ yuè lái yuè hǎo le.}\\
« Son chinois devient de meilleur en meilleur. »

\zh{这个城市越来越漂亮了。}\\
\emph{Zhège chéngshì yuè lái yuè piàoliang le.}\\
« Cette ville devient de plus en plus belle. »

\zh{上班的路越来越堵车了。}\\
\emph{Shàngbān de lù yuè lái yuè dǔchē le.}\\
« La route pour aller au travail est de plus en plus embouteillée. »
\end{exemplebox}

\begin{attention}[Piège n° 1 : jamais de 很 / 非常 / 太 après 越来越]
En français, on dit « de plus en plus \emph{très} cher ». En chinois, c'est impossible : \zh{越来越} porte déjà l'idée de degré croissant. On ne peut donc \textbf{jamais} mettre \zh{很}, \zh{非常} ou \zh{太} devant l'adjectif.

FAUX : *\zh{越来越很热} \quad *\zh{越来越非常贵} \quad *\zh{越来越太多}

CORRECT : \zh{越来越热} \quad \zh{越来越贵} \quad \zh{越来越多}

Même interdiction pour la négation : on dit \zh{越来越不好}, jamais *\zh{越来越很不好}.
\end{attention}

\courssub{Emploi 2 : 越来越 + verbe psychologique}

\zh{越来越} peut précéder les \cle{verbes d'état psychologique} : verbes de sentiment, de désir, de préférence, de cognition. Les plus utiles à votre niveau : \zh{喜欢} (aimer), \zh{爱} (aimer fort), \zh{想} (avoir envie de), \zh{怕} (craindre), \zh{希望} (espérer), \zh{担心} (s'inquiéter), \zh{了解} (connaître/comprendre), \zh{讨厌} (détester).

\begin{exemplebox}[越来越 + verbe psychologique]
\zh{他越来越想学中文。}\\
\emph{Tā yuè lái yuè xiǎng xué Zhōngwén.}\\
« Il a de plus en plus envie d'apprendre le chinois. »

\zh{我越来越喜欢喀麦隆的足球了。}\\
\emph{Wǒ yuè lái yuè xǐhuan Kāmàilóng de zúqiú le.}\\
« J'aime de plus en plus le football camerounais. »

\zh{她越来越怕考试。}\\
\emph{Tā yuè lái yuè pà kǎoshì.}\\
« Elle craint de plus en plus les examens. »

\zh{我们越来越了解中国了。}\\
\emph{Wǒmen yuè lái yuè liǎojiě Zhōngguó le.}\\
« Nous connaissons de mieux en mieux la Chine. »

\zh{妈妈越来越担心我的身体。}\\
\emph{Māma yuè lái yuè dānxīn wǒ de shēntǐ.}\\
« Maman s'inquiète de plus en plus pour ma santé. »
\end{exemplebox}

\begin{attention}[Quels verbes sont possibles ? Test du \zh{很}]
Seuls les verbes qui expriment un \textbf{état ou un sentiment} acceptent \zh{越来越}. Les verbes d'action ponctuelle ou de processus pur sont impossibles : on ne peut pas dire *\zh{越来越吃} ni *\zh{越来越去} ni *\zh{越来越写}.

\textbf{Test simple} : si le verbe peut être modifié par \zh{很} (\zh{很喜欢}, \zh{很想}, \zh{很怕}, \zh{很了解}), alors il accepte \zh{越来越}. Si *\zh{很吃}, *\zh{很去} sont faux, alors *\zh{越来越吃}, *\zh{越来越去} sont faux.
\end{attention}

\courssub{Le 了 final : constater le changement}

\begin{definition}[了 final après 越来越]
Le \zh{了} placé en fin de phrase est la \cle{particule de constat d'une nouvelle situation}. Il signale que le locuteur \emph{observe} le changement : « je constate que c'est devenu ainsi ». Il est facultatif mais très naturel, surtout à l'oral.
\end{definition}

Comparez :

\begin{exemplebox}[Avec ou sans 了]
\zh{天气越来越热。} \emph{Tiānqì yuè lái yuè rè.} — « Le temps devient de plus en plus chaud. » (constat neutre, tendance générale)

\zh{天气越来越热了。} \emph{Tiānqì yuè lái yuè rè le.} — « Le temps est devenu de plus en plus chaud (je le remarque). » (constat vécu, changement déjà perceptible)
\end{exemplebox}

\begin{attention}[Place du 了 : en toute fin de phrase !]
Il se met \textbf{après l'adjectif} et \textbf{après un éventuel complément d'objet} : \zh{我越来越喜欢这个城市了。} (\emph{Wǒ yuè lái yuè xǐhuan zhège chéngshì le.}). Ne le placez \textbf{jamais} entre \zh{越来越} et l'adjectif/verbe : *\zh{越来越了热} est faux.
\end{attention}

\courssub{La négation : 越来越不 + adjectif / verbe psych.}

Pour dire « de moins en moins », on insère \zh{不} entre \zh{越来越} et l'adjectif (ou le verbe psychologique) : \textbf{Sujet + 越来越不 + Adjectif / V. psych. (+ 了)}.

\begin{exemplebox}[越来越不 + adjectif / verbe]
\zh{他的身体越来越不好了。}\\
\emph{Tā de shēntǐ yuè lái yuè bù hǎo le.}\\
« Sa santé va de moins en moins bien. »

\zh{路越来越不方便了。}\\
\emph{Lù yuè lái yuè bù fāngbiàn le.}\\
« La route devient de moins en moins pratique. »

\zh{他越来越不喜欢吃肉了。}\\
\emph{Tā yuè lái yuè bù xǐhuan chī ròu le.}\\
« Il aime de moins en moins manger de la viande. »

\zh{学生们越来越不想上课了。}\\
\emph{Xuéshengmen yuè lái yuè bù xiǎng shàngkè le.}\\
« Les élèves ont de moins en moins envie d'aller en cours. »
\end{exemplebox}

\begin{aretenir}[Nuance : \zh{越来越不好} vs \zh{越来越坏}]
\zh{越来越不好} = « de moins en moins bon » (la qualité positive baisse).
\zh{越来越坏} = « de pire en pire » (la qualité négative augmente).
Les deux sont proches mais \zh{越来越不 + adj. positif} est souvent plus poli ou plus précis pour décrire une dégradation.
\end{aretenir}

\courssub{Comparaison avec le français (Structure 1)}

\begin{center}
\begin{tabularx}{\linewidth}{|X|X|X|}
\hline
\rowcolor{popblueL} Structure & Exemple (caractères + pinyin) & Traduction \\
\hline
Sujet + 越来越 + Adj. (+ 了) & \zh{东西越来越贵了。} Dōngxi yuè lái yuè guì le. & Les choses deviennent de plus en plus chères. \\
\hline
Sujet + 越来越 + V. psych. (+ 了) & \zh{他越来越想学中文。} Tā yuè lái yuè xiǎng xué Zhōngwén. & Il a de plus en plus envie d'apprendre le chinois. \\
\hline
Sujet + 越来越不 + Adj. (+ 了) & \zh{天气越来越不好了。} Tiānqì yuè lái yuè bù hǎo le. & Le temps devient de moins en moins bon. \\
\hline
越来越 + Adj. + 了 (sujet contextuel) & \zh{越来越贵了！} Yuè lái yuè guì le! & C'est de plus en plus cher ! \\
\hline
\end{tabularx}
\end{center}

Trois différences essentielles avec le français :
\begin{enumerate}
\item \textbf{Position figée} : En français, « de plus en plus » est mobile. En chinois, \zh{越来越} est \textbf{rivé} juste avant l'adjectif/verbe. *\zh{越来越，他来} est faux.
\item \textbf{Comparatifs irréguliers} : Le français a « de mieux en mieux », « de pire en pire ». Le chinois garde \zh{越来越} + adjectif simple : \zh{越来越好} = « de mieux en mieux », \zh{越来越坏} = « de pire en pire ».
\item \textbf{Le 了 final} : Aucun équivalent direct en français. Il faut le penser « constat d'une situation nouvelle » et non le traduire mot à mot.
\end{enumerate}

\courssub{Pour aller plus loin : 越 + Verbe + 越 + Adjectif (Corrélatif)}

\begin{aretenir}[Variante d'examen : 越V越Adj.]
Structure : \textbf{Sujet + 越 + Verbe + 越 + Adjectif}. Sens : « Plus on fait X, plus c'est Y » (corrélation entre action et résultat). Très prisée en rédaction et à l'oral (HSK 4).

\zh{中文越学越有意思。} \emph{Zhōngwén yuè xué yuè yǒu yìsi.} — « Plus on étudie le chinois, plus c'est intéressant. »

\zh{这本书越看越喜欢。} \emph{Zhè běn shū yuè kàn yuè xǐhuan.} — « Plus on lit ce livre, plus on l'aime. »

\zh{他越跑越快。} \emph{Tā yuè pǎo yuè kuài.} — « Plus il court, plus il va vite. »

\zh{越吃越胖。} \emph{Yuè chī yuè pàng.} — « Plus on mange, plus on grossit. »
\end{aretenir}

\begin{savaistu}[Étymologie du caractère 越]
\zh{越} (yuè) signifie à l'origine « dépasser, franchir, aller au-delà » (radical \zh{走} « marcher » + phonétique). On le retrouve dans \zh{越过} (franchir) et \zh{越南} (Yuènán, « Au-delà du Sud », le Vietnam). Dans \zh{越来越}, la répétition dessine une \textbf{montée continue} : chaque \zh{越} franchit le niveau précédent. C'est une image très chinoise de la progression : pas un saut brutal, mais une série de petits dépassements. Au Cameroun, on dirait : « Petit à petit, l'oiseau fait son nid » — \zh{越来越} raconte exactement ce mouvement patient.
\end{savaistu}

\courssub{Exercice résolu : Construction avec 越来越}

\begin{exoresolu}[Mise en application immédiate]
Énoncé : Traduisez en chinois les phrases suivantes en utilisant \zh{越来越} (ou \zh{越…越…} pour la dernière).\\
1. Mon chinois devient de meilleur en mieux.\\
2. Elle a de moins en moins envie de sortir.\\
3. Les élèves de Première A aiment de plus en plus ce cours.\\
4. La vie à Yaoundé devient de plus en plus chère (je le constate).\\
5. Plus on pratique, plus on devient fort.

\tcblower
Solution détaillée :

1. « Mon chinois » = \zh{我的汉语} ; « de meilleur en mieux » = \zh{越来越好}. Constat $\rightarrow$ \zh{了}.\\
\zh{我的汉语越来越好了。} \emph{Wǒ de Hànyǔ yuè lái yuè hǎo le.}

2. « De moins en moins » = \zh{越来越不} ; « avoir envie de sortir » = \zh{想出去}.\\
\zh{她越来越不想出去了。} \emph{Tā yuè lái yuè bù xiǎng chūqù le.}

3. Sujet : \zh{高一A班的学生们} (Première A) ; verbe \zh{喜欢} ; objet \zh{这门课} ; \zh{了} en fin.\\
\zh{高一A班的学生们越来越喜欢这门课了。} \emph{Gāo yī A bān de xuéshengmen yuè lái yuè xǐhuan zhè mén kè le.}

4. « La vie » = \zh{生活} ; « cher » = \zh{贵} ; constat $\rightarrow$ \zh{了}.\\
\zh{雅温得的生活越来越贵了。} \emph{Yǎwēndé de shēnghuó yuè lái yuè guì le.}

5. Corrélation action/résultat $\rightarrow$ \zh{越…越…}. « Pratiquer » = \zh{练习} ou \zh{练} ; « fort / fort en... » = \zh{厉害} ou \zh{棒}.\\
\zh{越练习越厉害。} \emph{Yuè liànxí yuè lìhai.} (Ou : \zh{越练越棒。})

Vérification : aucun \zh{很} après \zh{越来越}, \zh{了} en fin de phrase, \zh{不} bien placé.
\end{exoresolu}

\begin{attention}[Récapitulatif des erreurs fréquentes des francophones sur 越来越]
\begin{itemize}
\item *\zh{越来越很热} $\rightarrow$ \zh{越来越热} (jamais d'adverbe de degré après \zh{越来越}).
\item *\zh{我越来越了喜欢中文} $\rightarrow$ \zh{我越来越喜欢中文了} (\zh{了} en fin de phrase).
\item *\zh{越来越，天气热了} $\rightarrow$ \zh{天气越来越热了} (\zh{越来越} ne s'emploie jamais seul en tête de phrase).
\item *\zh{他越来越不很聪明} $\rightarrow$ \zh{他越来越不聪明} (la négation n'accepte pas \zh{很}).
\item *\zh{越来越吃得多} pour « il mange de plus en plus » $\rightarrow$ Structure incorrecte. Dire : \zh{他吃得越来越多} (complément de degré) ou \zh{他越来越爱吃} (verbe psychologique).
\end{itemize}
\end{attention}

\begin{aretenir}[L'essentiel de la structure 越来越]
\begin{itemize}
\item Schéma : \textbf{Sujet + 越来越 + Adjectif / Verbe psychologique (+ 了)}.
\item Sens : « de plus en plus », changement graduel dans le temps.
\item Négation : \zh{越来越不} + Adj/V.psych. = « de moins en moins ».
\item Interdiction absolue : \zh{很} / \zh{非常} / \zh{太} après \zh{越来越}.
\item \zh{了} final = constat du changement, facultatif mais naturel.
\item Variante corrélative (examen) : \zh{越 + V + 越 + Adj.} « plus on X, plus c'est Y ».
\end{itemize}
\end{aretenir}

\coursec{Structure 2 : 一方面……另一方面…… — Schéma et emplois}

\courssub{Le sens général : « d'un côté… de l'autre… »}

\begin{definition}[一方面……另一方面…… yī fāngmiàn… lìng yī fāngmiàn…]
Cette structure coordonnée sert à \textbf{présenter deux aspects d'une même réalité}, deux points de vue, deux raisons, ou deux conséquences qui coexistent. Elle correspond au français « d'un côté… de l'autre… », « à la fois… et… », « tout en… » ou « si… toutefois… ».
\end{definition}

Elle permet de \cle{nuancer son propos}, d'éviter les jugements trop tranchés, ce qui est très valorisé à l'oral (exposé, débat) et à l'écrit (rédaction argumentative).

\courssub{Le schéma de la structure}

\begin{propriete}[Schéma général]
\textbf{一方面，分句1； 另一方面，分句2。}

\begin{itemize}
\item \zh{一方面} (litt. « un côté/face ») introduit le premier aspect.
\item \zh{另一方面} (litt. « l'autre côté/face ») introduit le second aspect.
\item La ponctuation standard : \textbf{virgule après chaque connecteur}, \textbf{point-virgule} (ou point) entre les deux divisions principales.
\item Le sujet des deux divisions peut être le même ou différent.
\end{itemize}
\end{propriete}

\begin{popfigure}
\begin{tikzpicture}[
  bloc/.style={draw=popblue, fill=popblueL, rounded corners=3pt, minimum height=0.9cm, align=center, font=\small, text width=2.8cm},
  fleche/.style={-{Stealth[length=3mm]}, thick, popdark},
  brace/.style={decorate, decoration={brace, amplitude=5pt}, thick, poppurple}
]
\node[bloc, align=center] (y1) at (0,0){\zh{一方面}，\\ Sujet + Prédicat 1};
\node[bloc, align=center] (y2) at (7,0){\zh{另一方面}，\\ Sujet + Prédicat 2};
\draw[fleche] (y1) -- node[above, font=\small] {;} (y2);
\draw[brace] (y1.north west) -- (y1.north east) node[midway, above=5pt, font=\small\itshape, poppurple] {Aspect A};
\draw[brace] (y2.north west) -- (y2.north east) node[midway, above=5pt, font=\small\itshape, poppurple] {Aspect B};
\node[font=\small\itshape, popdark] at (3.5, -1.5) {Deux facettes d'un même sujet / situation};
\end{tikzpicture}
\legende{Structure binaire : 一方面 A ; 另一方面 B.}
\end{popfigure}

\courssub{Emploi 1 : Deux raisons / causes coexistantes}

Le sujet est souvent le même. On explique un fait par deux facteurs.

\begin{exemplebox}[Deux raisons]
\zh{他不来了。一方面他生病了，另一方面他没钱。}\\
\emph{Tā bù lái le. Yī fāngmiàn tā bìng le, lìng yī fāngmiàn tā méi qián.}\\
« Il ne vient plus. D'un côté, il est malade ; de l'autre, il n'a pas d'argent. »

\zh{雅温得交通堵塞。一方面车太多，另一方面路太窄。}\\
\emph{Yǎwēndé jiāotōng dǔsè. Yī fāngmiàn chē tài duō, lìng yī fāngmiàn lù tài zhǎi.}\\
« Les embouteillages à Yaoundé. D'un côté, trop de voitures ; de l'autre, les routes sont trop étroites. »
\end{exemplebox}

\courssub{Emploi 2 : Avantages et inconvénients / Pour et Contre}

Très utilisé pour donner un avis équilibré (ex: vie en ville vs campagne, études vs travail, technologie).

\begin{exemplebox}[Avantages / Inconvénients]
\zh{在杜阿拉住很方便。一方面商场多，另一方面交通堵车严重。}\\
\emph{Zài Dùālā zhù hěn fāngbiàn. Yī fāngmiàn shāngchǎng duō, lìng yī fāngmiàn jiāotōng dǔchē yánzhòng.}\\
« Habiter à Douala est pratique. D'un côté, il y a beaucoup de centres commerciaux ; de l'autre, les embouteillages sont graves. »

\zh{学中文很有用。一方面能找好工作，另一方面很难坚持。}\\
\emph{Xué Zhōngwén hěn yǒuyòng. Yī fāngmiàn néng zhǎo hǎo gōngzuò, lìng yī fāngmiàn hěn nán jiānchí.}\\
« Apprendre le chinois est utile. D'un côté, on peut trouver un bon travail ; de l'autre, c'est difficile à persévérer. »
\end{exemplebox}

\courssub{Emploi 3 : Sentiments ou actions contradictoires / simultanés}

Le sujet est souvent \zh{我} ou \zh{我们}. On exprime un dilemme intérieur.

\begin{exemplebox}[Sentiments contradictoires]
\zh{我想去中国。一方面想看长城，另一方面怕坐飞机。}\\
\emph{Wǒ xiǎng qù Zhōngguó. Yī fāngmiàn xiǎng kàn Chángchéng, lìng yī fāngmiàn pà zuò fēijī.}\\
« Je veux aller en Chine. D'un côté, je veux voir la Grande Muraille ; de l'autre, j'ai peur de prendre l'avion. »

\zh{这个周末我很纠结。一方面想复习，另一方面想去看足球赛。}\\
\emph{Zhège zhōumò wǒ hěn jiūjié. Yī fāngmiàn xiǎng fùxí, lìng yī fāngmiàn xiǎng qù kàn zúqiúsài.}\\
« Ce week-end, je suis tiraillé. D'un côté, je veux réviser ; de l'autre, je veux aller voir le match de foot. »
\end{exemplebox}

\courssub{Variante : 一面……一面…… (Actions simultanées)}

\begin{attention}[Ne pas confondre : \zh{一方面…另一方面…} vs \zh{一面…一面…}]
\begin{itemize}
\item \zh{一方面……另一方面……} = \textbf{Logique / Argumentation} : « D'un point de vue A… ; du point de vue B… ». Les deux clauses sont des \emph{propositions} (sujet + verbe).
\item \zh{一面……一面……} (yī miàn… yī miàn…) = \textbf{Simultanéité physique} : « Tout en faisant A, on fait B ». Sujet unique obligatoire. Souvent avec verbes d'action.
\end{itemize}
\end{attention}

\begin{exemplebox}[一面……一面…… (Simultanéité)]
\zh{他一面听音乐，一面做作业。}\\
\emph{Tā yī miàn tīng yīnyuè, yī miàn zuò zuòyè.}\\
« Il écoute de la musique \textbf{tout en} faisant ses devoirs. » (Deux actions en même temps).

\zh{妈妈一面做饭，一面打电话。}\\
\emph{Māma yī miàn zuòfàn, yī miàn dǎ diànhuà.}\\
« Maman cuisine tout en téléphonant. »
\end{exemplebox}

\begin{aretenir}[Différence clé]
\begin{center}
\begin{tabularx}{\linewidth}{|X|X|X|}
\hline
\rowcolor{popblueL} Structure & Sens & Exemple \\
\hline
\zh{一方面……另一方面……} & Argumentation, nuances, raisons, pour/contre & \zh{一方面想去，另一方面没钱。} (Dilemme / Raisons) \\
\hline
\zh{一面……一面……} & Actions simultanées (physiques) & \zh{一面吃饭，一面看电视。} (Manger EN REGARDANT la télé) \\
\hline
\end{tabularx}
\end{center}
\end{aretenir}

\courssub{Comparaison avec le français (Structure 2)}

\begin{center}
\begin{tabularx}{\linewidth}{|X|X|X|}
\hline
\rowcolor{popblueL} Structure chinoise & Exemple & Traduction française \\
\hline
一方面 A，另一方面 B & \zh{一方面累，另一方面开心。} & D'un côté fatigué, de l'autre content. / À la fois fatigué et content. \\
\hline
一方面 A，另一方面 B & \zh{一方面贵，另一方面好。} & C'est cher, mais c'est bon. / Si c'est cher, c'est pourtant bon. \\
\hline
一面 A，一面 B & \zh{他一面跑一面喊。} & Il court \textbf{tout en} criant. / Il court \textbf{en} criant. \\
\hline
\end{tabularx}
\end{center}

Pièges francophones :
\begin{enumerate}
\item \textbf{Traduire « d'un côté » par \zh{一面}} : *\zh{一面贵，一面好} $\rightarrow$ Faux (cela veut dire « c'est cher en même temps que c'est bon » physiquement). Correct : \zh{一方面贵，另一方面好}.
\item \textbf{Oublier la virgule après le connecteur} : *\zh{一方面他累了另一方面他高兴} $\rightarrow$ Correct : \zh{一方面，他累了；另一方面，他高兴。}
\item \textbf{Confondre « bien que » (concession) et « d'un côté »} : « Bien qu'il soit cher, je l'achète » = \zh{虽然贵，但我买} (concession), pas \zh{一方面贵，另一方面我买}.
\end{enumerate}

\coursec{Synthèse comparative et pièges croisés}

Les deux structures de cette leçon servent à \textbf{complexifier le discours} :
\begin{itemize}
\item \zh{越来越} : complexifie sur l'\textbf{axe temporel} (évolution, dynamique).
\item \zh{一方面……另一方面……} : complexifie sur l'\textbf{axe logique} (nuance, dialectique).
\end{itemize}

\begin{methode}[Comment choisir la bonne structure ?]
\begin{enumerate}
\item \textbf{Question :} Est-ce que je décris un \textbf{changement dans le temps} (hier vs aujourd'hui vs demain) ? $\rightarrow$ \zh{越来越} (ou \zh{越…越…}).
\item \textbf{Question :} Est-ce que je présente \textbf{deux arguments, deux raisons, deux sentiments contraires} coexistant \textbf{maintenant} ? $\rightarrow$ \zh{一方面……另一方面……}.
\item \textbf{Question :} Est-ce que je décris \textbf{deux actions physiques faites en même temps} par la même personne ? $\rightarrow$ \zh{一面……一面……}.
\end{enumerate}
\end{methode}

\begin{exemplebox}[Mélange des deux structures dans un texte cohérent]
\zh{现在的高中生压力很大。\textbf{一方面}，课程\textbf{越来越多}，\textbf{越来越难}；\textbf{另一方面}，家长希望\textbf{越来越高}。学生\textbf{一方面}想放松，\textbf{另一方面}不敢懈怠。}\\
\emph{Xiànzài de gāozhōngshēng yālì hěn dà. Yī fāngmiàn, kèchéng yuè lái yuè duō, yuè lái yuè nán; lìng yī fāngmiàn, jiāzhǎng xīwàng yuè lái yuè gāo. Xuésheng yī fāngmiàn xiǎng fàngsōng, lìng yī fāngmiàn bù gǎn xièdài.}\\
« Les lycéens d'aujourd'hui subissent une forte pression. D'un côté, les cours sont de plus en plus nombreux et de plus en plus difficiles ; de l'autre, les attentes des parents sont de plus en plus élevées. Les élèves, d'un côté, veulent se détendre ; de l'autre, ils n'osent pas relâcher leurs efforts. »
\end{exemplebox}

\coursec{Exercices d'application et de transformation}

\begin{exercice}[Exercice 1 : Complétez avec 越来越 / 越来越不 / 越…越…]
Mettez la structure adaptée (attention à la négation et au 了 final).
\begin{enumerate}
\item 这道菜 \_\_\_\_\_\_ 好吃，我要多吃一点。 (délicieux)
\item 最近我 \_\_\_\_\_\_ 想睡觉，可能太累了。 (avoir envie de dormir)
\item 中文 \_\_\_\_\_\_ 学 \_\_\_\_\_\_ 简单。 (corrélatif : plus on étudie, plus c'est simple)
\item 杜阿拉的房租 \_\_\_\_\_\_ 贵 \_\_\_\_\_\_，真买不起。 (de plus en plus cher, constat)
\item 妹妹 \_\_\_\_\_\_ 不 喜欢 吃 蔬菜 了。 (de moins en moins aimer)
\end{enumerate}

\end{exercice}

\begin{exercice}[Exercice 2 : Reliez les deux phrases par 一方面……另一方面……]
\begin{enumerate}
\item 这个手机很漂亮。 / 这个手机太贵了。
\item 我想去旅行。 / 我没有假期。
\item 住在农村空气好。 / 住在农村不方便 (交通/achats).
\item 学中文能了解中国文化。 / 学中文汉字很难写。
\end{enumerate}

\end{exercice}

\begin{exercice}[Exercice 3 : Traduction thème (Français $\rightarrow$ Chinois)]
\begin{enumerate}
\item Le climat au Cameroun devient de plus en plus sec (constat).
\item D'un côté, ce restaurant est bon marché ; de l'autre, c'est délicieux.
\item Plus tu parles chinois, plus tu progresses.
\item Elle a de moins en moins peur de parler en public.
\item Je suis occupé : d'un côté, je dois réviser ; de l'autre, je dois aider ma mère.
\end{enumerate}

\end{exercice}

\begin{exercice}[Exercice 4 : Transformation (Réécriture)]
Transformez les phrases suivantes en utilisant la structure indiquée entre parenthèses.
\begin{enumerate}
\item 他的中文好多了。 (用 \zh{越来越})
\item 虽然累，但是开心。 (用 \zh{一方面……另一方面……})
\item 他边走边看手机。 (用 \zh{一面……一面……})
\item 现在的年轻人不爱看书了。 (用 \zh{越来越不})
\item 这件衣服又贵又旧。 (用 \zh{一方面……另一方面……} pour exprimer le contraste « cher mais vieux/mauvais »)
\end{enumerate}

\end{exercice}

\begin{exercice}[Exercice 5 : Expression écrite guidée (Sujet de type Bac)]
\textbf{Sujet :} « La vie des jeunes à Yaoundé (ou Douala / Buea) change rapidement. »
Rédigez un court paragraphe (60-80 caractères chinois) en utilisant \textbf{au moins une fois} \zh{越来越} et \textbf{au moins une fois} \zh{一方面……另一方面……}.
\textit{Indices :} 手机, 网络, 学习, 压力, 娱乐, 方便, 近视, 交友.

\end{exercice}

\courssub{Corrigés des exercices}



\begin{corrige}[Exercice 1]
\begin{enumerate}
\item \zh{越来越} (adj. positif, progression) $\rightarrow$ \zh{这道菜越来越好吃，我要多吃一点。}
\item \zh{越来越} (verbe psychologique \zh{想}) $\rightarrow$ \zh{最近我越来越想睡觉，可能太累了。}
\item \zh{越} 学 \zh{越} 简单 (corrélatif) $\rightarrow$ \zh{中文学越越简单。}
\item \zh{越来越} 贵 \zh{了} (adj., constat) $\rightarrow$ \zh{杜阿拉的房租越来越贵了，真买不起。}
\item \zh{越来越不} (négation verbe psychologique) $\rightarrow$ \zh{妹妹越来越不喜欢吃蔬菜了。}
\end{enumerate}
\end{corrige}

\begin{corrige}[Exercice 2]
\begin{enumerate}
\item \zh{这个手机一方面很漂亮，另一方面太贵了。}
\item \zh{我想去旅行。一方面想去，另一方面没有假期。} (Sujet \zh{我} omis 2e fois) ou \zh{一方面我想去旅行，另一方面我没有假期。}
\item \zh{住在农村一方面空气好，另一方面不方便。}
\item \zh{学中文一方面能了解中国文化，另一方面汉字很难写。}
\end{enumerate}
\end{corrige}

\begin{corrige}[Exercice 3]
\begin{enumerate}
\item \zh{喀麦隆的气候越来越干燥了。} (\zh{干燥} gānzào = sec)
\item \zh{这家饭馆一方面便宜，另一方面很好吃。} (\zh{饭馆} fànguǎn = restaurant ; \zh{便宜} piányi = bon marché)
\item \zh{中文越说越进步。} ou \zh{越练习中文越进步。} (\zh{进步} jìnbù = progresser)
\item \zh{她越来越不怕公开演讲了。} (\zh{公开演讲} gōngkāi yǎnjiǎng = parler en public)
\item \zh{我很忙。一方面要复习，另一方面要帮妈妈。} (\zh{忙} máng = occupé ; \zh{帮} bāng = aider)
\end{enumerate}
\end{corrige}

\begin{corrige}[Exercice 4]
\begin{enumerate}
\item \zh{他的中文越来越好了。}
\item \zh{一方面他累了，另一方面他很开心。} (Ou : \zh{一方面累，另一方面开心。})
\item \zh{他一面走，一面看手机。}
\item \zh{现在的年轻人越来越不爱看书了。}
\item \zh{这件衣服一方面贵，另一方面旧。} (Ou : \zh{不好看} / \zh{质量不好} selon le sens voulu pour « vieux »).
\end{enumerate}
\end{corrige}

\begin{corrige}[Exercice 5 — Proposition de rédaction]
\zh{现在雅温得的年轻人生活变化大。一方面，手机和网络越来越方便，交友、娱乐都容易了；另一方面，学习压力越来越大，近视的同学也越来越多。我们一方面要用好网络，另一方面要保护眼睛，多运动。}\\
\emph{Xiànzài Yǎwēndé de niánqīng rén shēnghuó biànhuà dà. Yī fāngmiàn, shǒujī hé wǎngluò yuè lái yuè fāngbiàn, jiāoyǒu, yúlè dōu róngyì le; lìng yī fāngmiàn, xuéxí yālì yuè lái yuè dà, jìnshì de tóngxué yě yuè lái yuè duō. Wǒmen yī fāngmiàn yào yòng hǎo wǎngluò, lìng yī fāngmiàn yào bǎohù yǎnjing, duō yùndòng.}

\textbf{Analyse :} 3 occurrences de \zh{越来越} (\zh{越来越方便}、\zh{越来越大}、\zh{越来越多}), 2 occurrences de \zh{一方面……另一方面……}. Vocabulaire HSK 3-4. Structure cohérente.
\end{corrige}

\begin{savaistu}[Culture : 越 et la pensée dialectique chinoise]
L'association \zh{一方面……另一方面……} reflète une tendance culturelle chinoise à la \cle{dialectique} (辩证法) : les choses ne sont pas noires ou blanches, elles contiennent des contradictions (矛盾) qui coexistent et s'interpénètrent. Dire « c'est bien MAIS c'est cher » (\zh{虽然好，但是贵}) tranche ; dire « \zh{一方面好，另一方面贵} » maintient les deux vérités en tension. C'est une compétence clé pour l'argumentation au Bac et dans le supérieur (HSK 5-6).
\end{savaistu}

\begin{aretenir}[Bilan de la Leçon 24 — Les deux structures clés]
\begin{center}
\begin{tabularx}{\linewidth}{|X|X|X|X|}
\hline
\rowcolor{popblueL} Structure & Schéma & Sens clé & Piège majeur \\
\hline
\zh{越来越} & S + 越来越 + Adj/V.psych (+ 了) & Évolution temporelle : « de plus en plus » & *\zh{越来越很Adj} ; *\zh{了} au milieu \\
\hline
\zh{越V越Adj} & S + 越 + V + 越 + Adj & Corrélation : « plus on V, plus Adj » & Confondre avec \zh{越来越} (axe unique) \\
\hline
\zh{一方面…另一方面…} & 一方面，S+V ; 另一方面，S+V & Nuance / Argumentation : « d'un côté… de l'autre » & Confondre avec \zh{一面…一面…} (simultanéité) \\
\hline
\zh{一面…一面…} & S + 一面 + V1，一面 + V2 & Simultanéité : « tout en V1, V2 » & Sujet obligatoire unique ; verbes d'action \\
\hline
\end{tabularx}
\end{center}
\end{aretenir}

\coursec{Structure 2 : 一方面……另一方面……}

Après avoir appris à exprimer une évolution progressive avec \zh{越来越}, nous allons maintenant découvrir une structure qui permet d'\emph{organiser une argumentation} : présenter deux aspects, deux causes ou deux points de vue d'une même situation. C'est l'équivalent chinois du français « d'une part…, d'autre part… » : \cle{一方面……另一方面……}. Cette structure est très fréquente à l'oral comme à l'écrit, et elle sera précieuse pour vos exposés et vos rédactions.

\courssub{Observation guidée}

Lisons d'abord deux phrases. Observez bien : où se trouve le sujet ? Quels mots reviennent ? Comment les deux parties sont-elles reliées ?

\begin{exemplebox}[Phrases d'observation]
1. \zh{我想去中国，一方面可以学中文，另一方面可以看长城。}\\
\emph{Wǒ xiǎng qù Zhōngguó, yì fāngmiàn kěyǐ xué Zhōngwén, lìng yì fāngmiàn kěyǐ kàn Chángchéng.}\\
« Je veux aller en Chine : d'une part je pourrai apprendre le chinois, d'autre part je pourrai voir la Grande Muraille. »

2. \zh{住在城里一方面方便，另一方面太贵。}\\
\emph{Zhù zài chéng lǐ yì fāngmiàn fāngbiàn, lìng yì fāngmiàn tài guì.}\\
« Vivre en ville, d'une part c'est pratique, d'autre part c'est trop cher. »
\end{exemplebox}

\begin{center}
\begin{tabularx}{\linewidth}{|X|X|X|X|}\hline
\rowcolor{popblueL} Caractères & Pinyin & Nature & Traduction \\ \hline
一方面 & yì fāngmiàn & locution & d'une part \\ \hline
另一方面 & lìng yì fāngmiàn & locution & d'autre part \\ \hline
方面 & fāngmiàn & nom & aspect, côté \\ \hline
另 & lìng & adjectif & autre \\ \hline
方便 & fāngbiàn & adjectif & pratique, commode \\ \hline
贵 & guì & adjectif & cher, coûteux \\ \hline
长城 & Chángchéng & nom propre & la Grande Muraille \\ \hline
\end{tabularx}
\end{center}

\textbf{Ce que l'on remarque :} chaque phrase présente \emph{deux} informations sur \emph{un même sujet}. Le mot \zh{一方面} introduit le premier aspect, \zh{另一方面} introduit le second, et une virgule chinoise \zh{，} sépare les deux volets. Remarquez aussi que \zh{方面} seul est un nom (« aspect »), mais \zh{一方面} et \zh{另一方面} sont des \emph{locutions figées} : on ne peut ni les couper, ni changer leur ordre.

\courssub{Schéma et sens de la structure}

\begin{propriete}[Règle : 一方面……另一方面……]
La structure \cle{一方面……另一方面……} relie deux propositions parallèles qui présentent deux facettes d'un même fait : deux avantages, deux inconvénients, deux causes ou deux points de vue.\\
\textbf{Schéma :} (Sujet) + \zh{一方面} + proposition 1 \zh{，} \zh{另一方面} + proposition 2 \zh{。}\\
\textbf{Sens :} « d'une part…, d'autre part… » ; « à la fois… et… ».\\
\textbf{Ponctuation :} on met une virgule chinoise \zh{，} entre les deux volets.\\
\textbf{Condition :} les deux volets doivent être \emph{équilibrés} et de \emph{même nature} : deux avantages ensemble, ou deux causes ensemble, ou deux inconvénients ensemble — pas un avantage suivi d'une cause sans rapport.
\end{propriete}

\begin{popfigure}
\begin{tikzpicture}[mindmap, grow cyclic, every node/.style={concept, align=center, font=\small}, concept color=popblue!40, level 1/.append style={sibling angle=120, level distance=3.4cm}]
\node[align=center]{Situation\\主题}
child[concept color=popgreen!50] { node[align=center]{一方面\\aspect 1\\avantage / cause A} }
child[concept color=poporange!50] { node[align=center]{另一方面\\aspect 2\\avantage / cause B} }
child[concept color=poppurple!40] { node[align=center]{，\\virgule\\entre les\\deux volets} };
\end{tikzpicture}
\legende{Carte mentale : une même situation éclairée par deux aspects parallèles.}
\end{popfigure}

\courssub{La place du sujet}

La position du sujet dépend d'une question simple : \emph{le sujet est-il le même dans les deux volets ?}

\begin{propriete}[Position du sujet]
\textbf{Cas 1 — sujet commun aux deux propositions :} le sujet se place \emph{une seule fois, avant} \zh{一方面}.\\
\zh{学中文一方面有意思，另一方面很有用。}\\
\emph{Xué Zhōngwén yì fāngmiàn yǒu yìsi, lìng yì fāngmiàn hěn yǒuyòng.}\\
« Apprendre le chinois, d'une part c'est intéressant, d'autre part c'est très utile. »

\textbf{Cas 2 — sujets différents :} chaque proposition a son propre sujet, placé \emph{après} \zh{一方面} / \zh{另一方面}.\\
\zh{一方面老师帮助我，另一方面同学鼓励我。}\\
\emph{Yì fāngmiàn lǎoshī bāngzhù wǒ, lìng yì fāngmiàn tóngxué gǔlì wǒ.}\\
« D'une part le professeur m'aide, d'autre part les camarades m'encouragent. »
\end{propriete}

\begin{exemplebox}[Exemples variés]
1. \zh{在杜阿拉生活，一方面工作机会多，另一方面交通太拥挤。}\\
\emph{Zài Dù'ālā shēnghuó, yì fāngmiàn gōngzuò jīhuì duō, lìng yì fāngmiàn jiāotōng tài yōngjǐ.}\\
« Vivre à Douala, d'une part il y a beaucoup d'offres d'emploi, d'autre part la circulation est trop dense. »

2. \zh{我喜欢踢足球，一方面可以锻炼身体，另一方面可以交朋友。}\\
\emph{Wǒ xǐhuan tī zúqiú, yì fāngmiàn kěyǐ duànliàn shēntǐ, lìng yì fāngmiàn kěyǐ jiāo péngyou.}\\
« J'aime jouer au football : d'une part cela entretient le corps, d'autre part cela permet de se faire des amis. »

3. \zh{一方面我想去中国留学，另一方面我不想离开家人。}\\
\emph{Yì fāngmiàn wǒ xiǎng qù Zhōngguó liúxué, lìng yì fāngmiàn wǒ bù xiǎng líkāi jiārén.}\\
« D'une part je voudrais aller étudier en Chine, d'autre part je ne veux pas quitter ma famille. »

4. \zh{这个办法一方面省钱，另一方面省时间。}\\
\emph{Zhège bànfǎ yì fāngmiàn shěng qián, lìng yì fāngmiàn shěng shíjiān.}\\
« Cette méthode, d'une part fait économiser de l'argent, d'autre part fait gagner du temps. »
\end{exemplebox}

\courssub{Variantes et emploi argumentatif}

\begin{aretenir}[Variantes]
\begin{itemize}
\item Variante possible : on peut répéter \zh{一方面……一方面……} au lieu de \zh{一方面……另一方面……}, mais la forme avec \zh{另} (« autre ») est la plus claire et la plus fréquente à l'écrit.
\item La forme \zh{一方面……再一方面……} existe mais est \emph{rare} ; évitez-la à ce niveau.
\item \zh{另} signifie « autre » : \zh{另一方面} = littéralement « l'autre aspect ». Attention : on dit \zh{另一方面} et jamais *\zh{二方面}.
\end{itemize}
\end{aretenir}

Cette structure est l'outil de base de l'\emph{argumentation} en chinois : exposés, débats, rédactions. Elle sert à présenter :

\begin{itemize}
\item des \textbf{avantages et inconvénients} : \zh{网购一方面便宜，另一方面不能试衣服。} \emph{Wǎnggòu yì fāngmiàn piányi, lìng yì fāngmiàn bù néng shì yīfu.} « Les achats en ligne, d'une part c'est bon marché, d'autre part on ne peut pas essayer les vêtements. »
\item des \textbf{causes multiples} : \zh{他汉语进步很快，一方面因为他努力，另一方面因为他有一个中国朋友。} \emph{Tā Hànyǔ jìnbù hěn kuài, yì fāngmiàn yīnwèi tā nǔlì, lìng yì fāngmiàn yīnwèi tā yǒu yí ge Zhōngguó péngyou.} « Il progresse vite en chinois : d'une part parce qu'il travaille dur, d'autre part parce qu'il a un ami chinois. »
\item des \textbf{points de vue} : \zh{一方面家长希望孩子学医，另一方面孩子自己喜欢艺术。} \emph{Yì fāngmiàn jiāzhǎng xīwàng háizi xué yī, lìng yì fāngmiàn háizi zìjǐ xǐhuan yìshù.} « D'une part les parents espèrent que l'enfant fasse médecine, d'autre part l'enfant, lui, aime l'art. »
\end{itemize}

\begin{center}
\begin{tabularx}{\linewidth}{|X|X|X|}\hline
\rowcolor{popblueL} Structure & Exemple (caractères + pinyin) & Traduction \\ \hline
Sujet commun + 一方面…另一方面… & \zh{这个城市一方面大，另一方面很现代。} \emph{Zhège chéngshì yì fāngmiàn dà, lìng yì fāngmiàn hěn xiàndài.} & Cette ville, d'une part est grande, d'autre part est très moderne. \\ \hline
一方面 + sujet 1…另一方面 + sujet 2… & \zh{一方面爸爸支持我，另一方面妈妈担心我。} \emph{Yì fāngmiàn bàba zhīchí wǒ, lìng yì fāngmiàn māma dānxīn wǒ.} & D'une part papa me soutient, d'autre part maman s'inquiète pour moi. \\ \hline
Avec 因为 (causes) & \zh{一方面因为远，另一方面因为贵，我们没去。} \emph{Yì fāngmiàn yīnwèi yuǎn, lìng yì fāngmiàn yīnwèi guì, wǒmen méi qù.} & D'une part parce que c'était loin, d'autre part parce que c'était cher, nous n'y sommes pas allés. \\ \hline
\end{tabularx}
\end{center}

\courssub{Comparaison avec le français et pièges}

\begin{attention}[Erreurs fréquentes des francophones]
\begin{itemize}
\item \textbf{Ne pas oublier le second volet.} En français, on peut dire « d'une part, c'est pratique » et s'arrêter là. En chinois, \zh{一方面} appelle presque toujours \zh{另一方面} : la structure attend les deux volets. Phrase fautive : *\zh{一方面很方便。}
\item \textbf{Ne pas traduire mot à mot} « d'une part » par \zh{一个方面} : la locution figée est \zh{一方面} (sans classificateur \zh{个}).
\item \textbf{Équilibre des volets.} Ne mélangez pas un avantage et une cause sans rapport : *\zh{一方面很方便，另一方面因为我喜欢} est maladroit. Présentez deux éléments de même nature.
\item \textbf{Virgule obligatoire} \zh{，} entre les deux volets, même si le français varie sa ponctuation (deux-points, point-virgule).
\item \textbf{Ordre figé :} on ne dit jamais *\zh{方面一} ni *\zh{方面另一}.
\item \textbf{Ne pas confondre avec} \zh{越来越} : \zh{越来越} décrit \emph{un changement progressif} (une seule qualité qui évolue), tandis que \zh{一方面……另一方面……} présente \emph{deux aspects simultanés}. Comparez : \zh{房租越来越贵。} « Le loyer est de plus en plus cher » (évolution) et \zh{这个房子一方面大，另一方面便宜。} « Cette maison, d'une part est grande, d'autre part est bon marché » (deux aspects).
\end{itemize}
\end{attention}

En français, « d'une part… d'autre part… » est soutenu ; en chinois, \zh{一方面……另一方面……} s'emploie aussi bien à l'oral qu'à l'écrit, dans la conversation courante comme dans les dissertations. C'est donc une structure à maîtriser absolument pour l'exposé oral.

\begin{methode}[Construire une phrase avec 一方面……另一方面……]
\begin{enumerate}
\item \textbf{Identifier la situation} ou le thème unique dont on veut parler (ex. : vivre à Yaoundé).
\item \textbf{Lister deux aspects de même nature} : deux avantages, deux inconvénients ou deux causes (ex. : climat agréable / loyer élevé).
\item \textbf{Rédiger chaque volet} : sujet commun avant \zh{一方面}, ou sujet propre à chaque volet après chaque \zh{方面}.
\item \textbf{Vérifier l'équilibre} : les deux volets sont-ils parallèles ? La virgule \zh{，} est-elle en place ? Le second volet commence-t-il bien par \zh{另一方面} ?
\end{enumerate}
\end{methode}

\courssub{Texte de lecture}

\begin{textebox}[在城里还是在乡下？]
\zh{毕业以后，很多年轻人要做一个重要的选择：住在城里，还是住在乡下？住在城里一方面方便，另一方面太贵。城里有很多工作机会，交通也方便，可是房租越来越高。住在乡下，一方面空气好、生活安静，另一方面工作少、医院和学校也比较远。我的哥哥选择了城里，他说：“一方面我想挣钱，另一方面我想学新的东西。”我的姐姐不一样，她回到了乡下，因为她觉得，一方面乡下的生活更健康，另一方面她可以照顾父母。其实，每个选择都有两个方面，最重要的是知道自己真正想要什么。}\\
\emph{Bìyè yǐhòu, hěn duō niánqīngrén yào zuò yí ge zhòngyào de xuǎnzé : zhù zài chéng lǐ, háishi zhù zài xiāngxia ? Zhù zài chéng lǐ yì fāngmiàn fāngbiàn, lìng yì fāngmiàn tài guì. Chéng lǐ yǒu hěn duō gōngzuò jīhuì, jiāotōng yě fāngbiàn, kěshì fángzū yuèláiyuè gāo. Zhù zài xiāngxia, yì fāngmiàn kōngqì hǎo, shēnghuó ānjìng, lìng yì fāngmiàn gōngzuò shǎo, yīyuàn hé xuéxiào yě bǐjiào yuǎn. Wǒ de gēge xuǎnzé le chéng lǐ, tā shuō : « Yì fāngmiàn wǒ xiǎng zhèng qián, lìng yì fāngmiàn wǒ xiǎng xué xīn de dōngxi. » Wǒ de jiějie bù yíyàng, tā huí dào le xiāngxia, yīnwèi tā juéde, yì fāngmiàn xiāngxia de shēnghuó gèng jiànkāng, lìng yì fāngmiàn tā kěyǐ zhàogù fùmǔ. Qíshí, měi ge xuǎnzé dōu yǒu liǎng ge fāngmiàn, zuì zhòngyào de shì zhīdào zìjǐ zhēnzhèng xiǎng yào shénme.}\\
« Après les études, beaucoup de jeunes doivent faire un choix important : vivre en ville ou à la campagne ? Vivre en ville, d'une part c'est pratique, d'autre part c'est trop cher. En ville, il y a beaucoup d'emplois, les transports sont commodes, mais les loyers sont de plus en plus élevés. Vivre à la campagne, d'une part l'air est pur et la vie tranquille, d'autre part il y a peu de travail et les hôpitaux et écoles sont assez loin. Mon frère aîné a choisi la ville ; il dit : "D'une part je veux gagner de l'argent, d'autre part je veux apprendre des choses nouvelles." Ma sœur aînée, elle, est différente : elle est retournée à la campagne, car elle pense que d'une part la vie à la campagne est plus saine, d'autre part elle peut s'occuper de ses parents. En fait, chaque choix a deux faces ; le plus important est de savoir ce que l'on veut vraiment. »
\end{textebox}

\begin{center}
\begin{tabularx}{\linewidth}{|X|X|X|X|}\hline
\rowcolor{popblueL} Caractères & Pinyin & Nature & Traduction \\ \hline
毕业 & bìyè & verbe & terminer ses études \\ \hline
选择 & xuǎnzé & verbe / nom & choisir ; choix \\ \hline
乡下 & xiāngxia & nom & campagne \\ \hline
房租 & fángzū & nom & loyer \\ \hline
安静 & ānjìng & adjectif & calme, tranquille \\ \hline
挣钱 & zhèng qián & verbe & gagner de l'argent \\ \hline
照顾 & zhàogù & verbe & s'occuper de \\ \hline
其实 & qíshí & adverbe & en fait \\ \hline
\end{tabularx}
\end{center}

\textbf{Questions de compréhension :}
\begin{enumerate}
\item Quels sont les deux aspects de la vie en ville mentionnés dans le texte ?
\item Pourquoi la sœur aînée est-elle retournée à la campagne ? (Répondez avec \zh{一方面……另一方面……}.)
\item Relevez dans le texte une phrase construite avec \zh{越来越}, puis une phrase construite avec \zh{一方面……另一方面……}.
\end{enumerate}

\courssub{Exercice résolu et entraînement}

\begin{exoresolu}[Transformation de phrases]
Énoncé : combinez les deux phrases suivantes en une seule phrase avec \zh{一方面……另一方面……}.\\
\zh{坐火车旅行很便宜。坐火车旅行可以看风景。}\\
\emph{Zuò huǒchē lǚxíng hěn piányi. Zuò huǒchē lǚxíng kěyǐ kàn fēngjǐng.}
\tcblower
Solution détaillée :\\
Étape 1 : la situation commune est \zh{坐火车旅行} (voyager en train) ; le sujet est commun aux deux propositions, il se place donc une seule fois, avant \zh{一方面}.\\
Étape 2 : les deux aspects sont deux avantages (bon marché / voir le paysage) : ils sont de même nature, l'équilibre est respecté.\\
Étape 3 : on relie avec \zh{一方面……，另一方面……}.\\
Réponse : \zh{坐火车旅行一方面很便宜，另一方面可以看风景。}\\
\emph{Zuò huǒchē lǚxíng yì fāngmiàn hěn piányi, lìng yì fāngmiàn kěyǐ kàn fēngjǐng.}\\
« Voyager en train, d'une part c'est bon marché, d'autre part on peut admirer le paysage. »
\end{exoresolu}

\begin{exercice}[À vous de jouer]
1. Traduisez en chinois : « Apprendre le chinois, d'une part c'est difficile, d'autre part c'est passionnant. »\\
2. Combinez : \zh{雅温得的天气很舒服。雅温得的房租不便宜。} (sujet commun : \zh{雅温得})\\
3. Complétez avec deux sujets différents : \zh{一方面我……，另一方面我的弟弟……。}\\
4. Corrigez la phrase fautive : *\zh{一个方面中文有意思，另一个方面有用。}\\
5. Choisissez la bonne structure (\zh{越来越} ou \zh{一方面……另一方面……}) : « Le chinois est \dots{} difficile, mais d'une part les caractères sont beaux, d'autre part la grammaire est simple. »
\end{exercice}

\begin{corrige}[Exercice « À vous de jouer »]
1. \zh{学中文一方面很难，另一方面很有意思。} \emph{Xué Zhōngwén yì fāngmiàn hěn nán, lìng yì fāngmiàn hěn yǒu yìsi.}\\
2. \zh{雅温得一方面天气很舒服，另一方面房租不便宜。} \emph{Yǎwēndé yì fāngmiàn tiānqì hěn shūfu, lìng yì fāngmiàn fángzū bù piányi.}\\
3. Exemple : \zh{一方面我喜欢足球，另一方面我的弟弟喜欢篮球。} \emph{Yì fāngmiàn wǒ xǐhuan zúqiú, lìng yì fāngmiàn wǒ de dìdi xǐhuan lánqiú.}\\
4. Correction : \zh{中文一方面有意思，另一方面很有用。} — la locution figée est \zh{一方面} / \zh{另一方面}, sans classificateur \zh{个}.\\
5. \zh{中文越来越难，但是一方面汉字很美，另一方面语法很简单。} \emph{Zhōngwén yuèláiyuè nán, dànshì yì fāngmiàn Hànzì hěn měi, lìng yì fāngmiàn yǔfǎ hěn jiǎndān.} — \zh{越来越} pour l'évolution (« de plus en plus difficile »), \zh{一方面……另一方面……} pour les deux aspects.
\end{corrige}

\begin{aretenir}[L'essentiel de la structure]
\begin{itemize}
\item Schéma : (Sujet) + \zh{一方面} + proposition 1 \zh{，} \zh{另一方面} + proposition 2 \zh{。}
\item Sens : « d'une part…, d'autre part… » — deux facettes d'un même fait.
\item Sujet commun : placé avant \zh{一方面} ; sujets différents : un sujet dans chaque volet.
\item Les deux volets sont obligatoires, équilibrés et séparés par \zh{，}.
\item Ne pas confondre avec \zh{越来越} (évolution progressive d'une seule qualité).
\item Structure idéale pour l'exposé et l'argumentation (avantages/inconvénients, causes, points de vue).
\end{itemize}
\end{aretenir}

\coursec{Comparaison avec le français et pièges des francophones}

Dans la section précédente, nous avons construit pas à pas la structure \zh{一方面……，另一方面……} et vérifié qu'elle sert à présenter deux aspects parallèles d'une même situation. Nous savons donc maintenant \emph{former} nos deux structures. Reste une question décisive pour un élève francophone : que se passe-t-il quand le français et le chinois se ressemblent, et que se passe-t-il quand ils s'éloignent ? C'est l'objet de cette section : repérer les \cle{transferts positifs} (ce que le français nous offre gratuitement) et désamorcer les \cle{pièges} (ce que le français nous pousse à écrire de faux).

\courssub{越来越 et « de plus en plus » : un transfert presque parfait}

Observons d'abord deux phrases jumelles :

\zh{天气越来越热。} \emph{Tiānqì yuè lái yuè rè.} « Il fait de plus en plus chaud. »

\zh{他的汉语越来越好。} \emph{Tā de Hànyǔ yuè lái yuè hǎo.} « Son chinois est de mieux en mieux. »

Bonne nouvelle : l'ordre des constituants est \textbf{le même} en français et en chinois. Dans les deux langues, la marque de progression (\zh{越来越} / « de plus en plus ») se place \emph{juste avant} l'adjectif ou le verbe psychologique, et \emph{après} le sujet. Le francophone peut donc transposer mot à mot :

\begin{center}
\begin{tabularx}{\linewidth}{|X|X|X|}
\hline
\rowcolor{popblueL} \textbf{Français} & \textbf{Chinois correct} & \textbf{Erreur fréquente} \\
\hline
Il fait de plus en plus chaud. & 天气越来越热。 & \textcolor{poppink}{越来越天气热。} \\
\hline
Son chinois est de mieux en mieux. & 他的汉语越来越好。 & \textcolor{poppink}{他的汉语越来越很好。} \\
\hline
Yaoundé devient de plus en plus grand. & 雅温得越来越大。 & \textcolor{poppink}{雅温得很越来越大。} \\
\hline
J'aime de moins en moins le bruit. & 我越来越不喜欢吵闹。 & \textcolor{poppink}{我不越来越喜欢吵闹。} \\
\hline
\end{tabularx}
\end{center}

\begin{propriete}[Correspondance de position]
\zh{越来越} occupe exactement la place de « de plus en plus » : \textbf{Sujet + 越来越 + adjectif / verbe psychologique}. Le transfert depuis le français est donc \cle{positif} pour l'ordre des mots : il suffit de calquer la phrase française.
\end{propriete}

\courssub{Le piège du double degré : pourquoi 很 est exclu}

Le transfert positif s'arrête là où le français additionne les marques de degré et où le chinois les refuse. En français familier, on entend parfois « c'est de plus en plus \emph{très} compliqué » ; même à l'écrit, l'élève ressent le besoin de « renforcer » l'adjectif. En chinois, c'est impossible : \zh{越来越} contient déjà l'idée de degré croissant, et il est \textbf{incompatible} avec \zh{很} (très), \zh{太} (trop) et \zh{非常} (extrêmement).

\begin{exemplebox}[Faux et juste]
FAUX : \zh{*天气越来越很热。} (double marque de degré)\\
JUSTE : \zh{天气越来越热。} \emph{Tiānqì yuè lái yuè rè.} « Il fait de plus en plus chaud. »\\[0.4em]
FAUX : \zh{*他越来越太忙。}\\
JUSTE : \zh{他越来越忙。} \emph{Tā yuè lái yuè máng.} « Il est de plus en plus occupé. »\\[0.4em]
FAUX : \zh{*足球越来越非常受欢迎。}\\
JUSTE : \zh{足球越来越受欢迎。} \emph{Zúqiú yuè lái yuè shòu huānyíng.} « Le football est de plus en plus populaire. »
\end{exemplebox}

\begin{attention}[Pièges 1, 2 et 3 : les trois erreurs de placement et de cumul]
\begin{itemize}
\item \textbf{Piège 1 — 很 avant 越来越} : \zh{*他很越来越忙。} On ne peut pas « entourer » l'adjectif de deux marques de degré. Choisissez : soit \zh{他很忙} (il est très occupé), soit \zh{他越来越忙} (il est de plus en plus occupé), jamais les deux.
\item \textbf{Piège 2 — 很 après 越来越} : \zh{*越来越很忙} est le calque direct du français « de plus en plus très ». Après \zh{越来越}, l'adjectif reste \emph{nu}.
\item \textbf{Piège 3 — 越来越 avant le sujet} : \zh{*越来越他忙。} En français, « de plus en plus » peut commencer une phrase (« De plus en plus, il est fatigué ») ; en chinois, \zh{越来越} ne se détache jamais de l'adjectif et reste \emph{après} le sujet.
\end{itemize}
\end{attention}

\courssub{« De moins en moins » : 越来越不 + adjectif}

Le français dit « de moins en moins » en changeant « plus » en « moins ». Le chinois procède autrement : il garde \zh{越来越} intact et place la négation \zh{不} \textbf{entre} \zh{越来越} et l'adjectif ou le verbe.

\begin{propriete}[Schéma de la forme négative]
\textbf{Sujet + 越来越 + 不 + adjectif / verbe} = « de moins en moins… ».\\
La négation ne précède jamais \zh{越来越} : \zh{*不越来越} n'existe pas.
\end{propriete}

\zh{我越来越不喜欢大城市的吵闹。} \emph{Wǒ yuè lái yuè bù xǐhuan dà chéngshì de chǎonào.} « J'aime de moins en moins le bruit des grandes villes. »

\zh{他的病越来越不好，家人很担心。} \emph{Tā de bìng yuè lái yuè bù hǎo, jiārén hěn dānxīn.} « Sa santé va de moins en moins bien, sa famille est très inquiète. »

Remarquez que \zh{越来越不好} signifie « de moins en moins bon », c'est-à-dire « de pire en pire » : la négation porte sur l'adjectif, pas sur la progression.

\courssub{一方面……另一方面…… et « d'une part… d'autre part »}

Ici encore, la correspondance avec le français est presque directe : \zh{一方面} = « d'une part », \zh{另一方面} = « d'autre part ». Les deux volets sont séparés par une virgule chinoise ，et présentent deux aspects \emph{du même sujet}, sur le même plan logique.

\zh{在杜阿拉生活，一方面工作机会多，另一方面房租很贵。} \emph{Zài Dù'ālā shēnghuó, yī fāngmiàn gōngzuò jīhuì duō, lìng yī fāngmiàn fángzū hěn guì.} « Vivre à Douala : d'une part, les offres d'emploi sont nombreuses ; d'autre part, les loyers sont chers. »

\zh{学中文一方面很难，另一方面很有意思。} \emph{Xué Zhōngwén yī fāngmiàn hěn nán, lìng yī fāngmiàn hěn yǒu yìsi.} « Apprendre le chinois : d'une part c'est difficile, d'autre part c'est très intéressant. »

\begin{attention}[但是 n'est pas 另一方面]
Les francophones remplacent souvent le second volet par \zh{但是} (mais). Or les deux mots n'ont pas la même fonction :
\begin{itemize}
\item \zh{但是} marque une \textbf{opposition} : le second élément contredit ou nuance le premier. \zh{中文很难，但是很有意思。} « Le chinois est difficile, \emph{mais} intéressant. »
\item \zh{另一方面} marque un \textbf{parallélisme} : on ajoute un second aspect, sans le mettre en contradiction avec le premier. \zh{一方面很难，另一方面很有意思。} « D'une part difficile, d'autre part intéressant » — les deux constats coexistent.
\end{itemize}
Si vous avez ouvert avec \zh{一方面}, vous devez fermer avec \zh{另一方面} ; \zh{*一方面……，但是……} est une structure bancale.
\end{attention}

\begin{attention}[Piège 5 : les deux volets doivent être équilibrés]
\zh{一方面……另一方面……} présente deux aspects \emph{de même nature} : deux avantages, deux inconvénients, ou un avantage et un inconvénient. On ne peut pas opposer un avantage à une cause :\\
FAUX : \zh{*一方面足球很受欢迎，另一方面因为喀麦隆人喜欢运动。} (un constat contre une justification)\\
JUSTE : \zh{一方面足球很受欢迎，另一方面篮球也越来越流行。} \emph{Yī fāngmiàn zúqiú hěn shòu huānyíng, lìng yī fāngmiàn lánqiú yě yuè lái yuè liúxíng.} « D'une part le football est très populaire, d'autre part le basket-ball devient de plus en plus à la mode. »
\end{attention}

\begin{popfigure}
\begin{tikzpicture}[
  box/.style={draw=popblue, fill=popblueL, rounded corners=2pt, align=center, font=\small, inner sep=4pt},
  bad/.style={draw=poppink, fill=poppinkL, rounded corners=2pt, align=center, font=\small, inner sep=4pt},
  arr/.style={-{Stealth[length=2mm]}, thick, popdark}
]
\node[box, align=center] (fr) at (0,0){« de plus en plus »\\ + adjectif};
\node[box, right=3.2cm of fr, align=center] (zh){Sujet + 越来越\\ + adjectif nu};
\node[bad, below=1.1cm of zh, align=center] (err){interdit : 很, 太, 非常\\ pas avant le sujet};
\node[box, below=1.1cm of fr, align=center] (neg){« de moins en moins »\\ = 越来越 + 不 + adj.};
\draw[arr] (fr) -- node[above, font=\small]{transfert positif} (zh);
\draw[arr, poppink] (zh) -- (err);
\draw[arr] (fr) -- (neg);
\draw[arr, popgreen] (neg.east) -- (err.west |- neg.east) node[midway, above, font=\small, popgreen]{不 après 越来越};
\end{tikzpicture}
\legende{Schéma de transformation : du français « de plus en plus » au chinois 越来越, et formation du négatif « de moins en moins ».}
\end{popfigure}

\begin{methode}[S'auto-corriger en trois vérifications]
Avant de rendre une phrase avec \zh{越来越}, relisez-la et posez-vous trois questions :
\begin{enumerate}
\item \textbf{Après} \zh{越来越}, l'adjectif est-il nu ? Si vous voyez \zh{很}, \zh{太} ou \zh{非常}, supprimez-le.
\item \zh{越来越} est-il bien \emph{après} le sujet, jamais en tête de phrase ?
\item Pour une phrase négative, le \zh{不} est-il \emph{entre} \zh{越来越} et l'adjectif (\zh{越来越不喜欢}), jamais devant \zh{越来越} ?
\end{enumerate}
Et pour \zh{一方面……另一方面……} : le second volet commence-t-il bien par \zh{另一方面} (et non \zh{但是}), et les deux volets sont-ils de même nature ?
\end{methode}

\begin{exoresolu}[Corriger les phrases fautives]
Énoncé : les phrases suivantes ont été écrites par des élèves francophones. Repérez l'erreur et réécrivez chaque phrase correctement.\\
1. \zh{*他越来越很高。} \quad 2. \zh{*越来越我喜欢这个城市。} \quad 3. \zh{*我不越来越喜欢吃辣。} \quad 4. \zh{*一方面交通很方便，但是房租太贵。}
\tcblower
Solution détaillée :
\begin{enumerate}
\item Double marque de degré : \zh{很} est interdit après \zh{越来越}. Correction : \zh{他越来越高。} \emph{Tā yuè lái yuè gāo.} « Il devient de plus en plus grand. »
\item \zh{越来越} ne peut pas ouvrir la phrase : il doit suivre le sujet. Correction : \zh{我越来越喜欢这个城市。} \emph{Wǒ yuè lái yuè xǐhuan zhège chéngshì.} « J'aime de plus en plus cette ville. »
\item La négation mal placée : « de moins en moins » se dit \zh{越来越不}. Correction : \zh{我越来越不喜欢吃辣。} \emph{Wǒ yuè lái yuè bù xǐhuan chī là.} « J'aime de moins en moins manger pimenté. »
\item Mélange de structures : après \zh{一方面}, il faut \zh{另一方面}. Correction : \zh{一方面交通很方便，另一方面房租太贵。} \emph{Yī fāngmiàn jiāotōng hěn fāngbiàn, lìng yī fāngmiàn fángzū tài guì.} « D'une part les transports sont pratiques, d'autre part le loyer est trop cher. » (On pouvait aussi supprimer \zh{一方面} et garder \zh{但是}, mais pas mélanger les deux.)
\end{enumerate}
\end{exoresolu}

\begin{savaistu}[L'abréviation journalistique 另]
À l'écrit, surtout dans la presse et les textes administratifs chinois, \zh{另一方面} s'abrège parfois en un seul caractère : \zh{另} \emph{(lìng)}, comme dans \zh{一方面……，另……}. C'est une économie de moyens typique de la langue écrite chinoise, qui aime la concision. Au lycée, contentez-vous de \emph{reconnaître} cette forme si vous la rencontrez dans un article ; dans vos propres productions, utilisez toujours la forme complète \zh{另一方面}, seule correcte à l'oral et attendue à l'examen.
\end{savaistu}

\begin{aretenir}[L'essentiel de la section]
\begin{itemize}
\item \zh{越来越} = « de plus en plus » : même place qu'en français (après le sujet, avant l'adjectif) — transfert positif.
\item Après \zh{越来越}, adjectif nu : jamais \zh{很}, \zh{太}, \zh{非常}.
\item « De moins en moins » = \zh{越来越不} + adjectif ; \zh{*不越来越} n'existe pas.
\item \zh{一方面……另一方面……} = « d'une part… d'autre part » : parallélisme, volets équilibrés ; \zh{但是} exprime l'opposition et ne ferme jamais un volet ouvert par \zh{一方面}.
\end{itemize}
\end{aretenir}

\coursec{Exercices d'application et de transformation}

Après avoir analysé les différences fondamentales entre le français et le chinois pour exprimer la progression (\textbf{越来越}) et l'opposition coordonnée (\textbf{一方面……另一方面……}), nous passons maintenant à la mise en œuvre active. Cette section propose un parcours progressif : de la manipulation guidée à la production autonome, en passant par la traduction, la correction d'erreurs types et l'analyse d'un texte authentique. Chaque exercice est conçu pour ancrer les automatismes syntaxiques et prévenir les pièges identifiés précédemment.

\courssub{Exercice 1 : Complétion progressive — \zh{越来越} ou \zh{越来越不}}

\begin{methode}[Comment choisir entre \zh{越来越} et \zh{越来越不}]
\begin{enumerate}
\item Identifier l'adjectif ou le verbe d'état qui suit la structure.
\item Déterminer le sens de l'évolution : \textbf{augmentation} de la qualité (positif ou négatif) $\rightarrow$ \zh{越来越} + adjectif ; \textbf{diminution} ou disparition de la qualité $\rightarrow$ \zh{越来越不} + adjectif.
\item Vérifier la cohérence avec le marqueur aspectuel \zh{了} (souvent présent en fin de phrase pour marquer le changement d'état).
\item Relire la phrase à haute voix pour contrôler la prosodie (ton 4 sur \zh{不} $\rightarrow$ ton 2 devant un 4\up{e} ton).
\end{enumerate}
\end{methode}

\begin{exercice}[Complétez les phrases avec \zh{越来越} ou \zh{越来越不} selon le sens logique.]
\begin{enumerate}
\item 现在的手机 \_\_\_\_\_\_\_\_ 贵了，学生买不起。 \hfill (prix en hausse)
\item 这道题 \_\_\_\_\_\_\_\_ 难，我完全看不懂。 \hfill (difficulté croissante)
\item 妹妹长大了，\_\_\_\_\_\_\_\_ 爱哭了，变懂事了。 \hfill (diminution d'un comportement)
\item 雅温得的空气 \_\_\_\_\_\_\_\_ 干净，因为植树造林做得好。 \hfill (amélioration)
\item 考试临近，同学们 \_\_\_\_\_\_\_\_ 紧张，都在复习。 \hfill (augmentation du stress)
\item 这件衣服洗多了，\_\_\_\_\_\_\_\_ 红了，颜色褪了。 \hfill (perte de couleur)
\item 学中文 \_\_\_\_\_\_\_\_ 有意思，因为能认识更多汉字。 \hfill (intérêt grandissant)
\item 最近天气 \_\_\_\_\_\_\_\_ 冷，大家都穿上了外套。 \hfill (baisse de température)
\end{enumerate}
\end{exercice}

\begin{corrige}[Exercice 1]
\begin{enumerate}
\item \zh{越来越} — \zh{现在的手机越来越贵了，学生买不起。} (Xiànzài de shǒujī yuè lái yuè guì le, xuésheng mǎi bù qǐ.) — Les téléphones sont de plus en plus chers…
\item \zh{越来越} — \zh{这道题越来越难，我完全看不懂。} (Zhè dào tí yuè lái yuè nán, wǒ wánquán kàn bù dǒng.) — Cet exercice devient de plus en plus difficile…
\item \zh{越来越不} — \zh{妹妹长大了，越来越不爱哭了，变懂事了。} (Mèimei zhǎngdà le, yuè lái yuè bù ài kū le, biàn dǒngshì le.) — La petite sœur a grandi, elle pleure de moins en moins…
\item \zh{越来越} — \zh{雅温得的空气越来越干净，因为植树造林做得好。} (Yǎwēndé de kōngqì yuè lái yuè gānjìng, yīnwèi zhíshù zàolín zuò de hǎo.) — L'air de Yaoundé est de plus en plus pur…
\item \zh{越来越} — \zh{考试临近，同学们越来越紧张，都在复习。} (Kǎoshì línjìn, tóngxuémen yuè lái yuè jǐnzhāng, dōu zài fùxí.) — L'examen approche, les camarades sont de plus en plus nerveux…
\item \zh{越来越不} — \zh{这件衣服洗多了，越来越不红了，颜色褪了。} (Zhè jiàn yīfu xǐ duō le, yuè lái yuè bù hóng le, yánsè tuì le.) — Ce vêtement lavé trop souvent est de moins en moins rouge…
\item \zh{越来越} — \zh{学中文越来越有意思，因为能认识更多汉字。} (Xué Zhōngwén yuè lái yuè yǒu yìsi, yīnwèi néng rènshi gèng duō hànzì.) — Apprendre le chinois devient de plus en plus intéressant…
\item \zh{越来越} — \zh{最近天气越来越冷，大家都穿上了外套。} (Zuìjìn tiānqì yuè lái yuè lěng, dàjiā dōu chuān shàng le wàitào.) — Ces derniers temps, il fait de plus en plus froid…
\end{enumerate}
\end{corrige}

\begin{attention}[Piège \zh{不} + adjectif vs \zh{没} + verbe]
Ne confondez pas \zh{越来越不} + \textbf{adjectif/verbe d'état} (ex. \zh{不喜欢}、\zh{不累}) et la négation d'action passée avec \zh{没}. \textit{越来越} ne s'emploie qu'avec des adjectifs ou verbes d'état, jamais avec des verbes d'action purs (\zh{*越来越去北京} incorrect). Pour exprimer « aller de plus en plus souvent », utilisez \zh{去北京的次数越来越多} (le nombre de fois où je vais à Pékin augmente).
\end{attention}

\courssub{Exercice 2 : Thème — Traduction du français vers le chinois}

\begin{exercice}[Traduisez les phrases suivantes en chinois (caractères + pinyin). Utilisez \zh{越来越} ou \zh{越来越不}.]
\begin{enumerate}
\item Le prix du riz à Douala devient de plus en plus élevé.
\item Depuis le printemps, il fait de plus en plus chaud à Buea.
\item Mes résultats en chinois sont de moins en moins mauvais (de plus en plus bons).
\item Ce film est de moins en moins intéressant, je veux partir.
\item Avec l'âge, mon grand-père devient de plus en plus sage.
\end{enumerate}
\end{exercice}

\begin{corrige}[Exercice 2]
\begin{enumerate}
\item \zh{杜阿拉的大米价格越来越高了。} (Dù'ālā de dàmǐ jiàgé yuè lái yuè gāo le.)
\item \zh{春天以来，布埃亚的天气越来越热了。} (Chūntiān yǐlái, Bù'āiyà de tiānqì yuè lái yuè rè le.)
\item \zh{我的中文成绩越来越好了。} (Wǒ de Zhōngwén chéngjī yuè lái yuè hǎo le.) \textit{Note :} \zh{越来越不坏} est grammatical mais \zh{越来越好} est l'usage naturel pour exprimer une amélioration.
\item \zh{这部电影越来越没意思了，我想走。} (Zhè bù diànyǐng yuè lái yuè méi yìsi le, wǒ xiǎng zǒu.) \textit{Note :} \zh{没意思} = « sans intérêt », forme figée (verbe + objet lexicalisé) ; \zh{不有意思} ne s'emploie pas.
\item \zh{随着年龄增长，爷爷越来越有智慧了。} (Suízhe niánlíng zēngzhǎng, yéye yuè lái yuè yǒu zhìhuì le.) \textit{Structure :} \zh{随着……} (au fur et à mesure que…) + sujet + \zh{越来越} + adjectif + \zh{了}.
\end{enumerate}
\end{corrige}

\begin{exemplebox}[Vocabulaire clé de l'exercice 2]
\begin{center}
\begin{tabularx}{\linewidth}{|X|X|X|X|}
\hline
\rowcolor{popblueL} Caractères & Pinyin & Nature & Traduction \\
\hline
大米 & dàmǐ & n. & riz \\
价格 & jiàgé & n. & prix \\
布埃亚 & Bù'āiyà & n.p. & Buea \\
成绩 & chéngjī & n. & résultats, notes \\
好 & hǎo & adj. & bon \\
电影 & diànyǐng & n. & film \\
没意思 & méi yìsi & adj. & sans intérêt, ennuyeux \\
随着 & suízhe & prep. & au fur et à mesure que \\
年龄 & niánlíng & n. & âge \\
增长 & zēngzhǎng & v. & augmenter, croître \\
爷爷 & yéye & n. & grand-père paternel \\
智慧 & zhìhuì & n. & sagesse, intelligence \\
\hline
\end{tabularx}
\end{center}
\end{exemplebox}

\courssub{Exercice 3 : Transformation — \zh{一方面……另一方面……}}

\begin{methode}[Transformer deux phrases en une structure \zh{一方面……另一方面……}]
\begin{enumerate}
\item \textbf{Repérer le sujet commun} aux deux phrases (souvent le même nom ou pronom).
\item \textbf{Vérifier la nature des prédicats} : les deux aspects doivent être de même catégorie (deux avantages, deux inconvénients, deux caractéristiques, un pour / un contre).
\item \textbf{Construire la phrase unique} :
      \begin{itemize}
      \item Sujet + \zh{一方面} + prédicat 1 + \zh{，} + \zh{另一方面} + prédicat 2 + \zh{。}
      \item Si les sujets diffèrent mais le thème est commun : \zh{一方面} + Sujet A + prédicat 1 + \zh{，} + \zh{另一方面} + Sujet B + prédicat 2.
      \end{itemize}
\item \textbf{Ponctuation} : virgule chinoise \zh{，} entre les deux propositions ; point final \zh{。}.
\item \textbf{Relire} pour s'assurer que la logique « d'une part… d'autre part » est respectée (pas de contradiction interne non voulue).
\end{enumerate}
\end{methode}

\begin{popfigure}
\begin{tikzpicture}[
  node distance=1.2cm and 1.5cm,
  box/.style={draw=popblue, thick, rounded corners, fill=popblueL, text width=3.2cm, align=center, font=\small},
  arrow/.style={-Stealth, thick, popdark},
  labelbox/.style={font=\footnotesize, text=poppurple, align=center}
]
\node[box, align=center] (A){Phrase A :\\ \zh{学中文有意思。}\\ Xué Zhōngwén yǒu yìsi.};
\node[box, right=of A, align=center] (B){Phrase B :\\ \zh{学中文有用。}\\ Xué Zhōngwén yǒu yòng.};
\node[box, below=of A, align=center] (C){Étape 1 :\\ Sujet commun = \zh{学中文}\\ (Apprendre le chinois)};
\node[box, below=of B, align=center] (D){Étape 2 :\\ Même nature ?\\ Oui : deux qualités positives};
\node[box, below=of C, align=center] (E){Étape 3 :\\ Insertion\\ \zh{一方面……另一方面……}};
\node[box, below=of E, text width=5cm] (F) {Résultat :\\ \zh{学中文一方面有意思，另一方面有用。}\\ Xué Zhōngwén yì fāngmiàn yǒu yìsi, lìng fāngmiàn yǒu yòng.};

\draw[arrow] (A) -- (C);
\draw[arrow] (B) -- (D);
\draw[arrow] (C) -- (E);
\draw[arrow] (D) -- (E);
\draw[arrow] (E) -- (F);
\end{tikzpicture}
\legende{Organigramme de transformation : deux phrases simples $\rightarrow$ structure \zh{一方面……另一方面……}}
\end{popfigure}

\begin{exercice}[Transformez chaque paire de phrases en une seule phrase avec \zh{一方面……另一方面……}.]
\begin{enumerate}
\item \zh{我们学校很大。} / \zh{我们学校很美。}
\item \zh{住在雅温得交通方便。} / \zh{住在雅温得房租贵。}
\item \zh{做作业累。} / \zh{做作业有用。}
\item \zh{中国菜好吃。} / \zh{中国菜健康。}
\item \zh{学法语难。} / \zh{学法语重要。}
\end{enumerate}
\end{exercice}

\begin{corrige}[Exercice 3]
\begin{enumerate}
\item \zh{我们学校一方面很大，另一方面很美。} (Wǒmen xuéxiào yì fāngmiàn hěn dà, lìng fāngmiàn hěn měi.) — Notre école est à la fois grande et belle.
\item \zh{住在雅温得一方面交通方便，另一方面房租贵。} (Zhù zài Yǎwēndé yì fāngmiàn jiāotōng fāngbiàn, lìng fāngmiàn fángzū guì.) — Habiter à Yaoundé : d'une part le transport est pratique, d'autre part le loyer est cher.
\item \zh{做作业一方面累，另一方面有用。} (Zuò zuòyè yì fāngmiàn lèi, lìng fāngmiàn yǒu yòng.) — Faire les devoirs est d'un côté fatigant, de l'autre utile.
\item \zh{中国菜一方面好吃，另一方面健康。} (Zhōngguó cài yì fāngmiàn hǎochī, lìng fāngmiàn jiànkāng.) — La cuisine chinoise est d'une part délicieuse, d'autre part saine.
\item \zh{学法语一方面难，另一方面重要。} (Xué Fǎyǔ yì fāngmiàn nán, lìng fāngmiàn zhòngyào.) — Apprendre le français est d'une part difficile, d'autre part important.
\end{enumerate}
\end{corrige}

\begin{attention}[Erreurs fréquentes des francophones]
\begin{itemize}
\item \textbf{Oubli de la virgule} : \zh{*一方面……另一方面……} sans \zh{，} $\rightarrow$ incorrect. La virgule marque la pause prosodique obligatoire.
\item \textbf{Inversion \zh{一} / \zh{另}} : \zh{*另方面……一方面……} n'existe pas. L'ordre est figé : \zh{一方面} (premier aspect) puis \zh{另一方面} (deuxième aspect).
\item \textbf{Sujets différents sans thème commun} : \zh{*一方面我喜欢足球，另一方面他喜欢篮球} (incohérent). La structure exige un \textbf{thème unique} (même sujet ou même sujet de discussion).
\item \textbf{Confusion avec \zh{一边……一边……}} : \zh{一边……一边……} = actions simultanées (\zh{一边吃饭一边看电视}) ; \zh{一方面……另一方面……} = énumération d'arguments ou de caractéristiques statiques.
\end{itemize}
\end{attention}

\courssub{Exercice 4 : Correction d'erreurs — « Chasse aux fautes francophones »}

\begin{exercice}[Chaque phrase ci-dessous contient une ou plusieurs erreurs typiques. Corrigez-les et expliquez l'erreur.]
\begin{enumerate}
\item \zh{雅温得的天气越来越不热了。} (Contexte : on est en saison sèche, il fait de plus en plus chaud.)
\item \zh{学习汉字一方面难，另方面有意思。}
\item \zh{这本书越来越有意思，另一方面越来越贵。}
\item \zh{我的中文水平越来越不高，因为我天天复习。}
\item \zh{一方面爸爸是老师，另一方面妈妈是医生。} (Sujet : mes parents)
\end{enumerate}
\end{exercice}

\begin{corrige}[Exercice 4]
\begin{enumerate}
\item \textbf{Correction :} \zh{雅温得的天气越来越热了。} \\
      \textbf{Erreur :} \zh{不热} contredit le contexte (saison sèche = chaleur croissante). \zh{越来越} + adjectif négatif (\zh{不热}) exprime une diminution de la chaleur, or la réalité est une augmentation. \textit{Règle :} adapter l'adjectif à la réalité factuelle.
\item \textbf{Correction :} \zh{学习汉字一方面难，另一方面有意思。} \\
      \textbf{Erreur :} \zh{另方面} $\rightarrow$ forme inexistante. Le terme fixé est \zh{另一方面} (l'autre face/côté).
\item \textbf{Correction :} \zh{这本书越来越有意思，但也越来越贵。} \textit{ou} \zh{这本书一方面有意思，另一方面贵。} \\
      \textbf{Erreur :} Mélange maladroit des deux structures. \zh{越来越} décrit une évolution temporelle (diachronie) ; \zh{一方面……另一方面……} énumère des aspects synchrones (états actuels). On ne peut pas dire « \textit{d'une part} de plus en plus intéressant, \textit{d'autre part} de plus en plus cher » sans restructurer. Solutions : conjonction \zh{但} / \zh{可是} pour l'évolution, ou prédicats statiques pour \zh{一方面……}.
\item \textbf{Correction :} \zh{我的中文水平越来越高了，因为我天天复习。} \\
      \textbf{Erreur :} Incohérence logique : \zh{越来越不高} (de plus en plus bas) contredit la cause \zh{天天复习} (révision quotidienne $\rightarrow$ progrès). \zh{不} est en trop.
\item \textbf{Correction :} \zh{我爸爸是老师，妈妈是医生。} (Structure \zh{一方面……} inadaptée ici). \\
      \textbf{Erreur :} Sujet de la phrase \zh{我父母} (mes parents) mais prédicats portent sur \zh{爸爸} et \zh{妈妈} individuellement. La structure \zh{一方面……另一方面……} s'applique à \textbf{un même sujet} évalué sous deux angles. Pour décrire deux personnes, utiliser deux phrases ou une phrase coordonnée simple.
\end{enumerate}
\end{corrige}

\courssub{Exercice 5 : Version — Traduction vers le français}

\begin{textebox}[Paragraphe modèle — \textit{Notre école change}]
我们学校这几年变化很大。校园越来越大了，新建了教学楼和图书馆。一方面学生越来越多，今年新生有五百多人；另一方面老师也越来越多，外教有五位了。食堂饭菜一方面越来越好吃，另一方面价格还是不变。大家都说：我们学校越来越像大学了！
\par\medskip
\emph{Wǒmen xuéxiào zhè jǐ nián biànhuà hěn dà. Xiàoyuán yuè lái yuè dà le, xīn jiàn le jiàoxué lóu hé túshūguǎn. Yì fāngmiàn xuésheng yuè lái yuè duō, jīnnián xīnshēng yǒu wǔbǎi duō rén; lìng fāngmiàn lǎoshī yě yuè lái yuè duō, wàijiào yǒu wǔ wèi le. Shítáng fàncài yì fāngmiàn yuè lái yuè hǎochī, lìng fāngmiàn jiàgé háishì bù biàn. Dàjiā dōu shuō: wǒmen xuéxiào yuè lái yuè xiàng dàxué le!}
\par\medskip
Notre école a beaucoup changé ces dernières années. Le campus devient de plus en plus grand, on a construit de nouveaux bâtiments d'enseignement et une bibliothèque. D'une part, les élèves sont de plus en plus nombreux — cette année, les nouveaux inscrits sont plus de cinq cents ; d'autre part, les professeurs sont aussi de plus en plus nombreux, il y a maintenant cinq enseignants étrangers. À la cantine, les plats sont d'une part de plus en plus bons, d'autre part le prix n'a pas changé. Tout le monde dit : notre école ressemble de plus en plus à une université !
\end{textebox}

\begin{exercice}[Questions de compréhension et de traduction]
\begin{enumerate}
\item Traduisez le paragraphe ci-dessus en français naturel (version littéraire).
\item Identifiez et soulignez dans le texte chinois : a) \textbf{toutes} les occurrences de \zh{越来越} + adjectif/verbe d'état ; b) les deux structures \zh{一方面……另一方面……}.
\item Pourquoi l'auteur utilise-t-il \zh{还是不变} (hái shi bù biàn) à la fin ? Quel contraste cela crée-t-il ?
\item Réécrivez la dernière phrase (\zh{我们学校越来越像大学了}) en remplaçant \zh{像} par un synonyme structurel (indice : \zh{越来越} + adjectif).
\end{enumerate}
\end{exercice}

\begin{corrige}[Exercice 5]
\begin{enumerate}
\item \textit{Traduction proposée :} « Notre établissement a connu d'importantes transformations ces dernières années. Le campus s'agrandit sans cesse, avec la construction de nouveaux bâtiments pédagogiques et d'une bibliothèque. D'un côté, les effectifs élèves ne cessent d'augmenter — cette année, les nouveaux inscrits dépassent cinq cents ; de l'autre, le corps professoral s'étoffe également, comptant désormais cinq enseignants étrangers. À la cantine, les repas gagnent en qualité tout en maintenant leurs tarifs inchangés. Tous s'accordent à dire : notre lycée ressemble de plus en plus à une université ! »
\item a) \zh{越来越大} (l. 2), \zh{越来越多} (l. 3, deux fois : élèves, professeurs), \zh{越来越好吃} (l. 4), \zh{越来越像} (l. 5). \\
      b) Structure 1 (l. 3) : \zh{一方面学生越来越多……另一方面老师也越来越多}. Structure 2 (l. 4) : \zh{一方面越来越好吃，另一方面价格还是不变}.
\item \zh{还是不变} (« reste inchangé ») crée un \textbf{contraste pragmatique} : qualité en hausse (\zh{越来越好吃}) mais prix stable. C'est un argument de satisfaction (bon rapport qualité-prix). Structure : \zh{还是} (toujours/encore) + \zh{不变} (inchangé).
\item \zh{我们学校越来越现代化了} (Notre école devient de plus en plus modernisée) ou \zh{我们学校越来越正规了} (de plus en plus officielle/standardisée). \zh{现代化} et \zh{正规} fonctionnent comme adjectifs d'état après \zh{越来越}.
\end{enumerate}
\end{corrige}

\courssub{Exercice 6 : Production libre guidée — Ma ville / Mon école}

\begin{methode}[Réussir sa production écrite avec les deux structures]
\begin{enumerate}
\item \textbf{Choisir un sujet familier} : votre ville (Yaoundé, Douala, Buea, Ngaoundéré…), votre lycée, votre quartier.
\item \textbf{Lister 3--4 changements observés} (taille, population, transports, commerces, propreté, sécurité…).
\item \textbf{Sélectionner deux aspects pour \zh{一方面……另一方面……}} : soit deux évolutions parallèles (ex. \zh{人口多 / 房价高}), soit un pour / un contre (ex. \zh{方便 / 拥挤}).
\item \textbf{Rédiger 3 phrases minimum} :
      \begin{itemize}
      \item Phrase 1 : bilan global avec \zh{越来越} + adjectif + \zh{了}.
      \item Phrase 2 : analyse nuancée avec \zh{一方面……另一方面……}.
      \item Phrase 3 (optionnelle) : avis personnel ou conséquence avec \zh{越来越} + verbe d'état.
      \end{itemize}
\item \textbf{Relire} : vérifier l'ordre des mots, la ponctuation chinoise, l'accord des tons sur \zh{不} (4\up{e} ton $\rightarrow$ 2\up{e} ton devant un 4\up{e} ton).
\end{enumerate}
\end{methode}

\begin{exercice}[Production écrite — 3 phrases minimum]
Rédigez un court paragraphe (3--4 phrases) sur \textbf{votre ville} ou \textbf{votre école} en utilisant \textbf{au moins une fois} \zh{越来越} et \textbf{au moins une fois} \zh{一方面……另一方面……}.
\begin{itemize}
\item Écrivez en caractères chinois (simplifiés).
\item Ajoutez le pinyin avec tons sous chaque phrase.
\item Traduisez en français.
\item Vocabulaire suggéré : \zh{交通} (transport), \zh{环境} (environnement), \zh{商场} (centre commercial), \zh{公园} (parc), \zh{堵车} (embouteillage), \zh{绿化} (végétalisation), \zh{安静} (calme), \zh{热闹} (animé).
\end{itemize}
\end{exercice}

\begin{exoresolu}[Exemple de production — Ville de Douala]
\textbf{Énoncé :} Écrire 3 phrases sur Douala avec les deux structures.

\textbf{Production de l'élève (modèle) :}
\begin{enumerate}
\item \zh{杜阿拉这几年变化很大，城市越来越现代化了。} \\
      Dù'ālā zhè jǐ nián biànhuà hěn dà, chéngshì yuè lái yuè xiàndàihuà le. \\
      Douala a beaucoup changé ces années, la ville devient de plus en plus modernisée.
\item \zh{一方面交通越来越方便，新建了立交桥和快速路；另一方面高峰期越来越堵车，大家上班都很累。} \\
      Yì fāngmiàn jiāotōng yuè lái yuè fāngbiàn, xīn jiàn le lìjiāoqiáo hé kuàisùlù; lìng fāngmiàn gāofēngqī yuè lái yuè dǔchē, dàjiā shàngbān dōu hěn lèi. \\
      D'une part le transport devient de plus en plus pratique, on a construit des échangeurs et des voies rapides ; d'autre part aux heures de pointe c'est de plus en plus bouché, tout le monde est fatigué pour aller au travail.
\item \zh{不过海边公园越来越漂亮，周末越来越多人去散步。} \\
      Bùguò hǎibiān gōngyuán yuè lái yuè piàoliang, zhōumò yuè lái yuè duō rén qù sànbù. \\
      Mais le parc en bord de mer devient de plus en plus beau, le week-end de plus en plus de gens vont s'y promener.
\end{enumerate}

\textbf{Analyse didactique :}
\begin{itemize}
\item Phrase 1 : \zh{越来越现代化} (adjectif composé) + \zh{了} (changement d'état) — correct.
\item Phrase 2 : double structure \zh{一方面……另一方面……} avec sujets différents mais thème commun (transport) — correct. \zh{高峰期} (heure de pointe) comme sujet de la 2\up{e} proposition. \zh{越来越堵车} (verbe d'état \zh{堵车} = être embouteillé) — correct.
\item Phrase 3 : \zh{越来越漂亮} + \zh{越来越多人} (structure \zh{越来越多} + nom) — correct. \zh{不过} (mais) introduit une nuance positive.
\end{itemize}
\end{exoresolu}

\begin{aretenir}[Check-list finale avant de rendre votre production]
\begin{itemize}
\item \zh{越来越} est suivi d'un \textbf{adjectif} ou \textbf{verbe d'état} (pas de verbe d'action).
\item \zh{越来越不} + adjectif pour exprimer une diminution.
\item \zh{了} en fin de phrase si changement d'état par rapport au passé.
\item \zh{一方面……另一方面……} : virgule \zh{，} obligatoire, ordre figé, thème unique.
\item Caractères à soigner : \zh{越} (12 traits, clé \zh{走}), \zh{面} (9 traits, clé \zh{面}), \zh{方} (4 traits, clé \zh{方}).
\item Ponctuation chinoise pleine chasse : \zh{，} \zh{。} \zh{；} — pas d'espace.
\end{itemize}
\end{aretenir}

\begin{savaistu}[Le « \zh{越} » dans la culture chinoise : du Classique au moderne]
Le caractère \zh{越} (yuè) signifiait à l'origine « passer outre, franchir » (clé \zh{走} « marcher » + phonétique \zh{戉} « hache »). Dans les textes classiques (\zh{孟子}, \zh{荀子}), \zh{越} sert de préposition « vers, au-delà » (\zh{越国} = le pays de Yue, au sud du Yangtsé). La structure \zh{越……越……} (de plus en plus…) est une grammaticalisation tardive (dynastie Tang-Song) d'une construction comparative : « plus on franchit [la limite], plus… ». Aujourd'hui, \zh{越来越} est l'un des marqueurs d'aspect les plus fréquents du chinois oral (fréquence > 2000 occurrences / million de mots). À Yaoundé comme à Pékin, on l'entend dans toutes les bouches : \zh{越来越堵}, \zh{越来越贵}, \zh{越来越好} — la langue vivante en mouvement !
\end{savaistu}

\newpage

\coursec{Activité d'intégration}

\coursec{Leçon 24 — Grammaire : 越来越；一方面……另一方面……}

\courssub{Objectifs de l'activité}
\noindent
À l'issue de cette leçon, l'élève doit être capable de :
\begin{itemize}
    \item \textbf{Mobiliser} les structures \cle{越来越} (de plus en plus) et \cle{一方面……另一方面……} (d'un côté… de l'autre côté…) dans une argumentation cohérente, à l'oral comme à l'écrit.
    \item \textbf{Construire} des phrases grammaticalement correctes en respectant l'ordre des mots chinois (Sujet + 越来越 + Adjectif/Verbe d'état ; 一方面，Proposition 1；另一方面，Proposition 2。).
    \item \textbf{Exprimer} un point de vue nuancé sur un sujet de société (l'apprentissage du chinois au Cameroun) en articulant avantages et difficultés.
    \item \textbf{Identifier} et \textbf{corriger} les erreurs fréquentes des francophones (intrusion de \cle{很}, omission de la virgule, confusion \cle{另一方面}/\cle{别的方面}).
    \item \textbf{Écrire} correctement les caractères clés de la leçon en respectant l'ordre des traits et en identifiant leurs clés (radicaux).
\end{itemize}

\courssub{Situation déclenchante : Journal mural du Club de Chinois}
\noindent
\textbf{Contexte :} Au lycée bilingue de \emph{Nkolbisson} (Yaoundé), la « Semaine des Langues Vivantes » a pour thème \og L'ouverture au monde par les langues \fg{}. La section de chinois LV2 anime un débat : \og \textbf{Faut-il apprendre le chinois au lycée ?} \fg{}. Trois élèves affichent leur avis sur le journal mural. Ce texte authentique sert de modèle langagier : il contient naturellement les deux structures ciblées.

\begin{textebox}[Journal mural — Lycée de Nkolbisson, Yaoundé]
\zh{同学甲：现在学中文的同学越来越多了。一方面，中国和喀麦隆合作很多，找工作有机会；另一方面，汉字太难写，语法也不容易。} \\
\zh{Tóngxué Jiǎ: Xiànzài xué Zhōngwén de tóngxué yuè lái yuè duō le. Yī fāngmiàn, Zhōngguó hé Kāmàilóng hézuò hěn duō, zhǎo gōngzuò yǒu jīhuì; lìng yī fāngmiàn, Hànzì tài nán xiě, yǔfǎ yě bù róngyì.} \\
\emph{« De plus en plus de camarades apprennent le chinois. D'un côté, la Chine et le Cameroun coopèrent beaucoup, il y a des opportunités d'emploi ; de l'autre, les caractères sont trop difficiles à écrire, la grammaire n'est pas facile non plus. »}

\zh{同学乙：我同意。中文课越来越有意思，老师常组织比赛。但是一方面要背很多词，另一方面没时间复习。} \\
\zh{Tóngxué Yǐ: Wǒ tóngyì. Zhōngwén kè yuè lái yuè yǒu yìsi, lǎoshī cháng zǔzhī bǐsài. Dànshì yī fāngmiàn yào bèi hěn duō cí, lìng yī fāngmiàn méi shíjiān fùxí.} \\
\emph{« Je suis d'accord. Les cours de chinois sont de plus en plus intéressants, le prof organise souvent des concours. Mais d'un côté il faut apprendre beaucoup de mots, de l'autre on n'a pas le temps de réviser. »}

\zh{同学丙：我觉得学中文越来越重要。一方面可以去中国留学，另一方面能了解不同文化。我要坚持学下去！} \\
\zh{Tóngxué Bǐng: Wǒ juéde xué Zhōngwén yuè lái yuè zhòngyào. Yī fāngmiàn kěyǐ qù Zhōngguó liúxué, lìng yī fāngmiàn néng liǎojiě bùtóng wénhuà. Wǒ yào jiānchí xué xiàqù!} \\
\emph{« Je trouve qu'apprendre le chinois est de plus en plus important. D'un côté on peut aller étudier en Chine, de l'autre on peut comprendre une culture différente. Je vais persévérer ! »}
\end{textebox}

\begin{popfigure}
\begin{tikzpicture}[
    node distance=1.2cm and 1.5cm,
    every node/.style={align=center, font=\small},
    box1/.style={draw, rounded corners, thick, fill=popblueL, minimum width=3.5cm, minimum height=1.2cm},
    box2/.style={draw, rounded corners, thick, fill=popgreenL, minimum width=3.5cm, minimum height=1.2cm},
    box3/.style={draw, rounded corners, thick, fill=poporangeL, minimum width=4.5cm, minimum height=1cm},
    box4/.style={draw, rounded corners, thick, fill=poppinkL, minimum width=4.5cm, minimum height=1cm},
    arrow/.style={-Stealth, thick, popdark}
]
% Structure 1
\node[box1] (s1) {\textbf{Structure 1} : \zh{越来越}};
\node[box3, below=of s1] (sch1) {Schéma : \zh{Sujet + 越来越 + Adj / Verbe psychologique}};
\node[box4, below=of sch1, align=center] (ex1){Ex. : \zh{中文越来越有用。}\\ \emph{Zhōngwén yuè lái yuè yǒuyòng.}};

% Structure 2
\node[box2, right=of s1] (s2) {\textbf{Structure 2} : \zh{一方面……另一方面……}};
\node[box3, below=of s2] (sch2) {Schéma : \zh{一方面，句子1； 另一方面，句子2。}};
\node[box4, below=of sch2, align=center] (ex2){Ex. : \zh{一方面有用，另一方面难。}\\ \emph{Yī fāngmiàn yǒuyòng, lìng yī fāngmiàn nán.}};

\draw[arrow] (s1) -- (sch1);
\draw[arrow] (sch1) -- (ex1);
\draw[arrow] (s2) -- (sch2);
\draw[arrow] (sch2) -- (ex2);
\end{tikzpicture}
\legende{Résumé visuel des deux structures grammaticales de la leçon 24.}
\end{popfigure}

\courssub{Consignes de l'activité intégrée (50 min)}
\noindent
Le travail se déroule en \textbf{quatre phases successives}. Le professeur chronomètre chaque étape.

\begin{enumerate}
    \item \textbf{Phase 1 — Compréhension et repérage (10 min) — Individuel.}
    \begin{itemize}
        \item Lisez le document dans la \texttt{textebox} ci-dessus.
        \item Surlignez (ou recopiez) \textbf{toutes les occurrences} de \cle{越来越} et de \cle{一方面……另一方面……} (variantes incluses).
        \item Traduisez mentalement les trois paragraphes. Notez le vocabulaire nouveau : \zh{合作} (hézuò), \zh{坚持} (jiānchí), \zh{留学} (liúxué), \zh{声调} (shēngdiào).
    \end{itemize}

    \item \textbf{Phase 2 — Préparation du débat en groupes (20 min) — Groupes de 4.}
    \begin{itemize}
        \item Chaque groupe représente une \textbf{équipe « Pour »} ou « Contre » (attribuée par le prof).
        \item \textbf{Tâche A (Oral) :} Préparez une intervention de \textbf{1 min 30 max}. Incluez impérativement :
        \begin{itemize}
            \item \textbf{2 phrases} avec \cle{越来越} (ex. \zh{学中文的人越来越多} / \zh{中文越来越有用}).
            \item \textbf{2 phrases} avec \cle{一方面……另一方面……} (un argument pour, un contre ou une nuance).
        \end{itemize}
        \item \textbf{Tâche B (Écrit) :} Rédigez vos 4 phrases sur feuille A3 : \textbf{caractères + pinyin}. Vérifiez l'ordre des mots (fiche Ressources).
        \item Désignez un \textbf{porte-parole} et un \textbf{secrétaire} (écrira les structures au tableau).
    \end{itemize}

    \item \textbf{Phase 3 — Débat public et vote (15 min) — Classe entière.}
    \begin{itemize}
        \item Les porte-paroles interviennent (alternance Pour/Contre). Le secrétaire note les structures au tableau.
        \item La classe vote sur bulletin pour \textbf{l'argument le mieux construit grammaticalement} (structures justes, prononciation claire, tons respectés).
        \item Équipe gagnante : « point bonus » participation orale.
    \end{itemize}

    \item \textbf{Phase 4 — Trace écrite individuelle et auto-correction (15 min + devoirs).}
    \begin{itemize}
        \item \textbf{En classe :} Rédigez un \textbf{paragraphe de 4 à 5 phrases} (votre avis personnel) réutilisant \textbf{au moins 1 fois} \cle{越来越} et \textbf{au moins 1 fois} \cle{一方面……另一方面……}.
        \item \textbf{Correction guidée (pairs) :} Échangez les cahiers. Grille de correction (basée sur les \texttt{attention}) : \textit{Ordre des mots OK ? / Tons notés ? / Ponctuation chinoise ? / Sens logique ? / Caractères lisibles ?}
        \item \textbf{Version finale (devoirs) :} Recopiez proprement (caractères + pinyin + traduction FR) pour la prochaine séance. \textbf{Barème :} 5 pts grammaire / 3 pts vocabulaire-cohérence / 2 pts écriture-caractères.
    \end{itemize}
\end{enumerate}

\courssub{Ressources à mobiliser : Fiches de référence}

\begin{propriete}[Structure \cle{越来越} — « De plus en plus »]
\textbf{Schéma syntaxique :} \cle{Sujet + 越来越 + Adjectif / Verbe d'état psychologique} \\
\textbf{Emploi :} Indique un changement progressif, une évolution continue vers un degré supérieur. Le processus est en cours.
\begin{itemize}
    \item \textbf{Adjectifs :} \zh{越来越多} (de plus en plus nombreux), \zh{越来越难} (de plus en plus difficile), \zh{越来越重要} (de plus en plus important).
    \item \textbf{Verbes d'état / psychologiques :} \zh{越来越喜欢} (aimer de plus en plus), \zh{越来越想去} (vouloir aller de plus en plus), \zh{越来越觉得} (trouver de plus en plus que).
\end{itemize}
\textbf{Particule \cle{了} :} Fréquente en fin de phrase pour marquer le changement \textbf{constaté} (\zh{越来越多了}), absente pour une vérité générale ou une opinion (\zh{越来越重要}).
\textbf{Exemples :}
\begin{itemize}
    \item \zh{学中文的人越来越多。} (Xué Zhōngwén de rén yuè lái yuè duō.) — De plus en plus de gens apprennent le chinois.
    \item \zh{汉字越来越难。} (Hànzì yuè lái yuè nán.) — Les caractères sont de plus en plus difficiles.
    \item \zh{我越来越喜欢喀麦隆的菜。} (Wǒ yuè lái yuè xǐhuan Kāmàilóng de cài.) — J'aime de plus en plus la cuisine camerounaise.
    \item \zh{中文课越来越有意思了。} (Zhōngwén kè yuè lái yuè yǒu yìsi le.) — Les cours de chinois deviennent de plus en plus intéressants (constat).
\end{itemize}
\end{propriete}

\begin{propriete}[Structure \cle{一方面……另一方面……} — « D'un côté… de l'autre… »]
\textbf{Schéma syntaxique :} \cle{一方面，句子1； 另一方面，句子2。} \\
\textbf{Règles de ponctuation :} \textbf{Virgule obligatoire} après \cle{一方面} et après \cle{另一方面}. Point-virgule \textbf{；} ou point \textbf{。} entre les deux propositions principales.
\textbf{Emploi :} Présenter deux aspects contrastés, complémentaires ou contradictoires d'un \textbf{même sujet}. Équivalent à \og D'une part… d'autre part \fg{} ou \og D'un côté… de l'autre \fg{}.
\textbf{Variante courante :} \cle{一方面……另外一方面……} (yī fāngmiàn… lìng wài yī fāngmiàn… — légèrement plus formel).
\textbf{Exemples :}
\begin{itemize}
    \item \zh{一方面，中文课很有用；另一方面，作业太多。} (Yī fāngmiàn, Zhōngwén kè hěn yǒuyòng; lìng yī fāngmiàn, zuòyè tài duō.) — D'un côté le cours est utile ; de l'autre, trop de devoirs.
    \item \zh{去中国留学：一方面学费便宜，另一方面签证容易。} (Qù Zhōngguó liúxué: yī fāngmiàn xuéfèi piányi, lìng yī fāngmiàn qiānzhèng róngyì.) — Étudier en Chine : d'un côté frais bas, de l'autre visa facile.
    \item \zh{学中文一方面开阔视野，另一方面增加就业机会。} (Xué Zhōngwén yī fāngmiàn kāikuò shìyě, lìng yī fāngmiàn zēngjiā jiùyè jīhuì.) — Apprendre le chinois élargit l'horizon d'un côté, augmente les chances d'emploi de l'autre.
\end{itemize}
\end{propriete}

\begin{attention}[Pièges fréquents des francophones — À afficher au tableau]
\begin{itemize}
    \item \textbf{Intrusion de \cle{很} (hěn) après \cle{越来越}} : \textbf{FAUX} \zh{*越来越很难} / \zh{*越来越很好}. \textbf{JUSTE} \zh{越来越难} / \zh{越来越好}. \textit{Explication :} \cle{越来越} exprime déjà le degré élevé, \cle{很} est redondant et grammaticalement impossible ici.
    \item \textbf{Omission de la virgule} après \cle{一方面} et \cle{另一方面} : \textbf{FAUX} \zh{*一方面中文有用}。 \textbf{JUSTE} \zh{一方面，中文有用。} \textit{Règle :} Ce sont des connecteurs de phrase (adverbes conjonctifs), ils sont suivis d'une pause.
    \item \textbf{Confusion \cle{另一方面} vs \cle{别的方面}} : \cle{另一方面} = l'autre facette du \textbf{même} sujet. \cle{别的方面} = un \textbf{autre} sujet différent. \textit{Ex.} \zh{一方面累，另一方面快乐} (même sujet : le travail) vs \zh{别的方面我不懂} (autre domaine).
    \item \textbf{Ordre des mots avec \cle{越来越}} : Ne pas placer le sujet après. \textbf{FAUX} \zh{*越来越多的人学中文} (structure relative possible mais pas la structure de base apprise ici). \textbf{JUSTE} \zh{学中文的人越来越多} (Sujet \zh{学中文的人} + \cle{越来越多}).
    \item \textbf{Pinyin / Tons :} \cle{越来越} se lit \emph{yuè lái yuè} (4-2-4), pas \emph{yuè lái yuè} (4-4-4). \cle{一方面} = \emph{yī fāngmiàn} (1-1-4), \cle{另一方面} = \emph{lìng yī fāngmiàn} (4-1-1-4). Attention au changement de ton de \cle{一} (yī → yí avant ton 4, mais ici isolé ou avant \cle{方} ton 1 → reste \emph{yī}).
\end{itemize}
\end{attention}

\begin{aretenir}[Écriture des caractères clés — Ordre des traits et Clés (Radicaux)]
\begin{itemize}
    \item \zh{越} (yuè) — Clé : \zh{走} (zǒu, « marche ») en bas. Ordre : haut \zh{戈} (gē) → bas \zh{走} (traits horizontaux puis verticaux, enveloppement final).
    \item \zh{来} (lái) — Clé : \zh{木} (mù, « bois ») en bas (forme ancienne). Ordre : haut \zh{米} (mǐ) → bas \zh{木}. Attention : trait horizontal long final du \zh{木} traverse.
    \item \zh{方} (fāng) — Clé : \zh{方} (lui-même). Ordre : point → horizontal → vertical → horizontal fermant (boîte).
    \item \zh{面} (miàn) — Clé : \zh{面} (lui-même). Ordre : \zh{方} dans le haut → \zh{一} (trait horizontal) au milieu → quatre traits \zh{丨} verticaux descendants → trait horizontal final fermant.
    \item \zh{其} (qí) — Clé : \zh{八} (bā). Ordre : haut \zh{甘} (gān) → bas \zh{八}. Composant de \zh{其} dans \zh{期} (qī), \zh{棋} (qí), \zh{奇} (qí).
    \item \zh{他} (tā) — Clé : \zh{亻} (rén, « homme » à gauche). Ordre : \zh{亻} (traits 1-2) → \zh{也} (yě, traits 3-5). Règle : gauche avant droite.
\end{itemize}
\end{aretenir}

\begin{center}
\begin{tabularx}{\linewidth}{|X|X|X|X|}
\hline
\rowcolor{popblueL} \textbf{Caractères} & \textbf{Pinyin} & \textbf{Nature} & \textbf{Traduction} \\
\hline
合作 & hézuò & verbe / nom & coopérer / coopération \\
\hline
机会 & jīhuì & nom & opportunité, chance \\
\hline
留学 & liúxué & verbe & étudier à l'étranger \\
\hline
坚持 & jiānchí & verbe & persévérer, insister \\
\hline
有用 & yǒuyòng & adjectif & utile \\
\hline
签证 & qiānzhèng & nom & visa \\
\hline
学费 & xuéfèi & nom & frais de scolarité \\
\hline
便宜 & piányi & adjectif & pas cher, bon marché \\
\hline
声调 & shēngdiào & nom & tons (de la langue) \\
\hline
开阔视野 & kāikuò shìyě & loc. verbale & élargir ses horizons \\
\hline
就业 & jiùyè & verbe / nom & emploi, trouver du travail \\
\hline
\end{tabularx}
\end{center}

\courssub{Exercices d'application et de transformation (Travail en classe / Devoirs)}

\begin{exercice}[Exercice 1 : Transformation — \cle{越来越}]
\noindent
Transformez les phrases en utilisant \cle{越来越} + adjectif/verbe. Attention à l'ordre des mots et à l'absence de \cle{很}.
\begin{enumerate}
    \item 中文很难。 $\rightarrow$ \zh{中文越来越难。}
    \item 学中文的人很多。 $\rightarrow$ \zh{学中文的人越来越多。}
    \item 我喜欢吃 ndolé。 $\rightarrow$ \zh{我越来越喜欢吃 ndolé。}
    \item 现在的汉字很有意思。 $\rightarrow$ \zh{现在的汉字越来越有意思。}
    \item 中国企业在杜阿拉很多。 $\rightarrow$ \zh{中国企业在杜阿拉越来越多。}
\end{enumerate}
\end{exercice}

\begin{exercice}[Exercice 2 : Assemblage — \cle{一方面……另一方面……}]
\noindent
Reliez les deux propositions par \cle{一方面……另一方面……}. Pensez à la virgule et au point-virgule.
\begin{enumerate}
    \item Apprendre le chinois est utile. / Les caractères sont difficiles.
    \item Aller en Chine coûte peu cher. / Le visa est facile à obtenir.
    \item Le prof organise des jeux. / Il y a beaucoup de devoirs à la maison.
    \item On peut travailler pour une boîte chinoise. / Il faut parler couramment.
\end{enumerate}
\end{exercice}

\begin{exercice}[Exercice 3 : Traduction (Français $\rightarrow$ Chinois)]
\noindent
Traduisez en caractères + pinyin. Utilisez les structures de la leçon.
\begin{enumerate}
    \item De plus en plus d'élèves camerounais apprennent le chinois.
    \item D'un côté, la grammaire chinoise est simple ; de l'autre, les tons sont difficiles.
    \item Je trouve que le chinois devient de plus en plus important pour l'avenir.
    \item D'un côté, je veux aller en Chine ; de l'autre, je n'ai pas assez d'argent.
\end{enumerate}
\end{exercice}

\begin{exercice}[Exercice 4 : Production guidée — « Mon avis »]
\noindent
Complétez le paragraphe ci-dessous avec vos idées personnelles (vocabulaire du tableau + structures ciblées).
\zh{我认为在喀麦隆学中文 \underline{\hspace{3cm}}。} \\
\zh{\underline{\hspace{1cm}}，\underline{\hspace{3cm}}；\underline{\hspace{1cm}}，\underline{\hspace{3cm}}。} \\
\zh{不过，中文课 \underline{\hspace{3cm}}。} \\
\zh{所以我要 \underline{\hspace{3cm}}。}
\end{exercice}

\begin{corrige}[Exercices 1 à 4]
\textbf{Exercice 1 :}
1. \zh{中文越来越难。} (Zhōngwén yuè lái yuè nán.) \\
2. \zh{学中文的人越来越多。} (Xué Zhōngwén de rén yuè lái yuè duō.) \\
3. \zh{我越来越喜欢吃 ndolé。} (Wǒ yuè lái yuè xǐhuan chī ndolé.) \\
4. \zh{现在的汉字越来越有意思。} (Xiànzài de Hànzì yuè lái yuè yǒu yìsi.) \\
5. \zh{中国企业在杜阿拉越来越多。} (Zhōngguó qǐyè zài Dū'ālā yuè lái yuè duō.)

\textbf{Exercice 2 :}
1. \zh{一方面，学中文很有用；另一方面，汉字很难。} \\
2. \zh{一方面，去中国学费便宜；另一方面，签证容易。} \\
3. \zh{一方面，老师常组织游戏；另一方面，家庭作业太多。} \\
4. \zh{一方面，能在中资企业工作；另一方面，得说得流利。}

\textbf{Exercice 3 :}
1. \zh{喀麦隆学生学中文的越来越多。} (Kāmàilóng xuéshēng xué Zhōngwén de yuè lái yuè duō.) \\
2. \zh{一方面，中文语法简单；另一方面，声调很难。} (Yī fāngmiàn, Zhōngwén yǔfǎ jiǎndān; lìng yī fāngmiàn, shēngdiào hěn nán.) \\
3. \zh{我觉得中文对将来越来越重要。} (Wǒ juéde Zhōngwén duì jiānglái yuè lái yuè zhòngyào.) \\
4. \zh{一方面我想去中国；另一方面钱不够。} (Yī fāngmiàn wǒ xiǎng qù Zhōngguó; lìng yī fāngmiàn qián bù gòu.)

\textbf{Exercice 4 (Proposition de remplissage) :}
\zh{我认为在喀麦隆学中文 \textbf{越来越重要}。} \\
\zh{\textbf{一方面}，\textbf{中国企业在雅温得越来越多，找工作有机会}；\textbf{另一方面}，\textbf{汉字和声调太难，要花很多时间练习}。} \\
\zh{不过，中文课 \textbf{越来越有意思}。} \\
\zh{所以我要 \textbf{坚持学下去，将来去中国留学}。}
\end{corrige}

\courssub{Proposition de production et corrigé commenté (Phase 4 — Modèle élève)}

\begin{exoresolu}[Production modèle — Trace écrite individuelle (Niveau HSK 2-3)]
\textbf{Consigne :} Rédigez un paragraphe de 4-5 phrases donnant votre avis sur l'apprentissage du chinois au lycée, en utilisant \cle{越来越} et \cle{一方面……另一方面……}.

\tcblower

\textbf{Production élève (Modèle) :}

\zh{我认为在喀麦隆学中文越来越重要。} \\
\zh{Wǒ rènwéi zài Kāmàilóng xué Zhōngwén yuè lái yuè zhòngyào.} \\
\emph{« Je pense qu'au Cameroun, apprendre le chinois devient de plus en plus important. »}

\zh{一方面，中国企业在杜阿拉和雅温得越来越多，找工作的机会大。} \\
\zh{Yī fāngmiàn, Zhōngguó qǐyè zài Dū'ālā hé Yǎwēndé yuè lái yuè duō, zhǎo gōngzuò de jīhuì dà.} \\
\emph{« D'un côté, les entreprises chinoises sont de plus en plus nombreuses à Douala et Yaoundé, les opportunités d'emploi sont grandes. »}

\zh{另一方面，汉字和声调很难，需要花很多时间练习。} \\
\zh{Lìng yī fāngmiàn, Hànzì hé shēngdiào hěn nán, xūyào huā hěn duō shíjiān liànxí.} \\
\emph{« De l'autre côté, les caractères et les tons sont difficiles, il faut y consacrer beaucoup de temps pour s'exercer. »}

\zh{不过，中文课越来越有意思，我们常唱歌、比赛。} \\
\zh{Búguò, Zhōngwén kè yuè lái yuè yǒu yìsi, wǒmen cháng chànggē, bǐsài.} \\
\emph{« Cependant, les cours de chinois sont de plus en plus intéressants, on chante souvent, on fait des concours. »}

\zh{所以我要坚持学下去，将来想去中国留学。} \\
\zh{Suǒyǐ wǒ yào jiānchí xué xiàqù, jiānglái xiǎng qù Zhōngguó liúxué.} \\
\emph{« Donc je vais persévérer dans l'apprentissage, plus tard je veux aller étudier en Chine. »}

\textbf{Commentaire des choix de langue (Grille d'analyse pour l'élève) :}

\begin{enumerate}
    \item \textbf{Phrase 1} : \cle{越来越重要} — Structure correcte (Sujet \zh{在喀麦隆学中文} + \cle{越来越} + Adj). Pas de \cle{了} car énoncé général/opinion.
    \item \textbf{Phrase 2} : \cle{一方面} introduit l'argument positif. \textbf{Bonus :} Emploi d'une \textbf{deuxième occurrence de \cle{越来越多}} à l'intérieur de la proposition (\zh{企业……越来越多}) pour montrer la maîtrise de la progression. \cle{机会大} (collocation HSK 3 : « grande opportunité »).
    \item \textbf{Phrase 3} : \cle{另一方面} introduit la concession/la difficulté. Vocabulaire précis : \zh{声调} (tons), \zh{花时间} (passer du temps), \zh{练习}. Pas de \cle{了} car contrainte générale.
    \item \textbf{Phrase 4} : Transition \cle{不过} (mais/cependant). \textbf{Troisième occurrence de \cle{越来越有意思}} — montre la réutilisation spontanée. \zh{常} (souvent) placé avant le verbe.
    \item \textbf{Phrase 5} : Conclusion logique \cle{所以} (donc). \zh{坚持学下去} (persévérer à apprendre / continuer d'apprendre). \zh{将来} (à l'avenir) + \zh{想} (projet).
    \item \textbf{Global :} Ponctuation chinoise respectée (， 。). Pinyin noté pour l'auto-correction. Caractères lisibles (ordre des traits respecté).
\end{enumerate}
\end{exoresolu}

\courssub{Bilan de l'activité et perspectives}
\noindent
Cette activité d'intégration permet de \textbf{valider l'acquisition opérationnelle} des deux structures grammaticales du Module 3 dans une situation de communication authentique et motivante pour des lycéens camerounais.

\begin{itemize}
    \item \textbf{Sur le plan grammatical :} Le passage de la manipulation guidée (exercices 1-3) à la \textbf{production libre encadrée} (contraintes quantitatives : 2 phrases par structure en Phase 2, 1+1 en Phase 4) force l'élève à internaliser le schéma syntaxique. L'auto-correction croisée (Phase 4) développe la \textbf{métacognition linguistique} (repérer l'absence de \cle{很}, la place de la virgule, l'ordre Sujet-\cle{越来越}-Adj).
    \item \textbf{Sur le plan lexical :} Le contexte \og Chine-Cameroun \fg{} (entreprises à Douala/Yaoundé, études en Chine, visa, frais, ndolé) ancre le vocabulaire HSK dans une réalité vécue, favorisant la mémorisation à long terme.
    \item \textbf{Sur le plan des compétences (référentiel MINESEC) :}
    \begin{itemize}
        \item \emph{Compréhension de l'écrit :} Lecture du journal mural (texte authentique simulé, registre familier/standard).
        \item \emph{Expression orale :} Débat public, gestion du trac, prononciation des tons (surveillée sur \cle{越来越} \emph{yuè lái yuè} et \cle{一方面} \emph{yī fāngmiàn}).
        \item \emph{Expression écrite :} Paragraphe argumenté, cohérence énonciative (connecteurs \cle{不过}, \cle{所以}).
        \item \emph{Compétence socioculturelle :} Prise de conscience des enjeux réels de la relation bilatérale (emploi, culture, difficulté de la langue).
    \end{itemize}
    \item \textbf{Pistes de différenciation / Remédiation :}
    \begin{itemize}
        \item \textit{Élèves en difficulté :} Fournir une « phrase à trous » pour la trace écrite (ex. \zh{我觉得学中文越来越\_\_\_\_。一方面\_\_\_\_，另一方面\_\_\_\_。}) + banque de mots (重要、有用、难、多、机会、时间).
        \item \textit{Élèves avancés :} Demander d'ajouter une structure comparative \cle{比} (ex. \zh{学中文比学法语更有用}) ou la structure \cle{越……越……} (ex. \zh{越学越有意思}) en extension. Proposer un rôle de « modérateur » du débat pour synthétiser les arguments.
    \end{itemize}
    \item \textbf{Évaluation formative :} La grille de correction utilisée en Phase 4 (vérification par les pairs) sera collectée par le professeur pour identifier les erreurs récurrentes (ex. confusion \cle{另一方面} / \cle{别的方面}, omission de la virgule, intrusion de \cle{很}) et planifier une séance de \textbf{révision ciblée} (Module \texttt{revision}) avant l'évaluation sommative.
\end{itemize}

\noindent
\cle{Objectif atteint :} L'élève est désormais capable de \textbf{nuancer son propos} en chinois, outil indispensable pour le niveau B1 (HSK 4) visé en Terminale.

\newpage

\coursec{Méthodes et exercices résolus}

\courssub{Méthodes et exercices résolus}

\begin{methode}[Construire une phrase avec \cle{越来越} pour exprimer un changement progressif]
\textbf{Objectif :} Maîtriser l'ordre des mots et la nature du prédicat (adjectif ou verbe + complément) après \zh{越来越} (yuè lái yuè, « de plus en plus »).
\begin{enumerate}
    \item \textbf{Identifier le sujet} (personne, chose, situation ou phrase verbale).
    \item \textbf{Choisir le prédicat et placer \zh{越来越}} :
        \begin{itemize}
            \item \textbf{Adjectif qualificatif} (ex. : \zh{热} rè, \zh{忙} máng, \zh{好} hǎo, \zh{难} nán) 
            $\rightarrow$ \zh{Sujet + 越来越 + Adj.} \textit{(Pas de \zh{的} entre \zh{越来越} et l'adjectif).}
            \item \textbf{Verbe d'action + complément de degré} (ex. : \zh{跑得快} pǎo de kuài, \zh{写得好} xiě de hǎo) 
            $\rightarrow$ \zh{Sujet + Verbe + 得 + 越来越 + Complément.} \textit{\zh{越来越} s'insère \textbf{après} \zh{得}, devant l'adjectif/adverbe du complément.}
            \item \textbf{Verbe d'état psychologique / statif} (\zh{喜欢} xǐhuan, \zh{爱} ài, \zh{讨厌} tǎoyàn, \zh{知道} zhīdào) 
            $\rightarrow$ \zh{Sujet + 越来越 + Verbe (+ Objet).} \textit{Pas de \zh{得} car ce ne sont pas des verbes d'action dynamiques.}
        \end{itemize}
    \item \textbf{Position de \zh{越来越}} : Immédiatement devant le prédicat (adjectif) ou devant le complément de degré (après \zh{得}). Jamais en début de phrase avant le sujet, jamais après le prédicat.
    \item \textbf{Accord tonal} : \zh{越} (4e ton) \zh{来} (2e ton) \zh{越} (4e ton). Ne pas confondre \zh{越} (progression) et \zh{月} (yuè, lune/mois).
    \item \textbf{Particule \zh{了} en fin de phrase} : Fréquente (quasi systématique à l'oral) pour marquer le changement d'état constaté / nouvelle situation.
\end{enumerate}
\cle{Règle d'or} : \zh{越来越} est un bloc insécable. Pas de \zh{的} après \zh{越来越} devant un adjectif. Structure verbe : \zh{V + 得 + 越来越 + Compl.}
\end{methode}

\begin{exoresolu}[Application : Transformation et production avec \cle{越来越}]
\textbf{Énoncé :} 
\begin{enumerate}
    \item Transformer les phrases suivantes en utilisant \zh{越来越} (ajoutez \zh{了} si pertinent) :
    \begin{enumerate}
        \item 现在天气热。 
        \item 他跑得快。 
        \item 学中文有意思。 
        \item 喀麦隆的足球队强。 
    \end{enumerate}
    \item Traduire en chinois (caractères + pinyin) :
    \begin{enumerate}
        \item Yaoundé devient de plus en plus belle.
        \item Ma petite sœur mange de plus en plus vite.
        \item Étudier le chinois est de plus en plus difficile.
    \end{enumerate}
\end{enumerate}

\tcblower
\textbf{Solution détaillée :}

\textbf{1. Transformation :}
\begin{itemize}
    \item[a)] \zh{现在天气越来越热了。} (Xiànzài tiānqì yuè lái yuè rè le.) — « Maintenant le temps est de plus en plus chaud. » \textit{Adj. \zh{热} placé après \zh{越来越} ; \zh{了} marque le changement d'état.}
    \item[b)] \zh{他跑得越来越快了。} (Tā pǎo de yuè lái yuè kuài le.) — « Il court de plus en plus vite. » \textit{Structure V+\zh{得}+Compl. : \zh{越来越} s'insère \textbf{entre} \zh{得} et le complément \zh{快}.}
    \item[c)] \zh{学中文越来越有意思了。} (Xué Zhōngwén yuè lái yuè yǒu yìsi le.) — « Apprendre le chinois est de plus en plus intéressant. » \textit{Sujet = phrase verbale \zh{学中文} ; prédicat = \zh{有意思} (adj. composé).}
    \item[d)] \zh{喀麦隆的足球队越来越强了。} (Kāmàilóng de zúqiúduì yuè lái yuè qiáng le.) — « L'équipe de foot du Cameroun est de plus en plus forte. »
\end{itemize}

\textbf{2. Traduction (Thème) :}
\begin{itemize}
    \item[a)] \zh{雅温得越来越美了。} (Yǎwēndé yuè lái yuè měi le.) — \textit{Sujet propre \zh{雅温得} + \zh{越来越} + adj. \zh{美}.}
    \item[b)] \zh{我妹妹吃饭吃得越来越快了。} (Wǒ mèimei chīfàn chī de yuè lái yuè kuài le.) — \textit{Verbe \zh{吃} répété devant \zh{得} (structure V+O, V+de+Compl). \zh{越来越} après \zh{得}.}
    \item[c)] \zh{学中文越来越难了。} (Xué Zhōngwén yuè lái yuè nán le.) — \textit{Sujet verbal \zh{学中文} + \zh{越来越} + adj. \zh{难}.}
\end{itemize}

\textbf{Conclusion :} La position de \zh{越来越} est fixe : devant l'adjectif ou devant le complément de degré (après \zh{得}). L'erreur majeure est \zh{*他越来越跑得快} (ordre incorrect) ou \zh{*他跑越来越快得} (particule mal placée).
\end{exoresolu}

\begin{methode}[Structurer un argumentaire avec \cle{一方面……另一方面……}]
\textbf{Objectif :} Produire un énoncé complexe présentant deux facettes d'un même sujet (contraste, concession, énumération d'arguments opposés).
\begin{enumerate}
    \item \textbf{Analyser le sujet} : Identifier les deux aspects distincts (A et B) d'un même thème. Ex. : « Étudier en Chine : avantages / difficultés » ; « Réseaux sociaux : utiles / dangereux ».
    \item \textbf{Construire la première clause} : \zh{一方面} (yī fāng miàn) + Sujet 1 + Prédicat 1. \textit{Souvent aspect positif, cause, ou premier argument. \zh{一} se prononce \textbf{yí} (2e ton) devant \zh{面} (4e ton).}
    \item \textbf{Construire la deuxième clause} : \zh{另一方面} (lìng wài yī fāng miàn) + Sujet 2 + Prédicat 2. \textit{Souvent aspect négatif, conséquence, ou argument opposé. \zh{另} (lìng) 4e ton ; \zh{一} se prononce \textbf{yí} devant \zh{面}.}
    \item \textbf{Vérifier la position} : \zh{一方面} et \zh{autre方面} sont des \textbf{adverbiaux de thème} : ils ouvrent la clause (position initiale). Pas de sujet avant eux.
    \item \textbf{Ponctuation} : Point-virgule \zh{；} ou point \zh{。}, entre les deux clauses complètes. Virgule \zh{，} après \zh{一方面} et \zh{另一方面} si la clause est longue.
    \item \textbf{Cohérence} : Les deux aspects doivent porter sur le \textbf{même sujet général}. Éviter les digressions.
    \item \textbf{Variante} : \zh{一方面……另外一方面……} (plus insistant sur le second aspect, moins figé).
\end{enumerate}
\cle{Repère français} : « D'un côté... de l'autre côté » ; « En premier lieu... en second lieu » ; « Si... toutefois » ; « Certes... mais ».
\end{methode}

\begin{exoresolu}[Production écrite : Argumenter sur la vie d'étudiant]
\textbf{Énoncé :} Rédigez un court paragraphe (60-80 caractères) sur le thème : « Étudier le chinois au Cameroun : opportunités et défis ». Utilisez \zh{一方面……另一方面……} et au moins une fois \zh{越来越}. Vocabulaire imposé : \zh{机会} (jīhuì), \zh{困难} (kùnnán), \zh{坚持} (jiānchí).

\tcblower
\textbf{Solution détaillée :}

\textbf{Proposition de rédaction :}
\begin{textebox}[Production élève]
一方面，在喀麦隆学中文机会越来越多，可以去中国留学；\\
另一方面，语法和汉字有很多困难，但我们要坚持学习。\\
Yī fāng miàn, zài Kāmàilóng xué Zhōngwén jīhuì yuè lái yuè duō, kěyǐ qù Zhōngguó liúxué; \\
lìng wài yī fāng miàn, yǔfǎ hé hànzì yǒu hěnduō kùnnán, dàn wǒmen yào jiānchí xuéxí.
\end{textebox}
\textbf{Traduction :} D'un côté, au Cameroun, les opportunités d'apprendre le chinois sont de plus en plus nombreuses, on peut aller étudier en Chine ; de l'autre côté, la grammaire et les caractères comportent beaucoup de difficultés, mais nous devons persévérer dans l'étude.

\textbf{Analyse grammaticale :}
\begin{itemize}
    \item \zh{一方面} introduit la clause 1 (aspect positif : opportunités). Sujet implicite « nous/les étudiants ».
    \item \zh{机会越来越多} : structure \zh{越来越} + adj. \zh{多} (correct).
    \item \zh{另一方面} introduit la clause 2 (aspect négatif : difficultés). Transition claire.
    \item \zh{但我们要坚持学习} : coordination avec \zh{但} (dàn) pour marquer l'opposition interne à la 2e clause.
    \item Ponctuation : point-virgule \zh{；} entre les deux grandes clauses, point final \zh{。}
\end{itemize}

\textbf{Variante possible (plus formelle, niveau HSK 4) :}
\zh{一方面，中文在喀麦隆越来越受欢迎，提供了就业机会；另一方面，缺乏语言环境是一个大困难，需要同学们坚持不懈。}
(Yī fāng miàn, Zhōngwén zài Kāmàilóng yuè lái yuè shòu huān yíng, tígōng le jiùyè jīhuì; lìng wài yī fāng miàn, quēfá yǔyán huánjìng shì yí gè dà kùnnán, xūyào tóngxuémen jiānchí bùxiè.)
\end{exoresolu}

\begin{methode}[Réussir un exercice de transformation (Thème / Version / Restructuration)]
\textbf{Objectif :} Passer du français au chinois (thème) ou restructurer une phrase chinoise sans en changer le sens, en ciblant les deux structures du jour.
\begin{enumerate}
    \item \textbf{Analyser la structure source (français)} : Repérer « de plus en plus » $\rightarrow$ \zh{越来越} ; « d'un côté... de l'autre » $\rightarrow$ \zh{一方面……另一方面……}.
    \item \textbf{Identifier le noyau syntaxique chinois} : 
        \begin{itemize}
            \item « De plus en plus + adj » $\rightarrow$ \zh{越来越 + Adj}.
            \item « De plus en plus + adv » (vite, bien, tôt) $\rightarrow$ \zh{Verbe + 得 + 越来越 + Adv}.
            \item « D'un côté A, de l'autre B » $\rightarrow$ \zh{一方面 A，另一方面 B}.
        \end{itemize}
    \item \textbf{Gérer les mots-outils} : 
        \begin{itemize}
            \item \zh{了} (changement d'état) souvent nécessaire avec \zh{越来越}.
            \item \zh{得} (de) obligatoire devant le complément de degré si verbe d'action.
            \item Pas de \zh{的} après \zh{越来越}.
        \end{itemize}
    \item \textbf{Contrôler l'ordre des mots (Compléments avant verbe)} : Lieu/Temps précèdent le verbe en chinois (ex. \zh{在喀麦隆学} vs « apprendre au Cameroun »).
    \item \textbf{Relire à voix haute (tonal check)} : \zh{越(4) 来(2) 越(4)} ; \zh{一(2) 方(4) 面(4)} / \zh{另(4) 外(4) 一(2) 方(4) 面(4)}.
\end{enumerate}
\end{methode}

\begin{exoresolu}[Exercices de transformation type Bac]
\textbf{Énoncé :} 
\begin{enumerate}
    \item \textbf{Version (Chinois $\rightarrow$ Français) :} Traduire : \zh{一方面现在的手机功能越来越多，另一方面价格也越来越贵。}
    \item \textbf{Thème (Français $\rightarrow$ Chinois) :} « D'un côté, les professeurs sont de plus en plus exigeants ; de l'autre côté, les élèves sont de plus en plus stressés. »
    \item \textbf{Restructuration :} Réécrire la phrase suivante en utilisant \zh{一方面……另一方面……} : \zh{因为中文难，所以学生少；因为中文用处大，所以学生多。}
\end{enumerate}

\tcblower
\textbf{Solution détaillée :}

\textbf{1. Version :}
\textit{« D'un côté, les fonctions des téléphones actuels sont de plus en plus nombreuses ; de l'autre côté, les prix sont aussi de plus en plus chers. »}
\begin{itemize}
    \item \zh{现在的手机} (téléphones actuels) = Sujet 1.
    \item \zh{功能越来越多} (fonctions de plus en plus nombreuses) = Prédicat 1 (\zh{越来越} + adj \zh{多}).
    \item \zh{价格也越来越贵} (prix aussi de plus en plus chers) = Prédicat 2 (\zh{也} + \zh{越来越} + adj \zh{贵}).
\end{itemize}

\textbf{2. Thème :}
\zh{一方面，老师越来越严格；另一方面，学生压力越来越大。}
(Yī fāng miàn, lǎoshī yuè lái yuè yángé; lìng wài yī fāng miàn, xuésheng yālì yuè lái yuè dà.)
\begin{itemize}
    \item \textbf{Correction pédagogique :} « Les profs sont de plus en plus exigeants » $\rightarrow$ \zh{老师越来越严格} (Les profs sont de plus en plus stricts) ou \zh{老师的要求越来越严格} (Les exigences des profs...). On évite \zh{老师要求越来越严格} (ambigu : « les profs exigent de plus en plus strictement » $\rightarrow$ nécessite \zh{得}).
    \item \zh{压力越来越大} (pression de plus en plus grande) : collocation standard.
\end{itemize}

\textbf{3. Restructuration :}
\zh{一方面，因为中文难，所以学生少；另一方面，因为中文用处大，所以学生多。}
(Yī fāng miàn, yīnwèi Zhōngwén nán, suǒyǐ xuésheng shǎo; lìng wài yī fāng miàn, yīnwèi Zhōngwén yòngchù dà, suǒyǐ xuésheng duō.)
\begin{itemize}
    \item On conserve la causalité interne (\zh{因为……所以……}) à l'intérieur de chaque aspect.
    \item \zh{用处大} (yòngchù dà) = très utile / avoir beaucoup d'utilité.
\end{itemize}
\end{exoresolu}

\begin{methode}[Préparer une prise de parole continue (Exposé / Débat) avec les structures cibles]
\textbf{Objectif :} Mobiliser \zh{越来越} et \zh{一方面……另一方面……} dans une production orale fluide (1 à 2 min), notée au Bac (Épreuve orale / Contrôle continu).
\begin{enumerate}
    \item \textbf{Choisir un sujet d'actualité camerounais ou sino-camerounais} : Ex. « L'influence de la culture chinoise au Cameroun », « Les réseaux sociaux chez les jeunes », « L'apprentissage des langues étrangères », « Le football camerounais ».
    \item \textbf{Rédiger un canevas (pas un texte lu)} : 
        \begin{itemize}
            \item Introduction : Sujet + \zh{越来越} + adj (\zh{受欢迎} / \zh{重要} / \zh{现代化}).
            \item Développement : \zh{一方面} (avantages/exemples concrets : Yaoundé, Douala, Olembé, Japoma) + \zh{另一方面} (défis/risques).
            \item Conclusion : Opinion personnelle + \zh{坚持} / \zh{努力} / \zh{建议} / \zh{应该……才能……}.
        \end{itemize}
    \item \textbf{Entraîner les points de prononciation critiques} :
        \begin{itemize}
            \item Enchaînement \zh{越来越} (yuè lái yuè) : tons 4-2-4, liaison fluide, pas de pause.
            \item \zh{一方面} (yí fāng miàn) : \zh{一} $\rightarrow$ yí (devant 4e ton). \zh{另一方面} (lìng wài yí fāng miàn).
        \end{itemize}
    \item \textbf{S'enregistrer / S'auto-évaluer} : Vérifier l'absence de « euh » en français, la cohérence logique, l'usage correct de \zh{得} avec verbes d'action.
    \item \textbf{Anticiper les questions du jury} : « Pourquoi \zh{一方面} pas \zh{一个方面} ? » (Réponse : locution figée, \zh{一} côté « un/premier » figé, pas numéral).
\end{enumerate}
\end{methode}

\begin{exoresolu}[Simulation orale : Sujet « Le football camerounais »]
\textbf{Énoncé :} L'élève doit parler 1 min 30 sur : « L'évolution du football au Cameroun ». Il a 10 min de préparation. Proposez un canevas structuré et un extrait de production orale attendue (niveau HSK 3/4).

\tcblower
\textbf{Canevas de préparation (Fiche élève) :}
\begin{center}
\begin{tabularx}{\linewidth}{|X|X|X|}\hline
\rowcolor{popblueL} \textbf{Étape} & \textbf{Mots-clés / Structures cibles} & \textbf{Idées (Cameroun)} \\ \hline
Intro & \zh{越来越受欢迎} (yuè lái yuè shòu huān yíng) & Championnat MTN Elite One, Lions Indomptables, CAN. \\ \hline
Aspect 1 (\zh{一方面}) & \zh{基础设施越来越好} (jīchǔ shèshī yuè lái yuè hǎo) & Stades Japoma, Olembé ; centres de formation. \\ \hline
Aspect 2 (\zh{另一方面}) & \zh{管理还有困难} (guǎnlǐ hái yǒu kùnnán) & Âge des joueurs, organisation, arbitrage. \\ \hline
Conclusion & \zh{应该... 才能...} (yīnggāi... cáinéng...) & Investir dans la jeunesse, transparence. \\ \hline
\end{tabularx}
\end{center}

\textbf{Extrait de production orale attendue (Modèle) :}
\begin{textebox}[Modèle oral - 1min30]
大家好。今天我想谈谈喀麦隆的足球。\\
Dàjiā hǎo. Jīntiān wǒ xiǎng tántan Kāmàilóng de zúqiú.\\
\textit{Bonjour à tous. Aujourd'hui je veux parler du football camerounais.}

一方面，足球在喀麦隆\zh{越来越受欢迎}。新建的体育场，比如奥莱姆贝和贾波马，\zh{设施越来越现代化}。很多年轻人\zh{梦想成为}埃托奥或阿布巴卡尔。\\
Yī fāng miàn, zúqiú zài Kāmàilóng yuè lái yuè shòu huān yíng. Xīn jiàn de tǐyùchǎng, bǐrú Àoláimǔbèi hé Jiǎbōmǎ, shèshī yuè lái yuè xiàndàihuà. Hěnduō niánqīng rén mèngxiǎng chéngwéi Àituō'ào huò Ābùbākǎ'ěr.

另一方面，\zh{职业联赛的管理}还有\zh{不少问题}。比如年龄造假、比赛组织有时不规范。这\zh{影响了}国家队的成绩。\\
Lìng wài yī fāng miàn, zhíyè liánsài de guǎnlǐ hái yǒu bùshǎo wèntí. Bǐrú niánlíng zàojiǎ, bǐsài zǔzhī yǒushí bù guīfàn. Zhè yǐngxiǎng le guójiāduì de chéngjī.

总之，\zh{如果}我们\zh{想让}足球\zh{更强}，\zh{就必须}从青训和管理\zh{做起}。谢谢。\\
Zǒngzhī, rúguǒ wǒmen xiǎng ràng zúqiú gèng qiáng, jiù bìxū cóng qīngxùn hé guǎnlǐ zuòqǐ. Xièxie.
\end{textebox}

\textbf{Analyse didactique :}
\begin{itemize}
    \item \zh{越来越} utilisé 3 fois (adj \zh{受欢迎}, \zh{现代化}, structure V+得 implicite dans \zh{梦想成为} non concerné, mais \zh{越来越} modifie bien les adjectifs).
    \item \zh{一方面……另一方面……} structure le discours complet.
    \item Vocabulaire contextuel : \zh{奥莱姆贝} (Olembé), \zh{贾波马} (Japoma), \zh{年龄造假} (falsification d'âge), \zh{青训} (formation jeunes).
    \item Connecteurs logiques : \zh{比如} (par ex.), \zh{这影响了} (cela a influencé), \zh{总之/如果/就必须} (en somme/si/alors il faut).
\end{itemize}
\end{exoresolu}

\begin{methode}[Réviser et corriger ses erreurs types (Métacognition grammaticale)]
\textbf{Objectif :} Développer l'auto-correction systématique avant de rendre une copie ou de parler.
\begin{enumerate}
    \item \textbf{Check-list \zh{越来越} (cochez mentalement) :}
    \begin{itemize}
        \item[$\square$] Prédicat = Adjectif ? $\rightarrow$ \zh{Sujet + 越来越 + Adj.} (Ex. \zh{越来越好}). OK.
        \item[$\square$] Prédicat = Verbe d'action ? $\rightarrow$ \zh{Sujet + V + 得 + 越来越 + Compl.} (Ex. \zh{写得越来越好}). \textit{Pas \zh{*越来越写得好}.}
        \item[$\square$] Prédicat = Verbe statif (\zh{喜欢、爱、怕}) ? $\rightarrow$ \zh{Sujet + 越来越 + V.} (Ex. \zh{越来越喜欢}). Pas de \zh{得}.
        \item[$\square$] \zh{了} en fin de phrase pour le changement d'état ?
        \item[$\square$] Caractère \zh{越} (pied \zh{走}) et non \zh{月} (lune) ?
    \end{itemize}
    \item \textbf{Check-list \zh{一方面……另一方面……} :}
    \begin{itemize}
        \item[$\square$] Même thème général pour les deux clauses ?
        \item[$\square$] Position initiale de clause ? (Pas \zh{*我一方面觉得...} $\rightarrow$ \zh{一方面，我觉得...}).
        \item[$\square$] Ponctuation : \zh{；} ou \zh{。} entre clauses complètes (pas virgule seule).
        \item[$\square$] Distinction : \zh{一方面……另一方面……} (deux facettes d'UN sujet) vs \zh{一个方面……另一个方面……} (énumération d'éléments distincts).
    \end{itemize}
    \item \textbf{Relire à l'envers} (dernière phrase $\rightarrow$ première) pour couper l'anticipation du sens et voir la syntaxe pure.
\end{enumerate}
\end{methode}

\begin{exoresolu}[Atelier correction d'erreurs fréquentes (Corpus élève)]
\textbf{Énoncé :} Corrigez les 5 phrases suivantes (une erreur majeure par phrase). Réécrivez la phrase correcte (caractères + pinyin) et expliquez l'erreur en une phrase.

\begin{enumerate}
    \item \zh{我越来越喜欢吃辣了。} (Contexte : « J'aime de plus en plus manger épicé »).
    \item \zh{他跑越来越快得。} (Ordre des mots / Particule).
    \item \zh{一方面我喜欢中国菜，但是另一方面我不喜欢辣。} (Connecteur / Structure).
    \item \zh{学习中文一个方面难，另一个方面有用。} (Vocabulaire figé / Classificateur).
    \item \zh{雅温得的天气越来越热天。} (Redondance / Structure nominale).
\end{enumerate}

\tcblower
\textbf{Solution détaillée :}

\begin{enumerate}
    \item \textbf{Correction :} \zh{我越来越喜欢吃辣了。} (Wǒ yuè lái yuè xǐhuan chī là le.) \textbf{Phrase correcte !} 
    \textit{Explication :} \zh{喜欢} est un verbe d'état psychologique (statif), pas un verbe d'action dynamique ; il ne prend pas de complément de degré avec \zh{得}. \zh{越来越喜欢} est la structure correcte.
    
    \item \textbf{Correction :} \zh{他跑得越来越快了。} (Tā pǎo de yuè lái yuè kuài le.)
    \textit{Explication :} Structure V + \zh{得} + Complément de degré. \zh{越来越} se place \textbf{après} \zh{得}, devant l'adjectif \zh{快}. La particule \zh{得} ne va jamais en fin de phrase.
    
    \item \textbf{Correction :} \zh{一方面我喜欢中国菜，另一方面我不喜欢辣。} (Yī fāng miàn wǒ xǐhuan Zhōngguó cài, lìng wài yī fāng miàn wǒ bù xǐhuan là.)
    \textit{Explication :} \zh{一方面……另一方面……} structure déjà l'opposition (concession). Le connecteur \zh{但是} (dànshì) est redondant et alourdit la phrase ; on utilise soit la structure binaire, soit \zh{但是}, pas les deux ensemble.
    
    \item \textbf{Correction :} \zh{学习中文一方面难，另一方面有用。} (Xuéxí Zhōngwén yī fāng miàn nán, lìng wài yī fāng miàn yǒuyòng.)
    \textit{Explication :} L'expression figée pour contraster deux propriétés d'un même sujet est \zh{一方面……另一方面……}. \zh{一个方面 / 另一个方面} s'utilise pour énumérer des items distincts (ex: \zh{一个方面是语法，另一个方面是词汇}), pas pour opposer deux qualités du même objet.
    
    \item \textbf{Correction :} \zh{雅温得的天气越来越热了。} (Yǎwēndé de tiānqì yuè lái yuè rè le.)
    \textit{Explication :} \zh{热} (rè) est l'adjectif « chaud ». \zh{天气热} = « le temps est chaud ». Ajouter \zh{天} (tiān, ciel/jour) après \zh{热} crée une redondance syntaxique incorrecte (\zh{*chaud jour}). On dit \zh{天气越来越热} ou \zh{天越来越热} (sujet \zh{天}).
\end{enumerate}
\end{exoresolu}

\begin{attention}[Les 5 erreurs qui coûtent des points au Bac]
\begin{enumerate}
    \item \textbf{Ordre des mots avec \zh{越来越} + Verbe d'action :} Écrire \zh{*他越来越跑得快} ou \zh{*他跑越来越快得}. 
    \cle{Règle} : \zh{Sujet + Verbe + 得 + 越来越 + Adj/Adv}. (Correct : \zh{他跑得越来越快}).
    \item \textbf{Confusion \zh{的 / 得 / 地} après \zh{越来越} :} Écrire \zh{*越来越的快} ou \zh{*越来越地热}. 
    \cle{Règle} : \zh{越来越} + Adjectif direct (pas de particules). \zh{越来越} + Verbe + \zh{得} + Complément. Jamais \zh{的/地} immédiatement après \zh{越来越}.
    \item \textbf{Placement de \zh{一方面} / \zh{另一方面} :} Les mettre après le sujet (\zh{*我一方面觉得...}). 
    \cle{Règle} : Ces locutions sont des \textbf{adverbiaux de thème}, elles ouvrent la clause : \zh{一方面，我觉得... 另一方面，你觉得...}.
    \item \textbf{Caractères homophones / visuels :} Confondre \zh{越} (yuè, de plus en plus / dépasser) et \zh{月} (yuè, lune / mois) ; \zh{面} (miàn, face / côté / nouilles) et \zh{们} (men, pluriel). 
    \textit{Astuce mnémotechnique :} \zh{越} a \zh{走} (zǒu, pied/marche) en bas = mouvement/progression. \zh{月} = corps/lune (image statique).
    \item \textbf{Tons sandhi sur \zh{一} (yī) dans \zh{一方面} / \zh{另一方面} :} Dire \zh{[yī] fāng miàn} (ton 1 isolé) au lieu de \zh{[yí] fāng miàn} (ton 2 devant ton 4) ou \zh{[yì] fāng miàn} (ton 4 devant ton 1/2/3). 
    \cle{Règle Bac} : \zh{一} $\rightarrow$ \textbf{yí} (devant 4e ton \zh{面}), \textbf{yì} (devant 1/2/3e ton). \zh{另} (lìng) reste 4e ton. \textit{Entraînement oral obligatoire :} \zh{yí fāng miàn, lìng wài yí fāng miàn}.
\end{enumerate}
\end{attention}

\newpage

\coursec{Exercices}

\courssub{Je vérifie mes connaissances}

\begin{exercice}[QCM : la bonne structure]
Pour chaque phrase, choisis la bonne réponse (a, b ou c).

\begin{enumerate}
\item « Il fait de plus en plus chaud. » se dit :
\begin{enumerate}
\item \zh{天气越来越热。} \emph{Tiānqì yuè lái yuè rè.}
\item \zh{天气越来越很热。} \emph{Tiānqì yuè lái yuè hěn rè.}
\item \zh{天气越越来越热。} \emph{Tiānqì yuè yuè lái yuè rè.}
\end{enumerate}
\item « D'une part j'aime le chinois, d'autre part j'aime le football. » se dit :
\begin{enumerate}
\item \zh{一方面我喜欢汉语，我喜欢足球另一方面。}
\item \zh{一方面我喜欢汉语，另一方面我喜欢足球。}
\item \zh{我喜欢汉语一方面，另一方面我喜欢足球。}
\end{enumerate}
\item Dans \zh{学中文的人越来越多} \emph{(xué Zhōngwén de rén yuè lái yuè duō)}, \zh{越来越} signifie :
\begin{enumerate}
\item « de moins en moins »
\item « de plus en plus »
\item « de temps en temps »
\end{enumerate}
\item Quelle phrase est correcte ?
\begin{enumerate}
\item \zh{我越来越喜欢说汉语。}
\item \zh{我越来越很汉语。}
\item \zh{越来越我喜欢汉语说。}
\end{enumerate}
\end{enumerate}
\end{exercice}

\begin{exercice}[Vrai ou faux ? Justifie ta réponse]
Dis si chaque phrase est correcte (V) ou incorrecte (F) et justifie en une phrase en français.

\begin{enumerate}
\item \zh{他越来越很高。} \emph{Tā yuè lái yuè hěn gāo.}
\item \zh{一方面我想去中国，另一方面我想留在喀麦隆工作。} \emph{Yì fāngmiàn wǒ xiǎng qù Zhōngguó, lìng yì fāngmiàn wǒ xiǎng liú zài Kāmàilóng gōngzuò.}
\item \zh{雅温得的汽车越来越多。} \emph{Yǎwēndé de qìchē yuè lái yuè duō.}
\item \zh{一方面天气很热。} (phrase complète et correcte)
\end{enumerate}
\end{exercice}

\begin{exercice}[Vocabulaire : complète le tableau]
Complète le tableau suivant avec le pinyin et la traduction française.

\begin{center}\begin{tabularx}{\linewidth}{|X|X|X|X|}\hline
\rowcolor{popblueL} Caractères & Pinyin & Nature & Traduction \\ \hline
越来越 & \ldots\ldots\ldots & adverbe & \ldots\ldots\ldots \\ \hline
一方面 & \ldots\ldots\ldots & locution conjonctive & \ldots\ldots\ldots \\ \hline
另一方面 & \ldots\ldots\ldots & locution conjonctive & \ldots\ldots\ldots \\ \hline
方便 & \ldots\ldots\ldots & adjectif & \ldots\ldots\ldots \\ \hline
重要 & \ldots\ldots\ldots & adjectif & \ldots\ldots\ldots \\ \hline
\end{tabularx}\end{center}
\end{exercice}

\begin{exercice}[Associe caractères et pinyin]
Associe chaque expression en caractères (1 à 5) à son pinyin (a à e).

\begin{center}\begin{tabularx}{\linewidth}{|X|X|}\hline
\rowcolor{popblueL} Caractères & Pinyin \\ \hline
1. 越来越贵 & a. lìng yì fāngmiàn \\ \hline
2. 一方面 & b. yuè lái yuè hǎo \\ \hline
3. 另一方面 & c. yuè lái yuè guì \\ \hline
4. 越来越好 & d. yì fāngmiàn \\ \hline
5. 越来越多 & e. yuè lái yuè duō \\ \hline
\end{tabularx}\end{center}
\end{exercice}

\courssub{Je m'entraîne}

\begin{exercice}[Traduis en français]
Traduis les phrases suivantes en français.

\begin{enumerate}
\item \zh{杜阿拉的人越来越多了。} \emph{Dù'ālā de rén yuè lái yuè duō le.}
\item \zh{我的汉语说得越来越好。} \emph{Wǒ de Hànyǔ shuō de yuè lái yuè hǎo.}
\item \zh{一方面学习很忙，另一方面我还要帮妈妈做家务。} \emph{Yì fāngmiàn xuéxí hěn máng, lìng yì fāngmiàn wǒ hái yào bāng māma zuò jiāwù.}
\item \zh{天气越来越热，大家都想去海边。} \emph{Tiānqì yuè lái yuè rè, dàjiā dōu xiǎng qù hǎibiān.}
\end{enumerate}
\end{exercice}

\begin{exercice}[Traduis en chinois]
Traduis les phrases suivantes en chinois (caractères et pinyin).

\begin{enumerate}
\item Le français devient de plus en plus difficile.
\item De plus en plus de jeunes Camerounais apprennent le chinois.
\item D'une part ce restaurant n'est pas cher, d'autre part les plats sont très bons.
\item Ma ville est de plus en plus belle.
\end{enumerate}
\end{exercice}

\begin{exercice}[Transformation : fusionne les phrases]
Fusionne les deux phrases en une seule phrase avec \zh{一方面……另一方面……} ou transforme-les avec \zh{越来越}, selon la consigne entre parenthèses.

\begin{enumerate}
\item \zh{这个城市很方便。这个城市太贵。} (avec \zh{一方面……另一方面……})
\item \zh{他喜欢运动。他喜欢听音乐。} (avec \zh{一方面……另一方面……})
\item \zh{中文很难。} → dis que le chinois devient « de plus en plus difficile » (avec \zh{越来越})
\item \zh{我们学校的学生很多。} → dis que les élèves sont « de plus en plus nombreux » (avec \zh{越来越})
\end{enumerate}
\end{exercice}

\begin{exercice}[Texte à trous et remise en ordre]
\textbf{A. Texte à trous.} Complète le texte avec : \zh{越来越} (2 fois), \zh{一方面}, \zh{另一方面}.

\begin{textebox}[À compléter]
现在，学中文的人\ldots\ldots\ldots 多。\emph{Xiànzài, xué Zhōngwén de rén \ldots\ldots\ldots duō.}\\
\ldots\ldots\ldots ，很多中国公司在喀麦隆工作；\ldots\ldots\ldots ，去中国的奖学金也\ldots\ldots\ldots 多。\\
\emph{\ldots\ldots\ldots , hěn duō Zhōngguó gōngsī zài Kāmàilóng gōngzuò; \ldots\ldots\ldots , qù Zhōngguó de jiǎngxuéjīn yě \ldots\ldots\ldots duō.}
\end{textebox}

\textbf{B. Remise en ordre.} Remets les mots dans le bon ordre pour former des phrases correctes.

\begin{enumerate}
\item 喜欢 / 越来越 / 我 / 喀麦隆 / 的 / 菜
\item 一方面 / 太贵 / 这个 / 房子 / ， / 另一方面 / 很远 / 也
\item 汉语 / 越来越 / 说 / 得 / 他 / 好
\end{enumerate}
\end{exercice}

\courssub{Je me prépare au Bac}

\begin{exercice}[Type Bac — Compréhension de l'écrit et questions de langue]
Lis le texte suivant, puis réponds aux questions.

\begin{textebox}[我的城市在变化]
这几年，雅温得变化很大。楼房越来越高，马路上的汽车也越来越多。一方面，城市越来越方便：商店多了，学校多了，医院也多了；另一方面，问题也越来越多：房子越来越贵，交通越来越堵。我觉得，一方面我们应该欢迎这些变化，另一方面也应该想办法解决新的问题。\\
\emph{Zhè jǐ nián, Yǎwēndé biànhuà hěn dà. Lóufáng yuè lái yuè gāo, mǎlù shang de qìchē yě yuè lái yuè duō. Yì fāngmiàn, chéngshì yuè lái yuè fāngbiàn: shāngdiàn duō le, xuéxiào duō le, yīyuàn yě duō le; lìng yì fāngmiàn, wèntí yě yuè lái yuè duō: fángzi yuè lái yuè guì, jiāotōng yuè lái yuè dǔ. Wǒ juéde, yì fāngmiàn wǒmen yīnggāi huānyíng zhèxiē biànhuà, lìng yì fāngmiàn yě yīnggāi xiǎng bànfǎ jiějué xīn de wèntí.}
\end{textebox}

\textbf{Questions de compréhension} (réponds en français) :
\begin{enumerate}
\item Quels changements positifs l'auteur mentionne-t-il dans sa ville ?
\item Quels problèmes apparaissent selon l'auteur ?
\item Quelle est l'attitude de l'auteur face à ces changements ?
\end{enumerate}

\textbf{Questions de langue} :
\begin{enumerate}
\item Releve dans le texte trois occurrences de la structure \zh{越来越} et traduis-les en français.
\item Releve une occurrence de la structure \zh{一方面……另一方面……} et explique en français son emploi dans cette phrase.
\item Construis une phrase personnelle avec \zh{越来越} sur le thème de ta ville ou de ton école.
\end{enumerate}
\end{exercice}

\begin{exercice}[Type Bac — Thème et version]
\textbf{A. Thème.} Traduis en chinois (caractères et pinyin) :

« Le chinois devient de plus en plus important : d'une part beaucoup d'entreprises chinoises travaillent au Cameroun, d'autre part les bourses pour la Chine sont de plus en plus nombreuses. »

\textbf{B. Version.} Traduis en français le texte suivant :

\begin{textebox}[Version]
现在，学汉语的喀麦隆学生越来越多。一方面，汉语课越来越有意思；另一方面，会汉语的人找工作也越来越容易。所以，我觉得汉语越来越重要。\\
\emph{Xiànzài, xué Hànyǔ de Kāmàilóng xuésheng yuè lái yuè duō. Yì fāngmiàn, Hànyǔ kè yuè lái yuè yǒu yìsi; lìng yì fāngmiàn, huì Hànyǔ de rén zhǎo gōngzuò yě yuè lái yuè róngyì. Suǒyǐ, wǒ juéde Hànyǔ yuè lái yuè zhòngyào.}
\end{textebox}
\end{exercice}

\begin{exercice}[Type Bac — Correction d'erreurs et expression écrite]
\textbf{A. Correction d'erreurs.} La phrase suivante contient \textbf{trois fautes}. Identifie-les, corrige-les et explique chaque correction en français :

\begin{center}
\zh{*他越来越很喜欢中文，一方面他想中国，另一方面。}\\
\emph{Tā yuè lái yuè hěn xǐhuan Zhōngwén, yì fāngmiàn tā xiǎng Zhōngguó, lìng yì fāngmiàn.}
\end{center}

\textbf{B. Expression écrite.} Rédige en chinois un court texte intitulé \zh{我的城市在变化} \emph{(Wǒ de chéngshì zài biànhuà)} — « Ma ville change ».

Consignes :
\begin{itemize}
\item \textbf{5 phrases minimum} (environ 60 à 80 caractères) ;
\item emploie \textbf{au moins deux fois} la structure \zh{越来越} ;
\item emploie \textbf{au moins une fois} la structure \zh{一方面……另一方面……} ;
\item termine par une phrase donnant ton avis sur ces changements.
\end{itemize}
\end{exercice}

\newpage

\coursec{Corrigés des exercices}

\begin{corrige}[Exercice 1 — QCM : la bonne structure]
\begin{enumerate}
\item \textbf{Réponse a.} \zh{天气越来越热。} \emph{Tiānqì yuèláiyuè rè.} « Il fait de plus en plus chaud. »\\
Justification : après \zh{越来越}, l'adjectif seul suffit ; on ne met jamais \zh{很} (réponse b fautive) et l'ordre fixe est \zh{越来越} (réponse c fautive : \zh{越越来越} n'existe pas).
\item \textbf{Réponse b.} \zh{一方面我喜欢汉语，另一方面我喜欢足球。} \emph{Yìfāngmiàn wǒ xǐhuan Hànyǔ, lìngyífāngmiàn wǒ xǐhuan zúqiú.}\\
Justification : \zh{一方面} et \zh{另一方面} se placent en tête de chaque proposition (ou juste après le sujet), jamais en fin de phrase (a) ni après le verbe (c).
\item \textbf{Réponse b.} \zh{越来越} signifie « de plus en plus ». \zh{学中文的人越来越多} \emph{Xué Zhōngwén de rén yuèláiyuè duō} = « Les gens qui apprennent le chinois sont de plus en plus nombreux. »
\item \textbf{Réponse a.} \zh{我越来越喜欢说汉语。} \emph{Wǒ yuèláiyuè xǐhuan shuō Hànyǔ.} « J'aime de plus en plus parler chinois. »\\
Justification : b est fautive (\zh{很} interdit après \zh{越来越} et \zh{汉语} n'est pas un adjectif) ; c est fautive (ordre des mots bouleversé : le sujet \zh{我} doit précéder \zh{越来越}).
\end{enumerate}
\end{corrige}

\begin{corrige}[Exercice 2 — Vrai ou faux ? Justifie ta réponse]
\begin{enumerate}
\item \textbf{F.} \zh{*他越来越很高。} est incorrecte : après \zh{越来越}, l'adjectif ne peut pas être précédé de \zh{很} (ni de \zh{非常}, \zh{太}…). Correction : \zh{他越来越高。} \emph{Tā yuèláiyuè gāo.} « Il est de plus en plus grand. »
\item \textbf{V.} \zh{一方面我想去中国，另一方面我想留在喀麦隆工作。} \emph{Yìfāngmiàn wǒ xiǎng qù Zhōngguó, lìngyífāngmiàn wǒ xiǎng liú zài Kāmàilóng gōngzuò.} est correcte : les deux locutions encadrent bien deux propositions complètes et parallèles. « D'une part je veux aller en Chine, d'autre part je veux rester travailler au Cameroun. »
\item \textbf{V.} \zh{雅温得的汽车越来越多。} \emph{Yǎwēndé de qìchē yuèláiyuè duō.} « Il y a de plus en plus de voitures à Yaoundé. » Structure correcte : Sujet + \zh{越来越} + adjectif.
\item \textbf{F.} \zh{*一方面天气很热。} est incomplète : \zh{一方面} appelle obligatoirement son pendant \zh{另一方面} ; la structure fonctionne par paire. Correction possible : \zh{一方面天气很热，另一方面交通也很堵。} \emph{Yìfāngmiàn tiānqì hěn rè, lìngyífāngmiàn jiāotōng yě hěn dǔ.} « D'une part il fait très chaud, d'autre part la circulation est très bloquée. »
\end{enumerate}
\end{corrige}

\begin{corrige}[Exercice 3 — Vocabulaire : complète le tableau]
\begin{center}\begin{tabularx}{\linewidth}{|X|X|X|X|}\hline
\rowcolor{popblueL} Caractères & Pinyin & Nature & Traduction \\ \hline
越来越 & yuèláiyuè & adverbe & de plus en plus \\ \hline
一方面 & yìfāngmiàn & locution conjonctive & d'une part \\ \hline
另一方面 & lìngyífāngmiàn & locution conjonctive & d'autre part \\ \hline
方便 & fāngbiàn & adjectif & pratique, commode \\ \hline
重要 & zhòngyào & adjectif & important \\ \hline
\end{tabularx}\end{center}
Points de vigilance : \zh{越} porte le 4e ton (\emph{yuè}) ; dans \zh{一方面}, \zh{一} se prononce \emph{yì} (2e ton) par sandhi devant le 1er ton de \zh{方} ; \zh{另} se prononce \emph{lìng} (4e ton), à ne pas confondre avec \zh{另外} \emph{lìngwài} « en plus, d'ailleurs ».
\end{corrige}

\begin{corrige}[Exercice 4 — Associe caractères et pinyin]
\begin{center}\begin{tabularx}{\linewidth}{|X|X|X|}\hline
\rowcolor{popblueL} Caractères & Pinyin & Traduction \\ \hline
1. 越来越贵 & c. yuèláiyuè guì & de plus en plus cher \\ \hline
2. 一方面 & d. yìfāngmiàn & d'une part \\ \hline
3. 另一方面 & a. lìngyífāngmiàn & d'autre part \\ \hline
4. 越来越好 & b. yuèláiyuè hǎo & de mieux en mieux \\ \hline
5. 越来越多 & e. yuèláiyuè duō & de plus en plus nombreux \\ \hline
\end{tabularx}\end{center}
Réponses : 1-c ; 2-d ; 3-a ; 4-b ; 5-e.
\end{corrige}

\begin{corrige}[Exercice 5 — Traduis en français]
\begin{enumerate}
\item \zh{杜阿拉的人越来越多了。} \emph{Dūālā de rén yuèláiyuè duō le.} « Il y a de plus en plus de monde à Douala. » (Le \zh{了} final marque le changement de situation déjà constaté.)
\item \zh{我的汉语说得越来越好。} \emph{Wǒ de Hànyǔ shuō de yuèláiyuè hǎo.} « Je parle chinois de mieux en mieux. » (Structure : verbe + \zh{得} + \zh{越来越} + adjectif, pour évaluer la manière.)
\item \zh{一方面学习很忙，另一方面我还要帮妈妈做家务。} \emph{Yìfāngmiàn xuéxí hěn máng, lìngyífāngmiàn wǒ hái yào bāng māma zuò jiāwù.} « D'une part les études me prennent beaucoup de temps, d'autre part je dois encore aider ma mère à faire le ménage. »
\item \zh{天气越来越热，大家都想去海边。} \emph{Tiānqì yuèláiyuè rè, dàjiā dōu xiǎng qù hǎibiān.} « Il fait de plus en plus chaud, tout le monde veut aller au bord de la mer. »
\end{enumerate}
\end{corrige}

\begin{corrige}[Exercice 6 — Traduis en chinois]
\begin{enumerate}
\item \zh{法语越来越难。} \emph{Fǎyǔ yuèláiyuè nán.} (Pas de \zh{很} devant \zh{难}.)
\item \zh{学汉语的喀麦隆年轻人越来越多。} \emph{Xué Hànyǔ de Kāmàilóng niánqīngrén yuèláiyuè duō.} (Le groupe nominal \zh{学汉语的喀麦隆年轻人} « les jeunes Camerounais qui apprennent le chinois » précède \zh{越来越}.)
\item \zh{一方面这家饭店不贵，另一方面菜很好吃。} \emph{Yìfāngmiàn zhè jiā fàndiàn bù guì, lìngyífāngmiàn cài hěn hǎochī.} (\zh{贵} « cher » se nie par \zh{不} ; dans la seconde proposition, l'adjectif peut normalement être précédé de \zh{很}.)
\item \zh{我的城市越来越漂亮。} \emph{Wǒ de chéngshì yuèláiyuè piàoliang.}
\end{enumerate}
Points de vigilance : ne jamais écrire \zh{*越来越很…} ; ne pas traduire « d'une part… d'autre part… » par deux phrases séparées, la paire \zh{一方面……另一方面……} doit relier les deux propositions dans une même phrase.
\end{corrige}

\begin{corrige}[Exercice 7 — Transformation : fusionne les phrases]
\begin{enumerate}
\item \zh{一方面这个城市很方便，另一方面这个城市太贵。} \emph{Yìfāngmiàn zhège chéngshì hěn fāngbiàn, lìngyífāngmiàn zhège chéngshì tài guì.} « D'une part cette ville est très pratique, d'autre part elle est trop chère. » (On peut ne pas répéter le sujet dans la seconde proposition : \zh{……另一方面太贵。})
\item \zh{他一方面喜欢运动，另一方面喜欢听音乐。} \emph{Tā yìfāngmiàn xǐhuan yùndòng, lìngyífāngmiàn xǐhuan tīng yīnyuè.} « D'une part il aime le sport, d'autre part il aime écouter de la musique. » (Quand le sujet est identique, on peut placer \zh{一方面} après le sujet.)
\item \zh{中文越来越难。} \emph{Zhōngwén yuèláiyuè nán.} « Le chinois devient de plus en plus difficile. »
\item \zh{我们学校的学生越来越多。} \emph{Wǒmen xuéxiào de xuésheng yuèláiyuè duō.} « Les élèves de notre école sont de plus en plus nombreux. »
\end{enumerate}
\end{corrige}

\begin{corrige}[Exercice 8 — Texte à trous et remise en ordre]
\textbf{A. Texte à trous.} Texte complété :

现在，学中文的人\cle{越来越}多。\emph{Xiànzài, xué Zhōngwén de rén yuèláiyuè duō.}\\
\cle{一方面}，很多中国公司在喀麦隆工作；\cle{另一方面}，去中国的奖学金也\cle{越来越}多。\\
\emph{Yìfāngmiàn, hěn duō Zhōngguó gōngsī zài Kāmàilóng gōngzuò; lìngyífāngmiàn, qù Zhōngguó de jiǎngxuéjīn yě yuèláiyuè duō.}

Traduction : « Aujourd'hui, les gens qui apprennent le chinois sont de plus en plus nombreux. D'une part, beaucoup d'entreprises chinoises travaillent au Cameroun ; d'autre part, les bourses pour aller en Chine sont aussi de plus en plus nombreuses. »

\textbf{B. Remise en ordre.}
\begin{enumerate}
\item \zh{我越来越喜欢喀麦隆的菜。} \emph{Wǒ yuèláiyuè xǐhuan Kāmàilóng de cài.} « J'aime de plus en plus la cuisine camerounaise. » (Sujet + \zh{越来越} + verbe + complément.)
\item \zh{一方面这个房子太贵，另一方面也很远。} \emph{Yìfāngmiàn zhège fángzi tài guì, lìngyífāngmiàn yě hěn yuǎn.} « D'une part cette maison est trop chère, d'autre part elle est aussi très loin. »
\item \zh{他汉语说得越来越好。} \emph{Tā Hànyǔ shuō de yuèláiyuè hǎo.} « Il parle chinois de mieux en mieux. » (Thème \zh{汉语} en tête, puis verbe + \zh{得} + \zh{越来越好}.)
\end{enumerate}
\end{corrige}

\begin{corrige}[Exercice 9 — Type Bac : Compréhension de l'écrit et questions de langue]
\textbf{Barème indicatif (sur 10 points)} : compréhension 4 pts (1 + 1,5 + 1,5) ; questions de langue 6 pts (2 + 2 + 2).

\textbf{Questions de compréhension}
\begin{enumerate}
\item Changements positifs : la ville devient de plus en plus pratique — il y a davantage de magasins, d'écoles et d'hôpitaux ; les immeubles sont de plus en plus hauts.
\item Problèmes : les logements deviennent de plus en plus chers et la circulation est de plus en plus embouteillée ; les problèmes en général sont de plus en plus nombreux.
\item Attitude de l'auteur : elle est équilibrée et constructive — il faut d'une part accueillir favorablement ces changements, et d'autre part chercher des solutions aux nouveaux problèmes.
\end{enumerate}

\textbf{Questions de langue}
\begin{enumerate}
\item Trois occurrences possibles parmi : \zh{楼房越来越高} \emph{Lóufáng yuèláiyuè gāo} « les immeubles sont de plus en plus hauts » ; \zh{汽车也越来越多} \emph{Qìchē yě yuèláiyuè duō} « les voitures sont de plus en plus nombreuses » ; \zh{城市越来越方便} \emph{Chéngshì yuèláiyuè fāngbiàn} « la ville est de plus en plus pratique » ; \zh{房子越来越贵} \emph{Fángzi yuèláiyuè guì} « les logements sont de plus en plus chers » ; \zh{交通越来越堵} \emph{Jiāotōng yuèláiyuè dǔ} « la circulation est de plus en plus bloquée ».
\item Exemple : \zh{一方面，城市越来越方便……；另一方面，问题也越来越多……} \emph{Yìfāngmiàn, chéngshì yuèláiyuè fāngbiàn……; lìngyífāngmiàn, wèntí yě yuèláiyuè duō……}. Cette structure présente deux aspects complémentaires ou contrastés d'une même situation : ici, l'auteur oppose les avantages et les inconvénients du développement de la ville.
\item Exemple de production : \zh{我们学校的图书馆越来越好。} \emph{Wǒmen xuéxiào de túshūguǎn yuèláiyuè hǎo.} « La bibliothèque de notre école est de mieux en mieux. »
\end{enumerate}

\textbf{Points de vigilance} : en version, bien rendre \zh{越来越} par « de plus en plus / de mieux en mieux » et non par « très » ; relever la structure signifie recopier la proposition entière, pas seulement les deux caractères.
\end{corrige}

\begin{corrige}[Exercice 10 — Type Bac : Thème et version]
\textbf{Barème indicatif (sur 10 points)} : thème 5 pts ; version 5 pts.

\textbf{A. Thème — proposition de traduction :}

\zh{汉语越来越重要：一方面，很多中国公司在喀麦隆工作；另一方面，去中国的奖学金也越来越多。}\\
\emph{Hànyǔ yuèláiyuè zhòngyào: yìfāngmiàn, hěn duō Zhōngguó gōngsī zài Kāmàilóng gōngzuò; lìngyífāngmiàn, qù Zhōngguó de jiǎngxuéjīn yě yuèláiyuè duō.}

Points de vigilance : \zh{重要} seul après \zh{越来越} (pas de \zh{很}) ; \zh{公司} « entreprise » avec le classificateur optionnel \zh{家} si l'on compte ; « de plus en plus nombreuses » se rend par \zh{越来越多}, avec \zh{也} dans la seconde proposition pour marquer l'addition.

\textbf{B. Version — proposition de traduction :}

« Aujourd'hui, les élèves camerounais qui apprennent le chinois sont de plus en plus nombreux. D'une part, les cours de chinois sont de plus en plus intéressants ; d'autre part, les personnes qui savent le chinois trouvent du travail de plus en plus facilement. C'est pourquoi je pense que le chinois devient de plus en plus important. »

\zh{现在，学中文的喀麦隆学生越来越多。一方面，中文课越来越有意思；另一方面，会汉语的人找工作越来越容易。所以我觉得中文越来越重要。}\\
\emph{Xiànzài, xué Zhōngwén de Kāmàilóng xuésheng yuèláiyuè duō. Yìfāngmiàn, Zhōngwén kè yuèláiyuè yǒuyìsi; lìngyífāngmiàn, huì Hànyǔ de rén zhǎo gōngzuò yuèláiyuè róngyì. Suǒyǐ wǒ juéde Zhōngwén yuèláiyuè zhòngyào.}

Points de vigilance : \zh{有意思} = « intéressant » ; \zh{会汉语的人} = « les personnes qui savent (parler) chinois » ; \zh{容易} = « facile » ; \zh{所以} = « c'est pourquoi, donc ».
\end{corrige}

\begin{corrige}[Exercice 11 — Type Bac : Correction d'erreurs et expression écrite]
\textbf{Barème indicatif (sur 10 points)} : correction d'erreurs 4 pts (1 pt par faute identifiée et corrigée + 1 pt pour les explications) ; expression écrite 6 pts (respect des consignes 2 pts, exactitude grammaticale 2 pts, vocabulaire et caractères 2 pts).

\textbf{A. Correction d'erreurs.} Phrase fautive : \zh{*他越来越很喜欢中文，一方面他想中国，另一方面。}

\begin{enumerate}
\item \textbf{Faute 1 :} \zh{越来越很喜欢} → \zh{越来越喜欢}. Après \zh{越来越}, on ne peut pas ajouter \zh{很} : l'intensité croissante est déjà exprimée par la structure elle-même.
\item \textbf{Faute 2 :} \zh{他想中国} → \zh{他想去中国}. \zh{想} + nom signifie « manquer, penser à » ; pour dire « vouloir aller en Chine », il faut \zh{想} + verbe : \zh{想去中国}.
\item \textbf{Faute 3 :} \zh{另一方面} seul, sans proposition → il faut compléter, par exemple \zh{另一方面他也想留在喀麦隆}. La structure \zh{一方面……另一方面……} fonctionne obligatoirement par paire de propositions complètes.
\end{enumerate}

Phrase corrigée : \zh{他越来越喜欢中文，一方面他想去中国，另一方面他也想留在喀麦隆。}\\
\emph{Tā yuèláiyuè xǐhuan Zhōngwén, yìfāngmiàn tā xiǎng qù Zhōngguó, lìngyífāngmiàn tā yě xiǎng liú zài Kāmàilóng.}\\
« Il aime de plus en plus le chinois : d'une part il voudrait aller en Chine, d'autre part il voudrait aussi rester au Cameroun. »

\textbf{B. Expression écrite — proposition de rédaction modèle :}

\textbf{Modèle : 我的城市在变化}

\zh{这几年，我的城市变化很大。楼房越来越高，马路越来越宽。}\\
\emph{Zhè jǐ nián, wǒ de chéngshì biànhuà hěn dà. Lóufáng yuèláiyuè gāo, mǎlù yuèláiyuè kuān.}

\zh{一方面，生活越来越方便：商店多了，学校也多了；}\\
\emph{Yìfāngmiàn, shēnghuó yuèláiyuè fāngbiàn: shāngdiàn duō le, xuéxiào yě duō le;}

\zh{另一方面，汽车越来越多，交通越来越堵。}\\
\emph{Lìngyífāngmiàn, qìchē yuèláiyuè duō, jiāotōng yuèláiyuè dǔ.}

\zh{我觉得这些变化很好，但是我们也要解决新的问题。}\\
\emph{Wǒ juéde zhèxiē biànhuà hěn hǎo, dànshì wǒmen yě yào jiějué xīn de wèntí.}

« Ces dernières années, ma ville a beaucoup changé. Les immeubles sont de plus en plus hauts, les routes de plus en plus larges. D'une part, la vie devient de plus en plus pratique : il y a plus de magasins et plus d'écoles ; d'autre part, les voitures sont de plus en plus nombreuses et la circulation de plus en plus bloquée. Je trouve que ces changements sont une bonne chose, mais nous devons aussi résoudre les nouveaux problèmes. »

\textbf{Points de vigilance pour la rédaction} : vérifier que \zh{越来越} apparaît au moins deux fois sans \zh{很} ; vérifier que \zh{一方面……另一方面……} encadre deux propositions complètes ; la dernière phrase doit exprimer un avis (\zh{我觉得……}, \zh{我认为……}) ; soigner les caractères : ne pas confondre \zh{堵} \emph{dǔ} « bloqué » et \zh{都} \emph{dōu}, ni \zh{便} et \zh{更}.
\end{corrige}

\newpage

\coursec{Fiche bilan}

\begin{popfigure}
\begin{center}
\begin{tikzpicture}[mindmap, grow cyclic, every node/.style={concept, text=white, font=\small},
  root concept/.append style={concept color=popblue, font=\bfseries},
  level 1/.append style={level distance=3.2cm, sibling angle=60},
  level 1 concept/.append style={font=\small}]
\node[root concept, align=center]{Leçon 24\\ 越来越 / 一方面\ldots 另一方面\ldots}
  child[concept color=poporange]{ node[align=center]{越来越 + Adj.\\ de plus en plus} }
  child[concept color=popgreen]{ node[align=center]{越来越 + verbe\\ psychologique} }
  child[concept color=poppurple]{ node[align=center]{Pas de 很\\ après 越来越} }
  child[concept color=poppink]{ node[align=center]{一方面\ldots\\ 另一方面\ldots} }
  child[concept color=popgold]{ node[align=center]{Emploi : avantages\\ et inconvénients} }
  child[concept color=popblue!70]{ node[align=center]{Pièges\\ des francophones} };
\end{tikzpicture}
\end{center}
\legende{Carte mentale de la leçon 24 : les deux structures à maîtriser.}
\end{popfigure}

\begin{bilan}
\textbf{Lexique clé}
\begin{center}
\begin{tabularx}{\linewidth}{|X|X|X|X|}
\hline
\rowcolor{popblueL} Caractères & Pinyin & Nature & Traduction \\
\hline
越来越 & yuè lái yuè & adverbe & de plus en plus \\
\hline
方面 & fāngmiàn & nom & aspect, côté \\
\hline
一方面\ldots 另一方面\ldots & yì fāngmiàn\ldots lìng yì fāngmiàn\ldots & corrélation & d'une part\ldots d'autre part\ldots \\
\hline
进步 & jìnbù & verbe / nom & progresser ; progrès \\
\hline
发展 & fāzhǎn & verbe / nom & se développer ; développement \\
\hline
环境 & huánjìng & nom & environnement \\
\hline
污染 & wūrǎn & nom / verbe & pollution ; polluer \\
\hline
\end{tabularx}
\end{center}

\textbf{Règles de grammaire à connaître par cœur}
\begin{center}
\begin{tabularx}{\linewidth}{|X|X|X|}
\hline
\rowcolor{popblueL} Structure & Exemple (caractères + pinyin) & Traduction \\
\hline
越来越 + adjectif & 天气越来越热。Tiānqì yuè lái yuè rè. & Il fait de plus en plus chaud. \\
\hline
越来越 + verbe psychologique (喜欢, 想, 怕\ldots) & 我越来越喜欢汉语。Wǒ yuè lái yuè xǐhuan Hànyǔ. & J'aime de plus en plus le chinois. \\
\hline
Sujet + 一方面 + phrase 1, 另一方面 + phrase 2 & 一方面我想学汉语，另一方面我想了解中国文化。Yì fāngmiàn wǒ xiǎng xué Hànyǔ, lìng yì fāngmiàn wǒ xiǎng liǎojiě Zhōngguó wénhuà. & D'une part je veux apprendre le chinois, d'autre part je veux connaître la culture chinoise. \\
\hline
\end{tabularx}
\end{center}

\textbf{Méthodes express}
\begin{itemize}
\item Pour exprimer une \cle{évolution progressive} : repérer l'adjectif ou le verbe psychologique, placer \zh{越来越} juste devant, \textbf{sans} \zh{很}.
\item Pour présenter \cle{deux aspects parallèles} d'un sujet (avantages/inconvénients, pour/contre) : ouvrir avec \zh{一方面}, fermer avec \zh{另一方面}, séparer par une virgule chinoise ，.
\item Vérifier à voix haute : la phrase contient-elle un seul sujet ? Les deux « aspects » sont-ils bien de même nature (deux qualités, deux problèmes) ?
\end{itemize}
\end{bilan}

\begin{aretenir}[Checklist avant le Bac]
\begin{itemize}
\item[$\square$] Je sais que \zh{越来越} signifie « de plus en plus » et marque un changement progressif.
\item[$\square$] Je place \zh{越来越} immédiatement devant l'adjectif : \zh{越来越高} (de plus en plus grand/haut).
\item[$\square$] Je n'écris jamais \zh{越来越很热} : \zh{很} est interdit après \zh{越来越}.
\item[$\square$] Je sais que \zh{越来越} peut précéder un verbe psychologique (\zh{喜欢}, \zh{想}, \zh{害怕}) mais pas un verbe d'action simple.
\item[$\square$] Je n'ajoute pas \zh{了} à la fin d'une phrase en \zh{越来越} pour marquer le changement : la structure suffit.
\item[$\square$] Je construis \zh{一方面\ldots 另一方面\ldots} avec deux propositions complètes et équilibrées.
\item[$\square$] Je sais que \zh{另一方面} s'écrit avec \zh{另} (autre) et non avec un autre caractère homophone.
\item[$\square$] Je traduis correctement : \zh{一方面} = « d'une part », \zh{另一方面} = « d'autre part ».
\item[$\square$] Je sais utiliser cette corrélation pour structurer un paragraphe argumentatif (ex. : les avantages et les inconvénients d'Internet).
\item[$\square$] Je sépare les deux propositions par la virgule chinoise ， et je termine par 。
\item[$\square$] Je peux produire au moins deux phrases originales avec chaque structure, sur le Cameroun ou la Chine.
\item[$\square$] Je relis ma production en vérifiant l'ordre des mots : sujet + \zh{越来越} + adjectif.
\end{itemize}
\end{aretenir}

\findecours

\end{document}
